\input{preamble}

% OK, start here.
%
\begin{document}

\title{Courbes algébriques}


\maketitle

\phantomsection
\label{section-phantom}

\tableofcontents

\section{Introduction}
\label{section-introduction}

\noindent
Dans ce chapitre, nous développons quelques éléments de la théorie des courbes algébriques. Pour les courbes algébriques sur le corps des nombres complexes, on pourra consulter l'ouvrage \cite{ACGH}.

\medskip\noindent
Rappel des résultats déjà établis. Outre les résultats généraux de géométrie algébrique, nous avons déjà démontré plusieurs résultats propres aux courbes algébriques. En voici la liste. \begin{enumerate} \item Nous avons étudié les ouverts affines et les faisceaux inversibles amples sur les schémas noethériens de dimension $1$ dans Variétés, section \ref{varieties-section-dimension-one}. \item Nous avons vu qu'une courbe est soit affine, soit projective dans Variétés, section \ref{varieties-section-curves}. \item Nous avons étudié les degrés des modules localement libres sur les courbes propres dans Variétés, section \ref{varieties-section-divisors-curves}. \item Nous avons étudié le schéma de Picard d'une courbe projective non singulière sur un corps algébriquement clos dans Schémas de Picard des courbes, section \ref{pic-section-introduction}.
\end{enumerate}





\section{Courbes et corps de fonctions}
\label{section-curves-function-fields}

\noindent
Dans cette section, nous précisons les résultats de Variétés, section \ref{varieties-section-varieties-rational-maps} dans le cas des courbes.

\begin{lemma}
\label{lemma-extend-over-dvr}
Soit $k$ un corps. Soient $X$ une courbe et $Y$ une variété propre. Soient $U \subset X$ un ouvert non vide et $f : U \to Y$ un morphisme. Si $x \in X$ est un point fermé tel que $\mathcal{O}_{X, x}$ soit un anneau de valuation discrète, il existe un ouvert $U \subset U' \subset X$ contenant $x$ et un morphisme de variétés $f' : U' \to Y$ prolongeant $f$.
\end{lemma}

\begin{proof}
C'est un cas particulier de Morphismes, lemme \ref{morphisms-lemma-extend-across}.
\end{proof}

\begin{lemma}
\label{lemma-extend-over-normal-curve}
Soit $k$ un corps. Soient $X$ une courbe normale et $Y$ une variété propre. L'ensemble des applications rationnelles de $X$ dans $Y$ s'identifie à l'ensemble des morphismes $X \to Y$.
\end{lemma}

\begin{proof}
Toute application rationnelle de $X$ dans $Y$ se prolonge en un morphisme $X \to Y$ d'après le lemme \ref{lemma-extend-over-dvr}, puisque tout anneau local est un anneau de valuation discrète (par exemple d'après Variétés, lemme \ref{varieties-lemma-regular-point-on-curve}). Réciproquement, si deux morphismes $f,g: X \to Y$ sont équivalents en tant qu'applications rationnelles, alors $f = g$ d'après Morphismes, lemme \ref{morphisms-lemma-equality-of-morphisms}.
\end{proof}

\begin{lemma}
\label{lemma-flat}
Soit $k$ un corps. Soit $f : X \to Y$ un morphisme non constant de courbes sur $k$. Si $Y$ est normale, alors $f$ est plat.
\end{lemma}

\begin{proof}
Prenons $x \in X$ d'image $y \in Y$. Alors $\mathcal{O}_{Y, y}$ est soit un corps, soit un anneau de valuation discrète (Variétés, lemme \ref{varieties-lemma-regular-point-on-curve}). Puisque $f$ est non constant, il est dominant (car il envoie nécessairement le point générique de $X$ sur celui de $Y$). Il en résulte que $\mathcal{O}_{Y, y} \to \mathcal{O}_{X, x}$ est injectif (Morphismes, lemme \ref{morphisms-lemma-dominant-between-integral}). Ainsi $\mathcal{O}_{X, x}$ est sans torsion en tant que $\mathcal{O}_{Y, y}$-module; par conséquent, $\mathcal{O}_{X, x}$ est plat en tant que $\mathcal{O}_{Y, y}$-module d'après Compléments d'algèbre, lemme \ref{more-algebra-lemma-valuation-ring-torsion-free-flat}.
\end{proof}

\begin{lemma}
\label{lemma-finite}
Soit $k$ un corps. Soit $f : X \to Y$ un morphisme de schémas sur $k$. Supposons que \begin{enumerate} \item $Y$ soit séparé sur $k$, \item $X$ soit propre de dimension $\leq 1$ sur $k$, \item $f(Z)$ contienne au moins deux points pour toute composante irréductible $Z \subset X$ de dimension $1$. \end{enumerate} Alors $f$ est fini.
\end{lemma}

\begin{proof}
Le morphisme $f$ est propre d'après Morphismes, lemme \ref{morphisms-lemma-image-proper-scheme-closed}. Ainsi $f(X)$ est fermé et les images des points fermés sont fermées. Soit $y \in Y$ l'image d'un point fermé de $X$. Alors $f^{-1}(\{y\})$ est une partie fermée de $X$ ne contenant aucun point générique d'une composante irréductible de dimension $1$, d'après la condition (3). Il s'ensuit que $f^{-1}(\{y\})$ est fini. Par conséquent, $f$ est fini au-dessus d'un voisinage ouvert de $y$ d'après Compléments sur les morphismes, lemme \ref{more-morphisms-lemma-proper-finite-fibre-finite-in-neighbourhood} (si $Y$ est noethérien, on peut utiliser le résultat plus élémentaire Cohomologie des schémas, lemme \ref{coherent-lemma-proper-finite-fibre-finite-in-neighbourhood}). Puisque nous avons vu qu'il y a suffisamment de tels points $y$, la démonstration est achevée.
\end{proof}

\begin{lemma}
\label{lemma-extend-to-completion}
Soit $k$ un corps. Soit $X \to Y$ un morphisme de variétés, où $Y$ est propre et $X$ est une courbe. Il existe une factorisation $X \to \overline{X} \to Y$ dans laquelle $X \to \overline{X}$ est une immersion ouverte et $\overline{X}$ est une courbe projective.
\end{lemma}

\begin{proof}
Cela résulte immédiatement du lemme \ref{lemma-extend-over-dvr} et de Variétés, lemme \ref{varieties-lemma-reduced-dim-1-projective-completion}.
\end{proof}

\noindent
Voici le théorème principal de cette section. Un morphisme $f : X \to Y$ de variétés sera dit {\it constant} si son image $f(X)$ se réduit à un point $y$ de $Y$. Dans ce cas, $y$ est un point fermé de $Y$ (car l'image d'un point fermé de $X$ est un point fermé de $Y$).

\begin{theorem}
\label{theorem-curves-rational-maps}
{\emergencystretch=2em
Soit $k$ un corps. Les catégories suivantes sont canoniquement équivalentes : \begin{enumerate} \item La catégorie des extensions de corps de type fini $K/k$ de degré de transcendance $1$. \item La catégorie des courbes et des applications rationnelles dominantes. \item La catégorie des courbes projectives normales et des morphismes non constants. \item La catégorie des courbes projectives non singulières et des morphismes non constants. \item La catégorie des courbes projectives régulières et des morphismes non constants. \item La catégorie des courbes propres normales et des morphismes non constants.
\end{enumerate}
\par}
\end{theorem}

\begin{proof}
L'équivalence entre les catégories (1) et (2) est la restriction de l'équivalence de Variétés, théorème \ref{varieties-theorem-varieties-rational-maps}. En effet, une variété est une courbe si et seulement si son corps de fonctions est de degré de transcendance $1$; voir par exemple Variétés, lemme \ref{varieties-lemma-dimension-locally-algebraic}.

\medskip\noindent
Les catégories (3), (4), (5) et (6) coïncident. Tout d'abord, les termes « régulier » et « non singulier » sont synonymes; voir Propriétés, définition \ref{properties-definition-regular}. Pour les schémas noethériens de dimension $1$, normalité et régularité sont équivalentes (Propriétés, lemmes \ref{properties-lemma-regular-normal} et \ref{properties-lemma-normal-dimension-1-regular}). Voir Variétés, lemme \ref{varieties-lemma-regular-point-on-curve} pour le cas des courbes. Ainsi (3) coïncide avec (5). Enfin, (6) coïncide avec (3) d'après Variétés, lemme \ref{varieties-lemma-dim-1-proper-projective}.

\medskip\noindent
Si $f : X \to Y$ est un morphisme non constant de courbes projectives non singulières, alors $f$ envoie le point générique $\eta$ de $X$ sur le point générique $\xi$ de $Y$. On obtient donc un morphisme $k(Y) = \mathcal{O}_{Y, \xi} \to \mathcal{O}_{X, \eta} = k(X)$ dans la catégorie (1). Si deux morphismes $f,g: X \to Y$ induisent le même morphisme $k(Y) \to k(X)$, alors, par l'équivalence entre (1) et (2), $f$ et $g$ sont équivalents en tant qu'applications rationnelles; donc $f=g$ d'après le lemme \ref{lemma-extend-over-normal-curve}. Réciproquement, supposons donné un morphisme $k(Y) \to k(X)$ dans la catégorie (1). On obtient alors un morphisme $U \to Y$ pour un ouvert non vide $U \subset X$. D'après le lemme \ref{lemma-extend-over-dvr}, celui-ci se prolonge à $X$ tout entier et fournit un morphisme dans la catégorie (5). Nous obtenons ainsi un foncteur pleinement fidèle (5)$\to$(1).

\medskip\noindent
Pour achever la démonstration, il reste à montrer que tout objet $K/k$ de (1) est le corps de fonctions d'une courbe projective normale. Nous savons déjà que $K = k(X)$ pour une certaine courbe $X$. En remplaçant $X$ par sa normalisation (qui est une variété birationnelle à $X$), on peut supposer $X$ normale (Variétés, lemme \ref{varieties-lemma-normalization-locally-algebraic}). Choisissons alors $X \to \overline{X}$ avec $\overline{X} \setminus X = \{x_1, \ldots, x_n\}$ comme dans Variétés, lemme \ref{varieties-lemma-reduced-dim-1-projective-completion}. Puisque $X$ est normale et que chacun des anneaux locaux $\mathcal{O}_{\overline{X}, x_i}$ est normal, $\overline{X}$ est une courbe projective normale, comme voulu. (Remarque : on peut aussi commencer par compactifier à l'aide de Variétés, lemme \ref{varieties-lemma-dim-1-projective-completion}, puis normaliser à l'aide de Variétés, lemme \ref{varieties-lemma-normalization-locally-algebraic}. Cette méthode évite le recours au résultat quelque peu délicat Morphismes, lemme \ref{morphisms-lemma-relative-normalization-normal-codim-1}.)
\end{proof}

\begin{definition}
\label{definition-normal-projective-model}
Soit $k$ un corps. Soit $X$ une courbe. Un {\it modèle projectif non singulier de $X$} est un couple $(Y, \varphi)$, où $Y$ est une courbe projective non singulière et $\varphi : k(X) \to k(Y)$ un isomorphisme de corps de fonctions.
\end{definition}

\noindent
Un modèle projectif non singulier est déterminé à un unique isomorphisme près par le théorème \ref{theorem-curves-rational-maps}. On parle donc souvent du « modèle projectif non singulier ». On omet généralement $\varphi$ dans la notation. Attention : il se peut que $Y$ ne soit pas lisse sur $k$; le lemme \ref{lemma-nonsingular-model-smooth} montre toutefois que cela ne peut se produire qu'en caractéristique positive.

\begin{lemma}
\label{lemma-nonsingular-model-smooth}
Soit $k$ un corps. Soient $X$ une courbe et $Y$ le modèle projectif non singulier de $X$. Si $k$ est parfait, alors $Y$ est une courbe projective lisse.
\end{lemma}

\begin{proof}
Voir par exemple Variétés, lemme \ref{varieties-lemma-regular-point-on-curve}.
\end{proof}

\begin{lemma}
\label{lemma-smooth-models}
Soit $k$ un corps. Soit $X$ une courbe géométriquement irréductible sur $k$. Pour une extension de corps $K/k$, notons $Y_K$ un modèle projectif non singulier de $(X_K)_{red}$. \begin{enumerate} \item Si $X$ est propre, alors $Y_K$ est la normalisation de $X_K$. \item Il existe une extension $K/k$ finie purement inséparable telle que $Y_K$ soit lisse. \item Dès que $Y_K$ est lisse\footnote{Ou même géométriquement réduit.}, on a $H^0(Y_K, \mathcal{O}_{Y_K}) = K$. \item Étant donné un diagramme commutatif $$
\xymatrix{
\Omega & K' \ar[l] \\
K \ar[u] & k \ar[l] \ar[u]
}
$$ de corps tel que $Y_K$ et $Y_{K'}$ soient lisses, on a $Y_\Omega = (Y_K)_\Omega = (Y_{K'})_\Omega$.
\end{enumerate}
\end{lemma}

\begin{proof}
Soit $X'$ un modèle projectif non singulier de $X$. Alors $X'$ et $X$ possèdent des sous-schémas ouverts non vides isomorphes. En particulier, $X'$ est géométriquement irréductible, puisque $X$ l'est (nous omettons quelques détails). On peut donc supposer $X$ projective.

\medskip\noindent
Supposons $X$ propre. Alors $X_K$ est propre; sa normalisation $(X_K)^\nu$ est donc propre, étant finie sur un schéma propre (Variétés, lemme \ref{varieties-lemma-normalization-locally-algebraic} et Morphismes, lemmes \ref{morphisms-lemma-finite-proper} et \ref{morphisms-lemma-composition-proper}). D'autre part, $X_K$ est irréductible puisque $X$ est géométriquement irréductible. Ainsi $X_K^\nu$ est propre, normale, irréductible et birationnelle à $(X_K)_{red}$. Cela prouve (1), car toute courbe propre est projective (Variétés, lemme \ref{varieties-lemma-dim-1-proper-projective}).

\medskip\noindent
Démontrons (2). Puisque $X$ est propre et que (1) est établi, on peut appliquer Variétés, lemme \ref{varieties-lemma-finite-extension-geometrically-normal} pour trouver une extension $K/k$ finie purement inséparable telle que $Y_K$ soit géométriquement normale. Alors $Y_K$ est géométriquement régulière, puisque normalité et régularité sont équivalentes pour les courbes (Propriétés, lemme \ref{properties-lemma-normal-dimension-1-regular}). Ainsi $Y$ est une variété lisse d'après Variétés, lemme \ref{varieties-lemma-geometrically-regular-smooth}.

\medskip\noindent
Si $Y_K$ est géométriquement réduite, alors $Y_K$ est géométriquement intègre (Variétés, lemme \ref{varieties-lemma-geometrically-integral}) et l'on a $H^0(Y_K, \mathcal{O}_{Y_K}) = K$ d'après Variétés, lemme \ref{varieties-lemma-regular-functions-proper-variety}. Cela prouve (3), car toute variété lisse est géométriquement réduite (et même géométriquement régulière; voir Variétés, lemme \ref{varieties-lemma-geometrically-regular-smooth}).

\medskip\noindent
Si $Y_K$ est lisse, alors, pour toute extension $\Omega/K$, le changement de base $(Y_K)_\Omega$ est lisse sur $\Omega$ (Morphismes, lemme \ref{morphisms-lemma-base-change-smooth}). Il est donc clair que $Y_\Omega = (Y_K)_\Omega$. Cela prouve (4).
\end{proof}






\section{Séries linéaires}
\label{section-linear-series}

\noindent
Nous nous écartons de la présentation classique (voir la remarque \ref{remark-classical-linear-series}) en définissant les séries linéaires de la manière suivante.

\begin{definition}
\label{definition-linear-series}
Soit $k$ un corps. Soit $X$ un schéma propre de dimension $\leq 1$ sur $k$. Soient $d \geq 0$ et $r \geq 0$. Une {\it série linéaire de degré $d$ et de dimension $r$} est un couple $(\mathcal{L}, V)$, où $\mathcal{L}$ est un $\mathcal{O}_X$-module inversible de degré $d$ (Variétés, définition \ref{varieties-definition-degree-invertible-sheaf}) et $V \subset H^0(X, \mathcal{L})$ un sous-espace vectoriel sur $k$ de dimension $r + 1$. On dira, en abrégé, que $(\mathcal{L}, V)$ est une {\it $\mathfrak g^r_d$} sur $X$.
\end{definition}

\noindent
Nous utiliserons surtout cette notion lorsque $X$ est une courbe propre non singulière. En fait, la définition ci-dessus n'est qu'une façon de généraliser la définition classique d'une $\mathfrak g^r_d$. Par exemple, si $X$ est une courbe propre, on peut généraliser les séries linéaires en autorisant $\mathcal{L}$ à être un $\mathcal{O}_X$-module cohérent sans torsion de rang $1$. Sur une courbe non singulière, tout module cohérent sans torsion est localement libre; on retrouve donc notre notion pour les courbes propres non singulières.

\medskip\noindent
Le lemme suivant précise la signification géométrique des séries linéaires sur les courbes propres non singulières.

\begin{lemma}
\label{lemma-linear-series}
Soit $k$ un corps. Soit $X$ une courbe propre non singulière sur $k$. Soit $(\mathcal{L}, V)$ une $\mathfrak g^r_d$ sur $X$. Il existe alors un morphisme $$
\varphi : X \longrightarrow \mathbf{P}^r_k = \text{Proj}(k[T_0, \ldots, T_r])
$$ de variétés sur $k$ et une application $\alpha : \varphi^*\mathcal{O}_{\mathbf{P}^r_k}(1) \to \mathcal{L}$ tels que $\varphi^*T_0, \ldots, \varphi^*T_r$ aient pour images une base de $V$ par $\alpha$.
\end{lemma}

\begin{proof}
Soit $s_0, \ldots, s_r \in V$ une base sur $k$. Puisque $X$ est non singulière, l'image $\mathcal{L}' \subset \mathcal{L}$ de l'application $s_0, \ldots, s_r : \mathcal{O}_X^{\oplus r + 1} \to \mathcal{L}$ est un $\mathcal{O}_X$-module inversible, par exemple d'après Diviseurs, lemme \ref{divisors-lemma-torsion-free-over-regular-dim-1}. On utilise ensuite Constructions, lemme \ref{constructions-lemma-projective-space} pour obtenir un morphisme $$
\varphi = \varphi_{(\mathcal{L}', (s_0, \ldots, s_r))} :
X \longrightarrow \mathbf{P}^r_k
$$ comme dans l'énoncé du lemme.
\end{proof}

\begin{lemma}
\label{lemma-linear-series-trivial-existence}
Soit $k$ un corps. Soit $X$ un schéma propre de dimension $\leq 1$ sur $k$. Si $X$ possède une $\mathfrak g^r_d$, alors $X$ possède une $\mathfrak g^s_d$ pour tout $0 \leq s \leq r$.
\end{lemma}

\begin{proof}
En effet, un espace vectoriel $V$ de dimension $r + 1$ sur $k$ possède un sous-espace vectoriel de dimension $s + 1$ pour tout $0 \leq s \leq r$.
\end{proof}

\begin{lemma}
\label{lemma-g1d}
Soit $k$ un corps. Soit $X$ une courbe propre non singulière sur $k$. Soit $(\mathcal{L}, V)$ une $\mathfrak g^1_d$ sur $X$. Alors le morphisme $\varphi : X \to \mathbf{P}^1_k$ du lemme \ref{lemma-linear-series} vérifie l'une des deux propriétés suivantes : \begin{enumerate} \item il est non constant et de degré $\leq d$; \item il se factorise par un point fermé de $\mathbf{P}^1_k$ et, dans ce cas, $H^0(X, \mathcal{O}_X) \not = k$.
\end{enumerate}
\end{lemma}

\begin{proof}
D'après le lemme \ref{lemma-linear-series}, on voit que $\mathcal{L}' = \varphi^*\mathcal{O}_{\mathbf{P}^1_k}(1)$ admet une application non nulle $\mathcal{L}' \to \mathcal{L}$. Ainsi, d'après Variétés, lemme \ref{varieties-lemma-check-invertible-sheaf-trivial}, on a $0 \leq \deg(\mathcal{L}') \leq d$. Si $\deg(\mathcal{L}') = 0$, le même lemme donne $\mathcal{L}' \cong \mathcal{O}_X$; comme on dispose de deux sections linéairement indépendantes, on est dans le cas (2). Si $\deg(\mathcal{L}') > 0$, alors $\varphi$ est non constant (car l'image réciproque d'un module inversible par un morphisme constant est triviale). On a donc $$
\deg(\mathcal{L}') =
\deg(X/\mathbf{P}^1_k) \deg(\mathcal{O}_{\mathbf{P}^1_k}(1))
$$ d'après Variétés, lemme \ref{varieties-lemma-degree-pullback-map-proper-curves}. Cela achève la démonstration, puisque le degré de $\mathcal{O}_{\mathbf{P}^1_k}(1)$ est $1$.
\end{proof}

\begin{lemma}
\label{lemma-grd-inequalities}
{\emergencystretch=2em
Soit $k$ un corps. Soit $X$ une courbe propre sur $k$ telle que $H^0(X, \mathcal{O}_X) = k$. Si $X$ possède une $\mathfrak g^r_d$, alors $r \leq d$. En cas d'égalité, on a $H^1(X, \mathcal{O}_X) = 0$, c'est-à-dire que le genre de $X$ (définition \ref{definition-genus}) est $0$.
\par}
\end{lemma}

\begin{proof}
Soit $(\mathcal{L}, V)$ une $\mathfrak g^r_d$. Puisque cela ne fait qu'augmenter $r$, on peut supposer $V = H^0(X, \mathcal{L})$. Choisissons un élément non nul $s \in V$. Le schéma des zéros de $s$ est alors un diviseur de Cartier effectif $D \subset X$, on a $\mathcal{L} = \mathcal{O}_X(D)$ et l'on dispose d'une suite exacte courte $$
0 \to \mathcal{O}_X \to \mathcal{L} \to \mathcal{L}|_D \to 0
$$; voir Diviseurs, lemme \ref{divisors-lemma-characterize-OD} et remarque \ref{divisors-remark-ses-regular-section}. D'après Variétés, lemme \ref{varieties-lemma-degree-effective-Cartier-divisor}, on a $\deg(D) = \deg(\mathcal{L}) = d$. Puisque $D$ est un schéma artinien, on a $\mathcal{L}|_D \cong \mathcal{O}_D$\footnote{Dans notre cas, cela résulte de Diviseurs, lemme \ref{divisors-lemma-finite-trivialize-invertible-upstairs}, car $D \to \Spec(k)$ est fini.}. Ainsi $$
\dim_k H^0(D, \mathcal{L}|_D) = \dim_k H^0(D, \mathcal{O}_D) = \deg(D) = d
$$ D'autre part, par hypothèse, $\dim_k H^0(X, \mathcal{O}_X) = 1$ et $\dim H^0(X, \mathcal{L}) = r + 1$. On en conclut que $r + 1 \leq 1 + d$, c'est-à-dire $r \leq d$, comme dans le lemme.

\medskip\noindent
Supposons qu'il y ait égalité. Alors $H^0(X, \mathcal{L}) \to H^0(X, \mathcal{L}|_D)$ est surjective. Si nous savions que $H^1(X, \mathcal{L})$ est nul, la suite exacte longue de cohomologie donnerait la nullité de $H^1(X, \mathcal{O}_X)$ et la démonstration serait achevée. Nous allons remplacer $\mathcal{L}$ par une puissance assez grande dont la cohomologie s'annule. Puisque $\mathcal{L}|_D$ est le module inversible trivial (voir ci-dessus), on peut trouver une section $t$ de $\mathcal{L}$ dont la restriction à $D$ engendre $\mathcal{L}|_D$. Considérons l'application de multiplication $$
\mu :
H^0(X, \mathcal{L}) \otimes_k H^0(X, \mathcal{L})
\longrightarrow
H^0(X, \mathcal{L}^{\otimes 2})
$$ et la suite exacte courte $$
0 \to \mathcal{L} \xrightarrow{s}
\mathcal{L}^{\otimes 2} \to \mathcal{L}^{\otimes 2}|_D \to 0
$$ Puisque $H^0(\mathcal{L}) \to H^0(\mathcal{L}|_D)$ est surjective et que $t$ a pour image une trivialisation de $\mathcal{L}|_D$, on voit que $\mu(H^0(X, \mathcal{L}) \otimes t)$ fournit un sous-espace de $H^0(X, \mathcal{L}^{\otimes 2})$ qui se projette surjectivement sur les sections globales de $\mathcal{L}^{\otimes 2}|_D$. On a donc $$
\dim H^0(X, \mathcal{L}^{\otimes 2}) = r + 1 + d = 2r + 1 =
\deg(\mathcal{L}^{\otimes 2}) + 1
$$ Ainsi $\mathcal{L}^{\otimes 2}$ possède la même propriété que $\mathcal{L}$ : la dimension de l'espace des sections globales est égale au degré augmenté de un. Puisque $\mathcal{L}$ est ample (Variétés, lemme \ref{varieties-lemma-ample-curve}), il existe un entier $n_0$ tel que $\mathcal{L}^{\otimes n}$ ait une cohomologie $H^1$ nulle pour tout $n \geq n_0$ (Cohomologie des schémas, lemme \ref{coherent-lemma-coherent-proper-ample}). En appliquant l'argument précédent à $\mathcal{L}^{\otimes n}$ avec $n = 2^m$ pour un entier $m$ assez grand, on obtient donc le lemme.
\end{proof}

\begin{remark}[Définition classique] \label{remark-classical-linear-series} Soit $X$ une courbe projective lisse sur un corps algébriquement clos $k$. Deux diviseurs de Cartier effectifs $D, D' \subset X$ sont dits {\it linéairement équivalents} si et seulement si $\mathcal{O}_X(D) \cong \mathcal{O}_X(D')$ en tant que $\mathcal{O}_X$-modules. Puisque $\Pic(X) = \text{Cl}(X)$ (Diviseurs, lemme \ref{divisors-lemma-local-rings-UFD-c1-bijective}), on voit que $D$ et $D'$ sont linéairement équivalents si et seulement si les diviseurs de Weil associés à $D$ et $D'$ définissent le même élément de $\text{Cl}(X)$. Pour un diviseur de Cartier effectif $D \subset X$ de degré $d$, le {\it système linéaire complet}, ou la {\it série linéaire complète}, $|D|$ de $D$ est l'ensemble des diviseurs de Cartier effectifs $E \subset X$ linéairement équivalents à $D$. Autrement dit, $|D|$ est l'ensemble des points fermés de la fibre du morphisme $$
\gamma_d :
\underline{\Hilbfunctor}^d_{X/k}
\longrightarrow
\underline{\Picardfunctor}^d_{X/k}
$$ (Schémas de Picard des courbes, lemme \ref{pic-lemma-picard-pieces}) au-dessus du point fermé correspondant à $\mathcal{O}_X(D)$. Cela munit $|D|$ d'une structure naturelle de schéma, et l'on a $|D| \cong \mathbf{P}^m_k$ avec $m + 1 = h^0(\mathcal{O}_X(D))$. Plus canoniquement, on a en fait $$
|D| = \mathbf{P}(H^0(X, \mathcal{O}_X(D))^\vee)
$$, où $(-)^\vee$ désigne le dual sur $k$ et où $\mathbf{P}$ est défini dans Constructions, exemple \ref{constructions-example-projective-space}. Dans ce langage, un {\it système linéaire}, ou une {\it série linéaire}, sur $X$ est une sous-variété fermée $L \subset |D|$ définissable par des équations linéaires. Si $L$ est de dimension $r$, alors $L = \mathbf{P}(V^\vee)$, où $V \subset H^0(X, \mathcal{O}_X(D))$ est un sous-espace vectoriel de dimension $r + 1$. Ainsi la série linéaire classique $L \subset |D|$ correspond à la série linéaire $(\mathcal{O}_X(D), V)$ définie ci-dessus.
\end{remark}







\section{Dualité}
\label{section-duality}

\noindent
Dans cette section, nous explicitons, pour les schémas propres de dimension $1$ sur un corps, les conséquences des résultats très généraux sur les complexes dualisants et la dualité. Pour une discussion analogue concernant les schémas propres de dimension supérieure sur un corps, voir Dualité pour les schémas, section \ref{duality-section-duality-proper-over-field}.

\begin{lemma}
\label{lemma-duality-dim-1}
Soit $X$ un schéma propre de dimension $\leq 1$ sur un corps $k$. Il existe un complexe dualisant $\omega_X^\bullet$ possédant les propriétés suivantes : \begin{enumerate} \item $H^i(\omega_X^\bullet)$ n'est non nul que pour $i = -1, 0$; \item $\omega_X = H^{-1}(\omega_X^\bullet)$ est un module cohérent de Cohen-Macaulay dont le support est constitué des composantes irréductibles de dimension $1$; \item pour $x \in X$ fermé, le module $H^0(\omega_{X, x}^\bullet)$ est non nul si et seulement si l'une des conditions suivantes est vérifiée : \begin{enumerate} \item $\dim(\mathcal{O}_{X, x}) = 0$, ou \item $\dim(\mathcal{O}_{X, x}) = 1$ et $\mathcal{O}_{X, x}$ n'est pas de Cohen-Macaulay; \end{enumerate} \item pour $K \in D_\QCoh(\mathcal{O}_X)$, il existe des isomorphismes fonctoriels\footnote{Cette propriété caractérise $\omega_X^\bullet$ dans $D_\QCoh(\mathcal{O}_X)$ à un unique isomorphisme près, par le lemme de Yoneda. Puisque $\omega_X^\bullet$ appartient à $D^b_{\textit{Coh}}(\mathcal{O}_X)$, il suffit en fait de considérer $K \in D^b_{\textit{Coh}}(\mathcal{O}_X)$.} $$
\Ext^i_X(K, \omega_X^\bullet) = \Hom_k(H^{-i}(X, K), k)
$$ compatibles aux décalages et aux triangles distingués; \item il existe des isomorphismes fonctoriels $\Hom(\mathcal{F}, \omega_X) = \Hom_k(H^1(X, \mathcal{F}), k)$ pour $\mathcal{F}$ quasi-cohérent sur $X$; \item si $X \to \Spec(k)$ est lisse de dimension relative $1$, alors $\omega_X \cong \Omega_{X/k}$.
\end{enumerate}
\end{lemma}

\begin{proof}
Notons $f : X \to \Spec(k)$ le morphisme structural. Partons du complexe dualisant relatif $$
\omega_X^\bullet = \omega_{X/k}^\bullet = a(\mathcal{O}_{\Spec(k)})
$$ décrit dans Dualité pour les schémas, remarque \ref{duality-remark-relative-dualizing-complex}. La propriété (4) est alors vraie par construction, puisque $a$ est l'adjoint à droite de $f_* : D_\QCoh(\mathcal{O}_X) \to D(\mathcal{O}_{\Spec(k)})$. Puisque $f$ est propre, on a $f^!(\mathcal{O}_{\Spec(k)}) = a(\mathcal{O}_{\Spec(k)})$ par définition; voir Dualité pour les schémas, section \ref{duality-section-upper-shriek}. Ainsi $\omega_X^\bullet$ et $\omega_X$ sont ceux de Dualité pour les schémas, exemple \ref{duality-example-proper-over-local} et exemple \ref{duality-example-equidimensional-over-field}. Les assertions (1) et (2) résultent de Dualité pour les schémas, lemme \ref{duality-lemma-vanishing-good-dualizing}. Pour un point fermé $x \in X$, le complexe $\omega_{X, x}^\bullet$ est un complexe dualisant normalisé sur $\mathcal{O}_{X, x}$; voir Dualité pour les schémas, lemme \ref{duality-lemma-good-dualizing-normalized}. L'assertion (3) résulte alors de Complexes dualisants, lemme \ref{dualizing-lemma-apply-CM}. L'assertion (5) résulte de Dualité pour les schémas, lemme \ref{duality-lemma-dualizing-module-proper-over-A} lorsque $\mathcal{F}$ est cohérent; dans le cas général, on explicite (4) pour $K = \mathcal{F}[0]$ et $i = -1$. L'assertion (6) résulte de Dualité pour les schémas, lemme \ref{duality-lemma-smooth-proper}.
\end{proof}

\begin{lemma}
\label{lemma-duality-dim-1-CM}
{\small\emergencystretch=3em
Soit $X$ un schéma propre sur un corps $k$, de Cohen-Macaulay et équidimensionnel de dimension $1$. Le module $\omega_X$ du lemme \ref{lemma-duality-dim-1} possède les propriétés suivantes : \begin{enumerate} \item $\omega_X$ est un module dualisant sur $X$ (Dualité pour les schémas, section \ref{duality-section-dualizing-module}); \item $\omega_X$ est un module cohérent de Cohen-Macaulay dont le support est $X$; \item il existe des isomorphismes fonctoriels $\Ext^i_X(K, \omega_X[1]) = \Hom_k(H^{-i}(X, K), k)$ compatibles aux décalages pour $K \in D_\QCoh(X)$; \item il existe des isomorphismes fonctoriels $\Ext^{1 + i}(\mathcal{F}, \omega_X) = \Hom_k(H^{-i}(X, \mathcal{F}), k)$ pour $\mathcal{F}$ quasi-cohérent sur $X$.
\end{enumerate}
\par}
\end{lemma}

\begin{proof}
Rappelons que, d'après la démonstration du lemme \ref{lemma-duality-dim-1}, le module $\omega_X$ est celui de Dualité pour les schémas, exemple \ref{duality-example-proper-over-local} et qu'il est donc dualisant. Les autres assertions résultent du lemme \ref{lemma-duality-dim-1} et du fait que $\omega_X^\bullet = \omega_X[1]$, puisque $X$ est de Cohen-Macaulay (Dualité pour les schémas, lemme \ref{duality-lemma-dualizing-module-CM-scheme}).
\end{proof}

\begin{remark}
\label{remark-rework-duality-locally-free}
Soit $X$ un schéma propre de dimension $\leq 1$ sur un corps $k$. Soient $\omega_X^\bullet$ et $\omega_X$ comme dans le lemme \ref{lemma-duality-dim-1}. Si $\mathcal{E}$ est un $\mathcal{O}_X$-module fini localement libre de dual $\mathcal{E}^\vee$, on a des isomorphismes canoniques $$
\Hom_k(H^{-i}(X, \mathcal{E}), k) =
H^i(X, \mathcal{E}^\vee \otimes_{\mathcal{O}_X}^\mathbf{L} \omega_X^\bullet)
$$ Cela résulte du lemme et de Cohomologie, lemme \ref{cohomology-lemma-dual-perfect-complex}. Si $X$ est de Cohen-Macaulay et équidimensionnel de dimension $1$, on a des isomorphismes canoniques $$
\Hom_k(H^{-i}(X, \mathcal{E}), k) =
H^{1 + i}(X, \mathcal{E}^\vee \otimes_{\mathcal{O}_X} \omega_X)
$$ d'après le lemme \ref{lemma-duality-dim-1-CM}. En particulier, si $\mathcal{L}$ est un $\mathcal{O}_X$-module inversible, on a $$
\dim_k H^0(X, \mathcal{L}) =
\dim_k H^1(X, \mathcal{L}^{\otimes -1} \otimes_{\mathcal{O}_X} \omega_X)
$$ et
$$
\dim_k H^1(X, \mathcal{L}) =
\dim_k H^0(X, \mathcal{L}^{\otimes -1} \otimes_{\mathcal{O}_X} \omega_X)
$$
\end{remark}

\noindent
Vérifions une propriété de cohérence du complexe dualisant.

\begin{lemma}
\label{lemma-sanity-check-duality}
Soit $X$ un schéma propre de dimension $\leq 1$ sur un corps $k$. Soient $\omega_X^\bullet$ et $\omega_X$ comme dans le lemme \ref{lemma-duality-dim-1}. \begin{enumerate} \item Si $X \to \Spec(k)$ se factorise sous la forme $X \to \Spec(k') \to \Spec(k)$ pour un corps $k'$, alors $\omega_X^\bullet$ et $\omega_X$ vérifient les propriétés (4), (5), (6) en remplaçant $k$ par $k'$. \item Si $K/k$ est une extension de corps, les images réciproques de $\omega_X^\bullet$ et $\omega_X$ par le changement de base $X_K$ sont celles du lemme \ref{lemma-duality-dim-1} pour le morphisme $X_K \to \Spec(K)$.
\end{enumerate}
\end{lemma}

\begin{proof}
Notons $f : X \to \Spec(k)$ le morphisme structural. L'assertion (1) signifie précisément que $\omega_X^\bullet$ et $\omega_X$ sont ceux du lemme \ref{lemma-duality-dim-1} pour le morphisme $f' : X \to \Spec(k')$. Dans la démonstration du lemme \ref{lemma-duality-dim-1}, nous avons posé $\omega_X^\bullet = a(\mathcal{O}_{\Spec(k)})$, où $a$ est l'adjoint à droite de Dualité pour les schémas, lemme \ref{duality-lemma-twisted-inverse-image} pour $f$. Il faut donc montrer que $a(\mathcal{O}_{\Spec(k)}) \cong a'(\mathcal{O}_{\Spec(k)})$, où $a'$ est l'adjoint à droite de Dualité pour les schémas, lemme \ref{duality-lemma-twisted-inverse-image} pour $f'$. Puisque $k' \subset H^0(X, \mathcal{O}_X)$, l'extension $k'/k$ est finie (Cohomologie des schémas, lemme \ref{coherent-lemma-proper-over-affine-cohomology-finite}). Par unicité des adjoints, on a $a = a' \circ b$, où $b$ est l'adjoint à droite de Dualité pour les schémas, lemme \ref{duality-lemma-twisted-inverse-image} pour $g : \Spec(k') \to \Spec(k)$. Autrement dit, on a $f^! = (f')^! \circ g^!$. Il suffit donc de montrer que $\Hom_k(k', k) \cong k'$ en tant que $k'$-modules; voir Dualité pour les schémas, exemple \ref{duality-example-affine-twisted-inverse-image}. C'est le cas, car ce sont des espaces vectoriels sur $k'$ de même dimension, à savoir $1$.

\medskip\noindent
Démontrons (2). L'assertion résulte de la compatibilité de $a$ au changement de base, d'après Dualité pour les schémas, lemme \ref{duality-lemma-more-base-change}. Voir la discussion dans Dualité pour les schémas, remarque \ref{duality-remark-relative-dualizing-complex}.
\end{proof}

\begin{lemma}
\label{lemma-closed-immersion-dim-1-CM}
Soit $X$ un schéma propre de dimension $\leq 1$ sur un corps $k$. Soit $i : Y \to X$ une immersion fermée. Soient $\omega_X^\bullet$, $\omega_X$, $\omega_Y^\bullet$, $\omega_Y$ comme dans le lemme \ref{lemma-duality-dim-1}. Alors \begin{enumerate} \item $\omega_Y^\bullet = R\SheafHom(\mathcal{O}_Y, \omega_X^\bullet)$; \item $\omega_Y = \SheafHom(\mathcal{O}_Y, \omega_X)$ et $i_*\omega_Y = \SheafHom_{\mathcal{O}_X}(i_*\mathcal{O}_Y, \omega_X)$.
\end{enumerate}
\end{lemma}

\begin{proof}
Notons $g : Y \to \Spec(k)$ et $f : X \to \Spec(k)$ les morphismes structuraux. Alors $g = f \circ i$. Notons $a, b, c$ l'adjoint à droite de Dualité pour les schémas, lemme \ref{duality-lemma-twisted-inverse-image} pour $f, g, i$. Alors $b = c \circ a$ par unicité des adjoints à droite et parce que $Rg_* = Rf_* \circ Ri_*$. Dans la démonstration du lemme \ref{lemma-duality-dim-1}, nous avons posé $\omega_X^\bullet = a(\mathcal{O}_{\Spec(k)})$ et $\omega_Y^\bullet = b(\mathcal{O}_{\Spec(k)})$. Ainsi $\omega_Y^\bullet = c(\omega_X^\bullet)$, ce qui entraîne (1) d'après Dualité pour les schémas, lemme \ref{duality-lemma-twisted-inverse-image-closed}. Puisque $\omega_X = H^{-1}(\omega_X^\bullet)$ et $\omega_Y = H^{-1}(\omega_Y^\bullet)$, on en conclut que $\omega_Y = \SheafHom(\mathcal{O}_Y, \omega_X)$. Cela implique $i_*\omega_Y = \SheafHom_{\mathcal{O}_X}(i_*\mathcal{O}_Y, \omega_X)$ d'après Dualité pour les schémas, lemme \ref{duality-lemma-sheaf-with-exact-support-ext}.
\end{proof}

\begin{lemma}
\label{lemma-closed-subscheme-reduced-gorenstein}
Soit $X$ un schéma propre sur un corps $k$, de Gorenstein, réduit et équidimensionnel de dimension $1$. Soit $i : Y \to X$ un sous-schéma fermé réduit équidimensionnel de dimension $1$. Soit $j : Z \to X$ l'adhérence schématique de $X \setminus Y$. Alors \begin{enumerate} \item $Y$ et $Z$ sont de Cohen-Macaulay; \item si $\mathcal{I} \subset \mathcal{O}_X$, resp.\ $\mathcal{J} \subset \mathcal{O}_X$ est le faisceau d'idéaux de $Y$, resp.\ $Z$ dans $X$, alors $$
\mathcal{I} = i_*\mathcal{I}'
\quad\text{et}\quad
\mathcal{J} = j_*\mathcal{J}'
$$, où $\mathcal{I}' \subset \mathcal{O}_Z$, resp.\ $\mathcal{J}' \subset \mathcal{O}_Y$ est le faisceau d'idéaux de $Y \cap Z$ dans $Z$, resp.\ $Y$; \item $\omega_Y = \mathcal{J}'(i^*\omega_X)$ et $i_*(\omega_Y) = \mathcal{J}\omega_X$; \item $\omega_Z = \mathcal{I}'(i^*\omega_X)$ et $i_*(\omega_Z) = \mathcal{I}\omega_X$; \item on a les suites exactes courtes suivantes : \begin{align*}
0 \to \omega_X \to i_*i^*\omega_X \oplus j_*j^*\omega_X \to
\mathcal{O}_{Y \cap Z} \to 0 \\
0 \to i_*\omega_Y \to \omega_X \to j_*j^*\omega_X \to 0 \\
0 \to j_*\omega_Z \to \omega_X \to i_*i^*\omega_X \to 0 \\
0 \to i_*\omega_Y \oplus j_*\omega_Z \to \omega_X \to
\mathcal{O}_{Y \cap Z} \to 0 \\
0 \to \omega_Y \to i^*\omega_X \to \mathcal{O}_{Y \cap Z} \to 0 \\
0 \to \omega_Z \to j^*\omega_X \to \mathcal{O}_{Y \cap Z} \to 0
\end{align*} \end{enumerate} Ici $\omega_X$, $\omega_Y$, $\omega_Z$ sont ceux du lemme \ref{lemma-duality-dim-1}.
\end{lemma}

\begin{proof}
Un schéma noethérien réduit de dimension $1$ est de Cohen-Macaulay; (1) est donc vrai. Puisque $X$ est réduit, on a $X = Y \cup Z$ au sens schématique. Avec les notations de Morphismes, lemme \ref{morphisms-lemma-scheme-theoretic-union}, et d'après son énoncé, on a une suite exacte courte $$
0 \to \mathcal{O}_X \to
\mathcal{O}_Y \oplus \mathcal{O}_Z \to \mathcal{O}_{Y \cap Z} \to 0
$$ Puisque $\mathcal{J} = \Ker(\mathcal{O}_X \to \mathcal{O}_Z)$, $\mathcal{J}' = \Ker(\mathcal{O}_Y \to \mathcal{O}_{Y \cap Z})$, $\mathcal{I} = \Ker(\mathcal{O}_X \to \mathcal{O}_Y)$ et $\mathcal{I}' = \Ker(\mathcal{O}_Z \to \mathcal{O}_{Y \cap Z})$, une chasse au diagramme donne (2). Remarquons que $\mathcal{I} + \mathcal{J}$ est le faisceau d'idéaux de $Y \cap Z$ et que $\mathcal{I} \cap \mathcal{J} = 0$. On dispose donc des suites exactes suivantes : \begin{align*}
0 \to \mathcal{O}_X \to \mathcal{O}_Y \oplus \mathcal{O}_Z \to
\mathcal{O}_{Y \cap Z} \to 0 \\
0 \to \mathcal{J} \to \mathcal{O}_X \to \mathcal{O}_Z \to 0 \\
0 \to \mathcal{I} \to \mathcal{O}_X \to \mathcal{O}_Y \to 0 \\
0 \to \mathcal{J} \oplus \mathcal{I} \to \mathcal{O}_X \to
\mathcal{O}_{Y \cap Z} \to 0 \\
0 \to \mathcal{J}' \to \mathcal{O}_Y \to \mathcal{O}_{Y \cap Z} \to 0 \\
0 \to \mathcal{I}' \to \mathcal{O}_Z \to \mathcal{O}_{Y \cap Z} \to 0
\end{align*} Puisque $X$ est de Gorenstein, $\omega_X$ est un $\mathcal{O}_X$-module inversible (Dualité pour les schémas, lemme \ref{duality-lemma-gorenstein}). Puisque $Y \cap Z$ est de dimension $0$, on a $\omega_X|_{Y \cap Z} \cong \mathcal{O}_{Y \cap Z}$. Une fois (3) et (4) démontrés, on obtient donc les suites exactes courtes du lemme en tensorisant la suite exacte courte ci-dessus par le module inversible $\omega_X$. Par symétrie, il suffit de démontrer (3), et, d'après (2), il suffit de montrer que $i_*(\omega_Y) = \mathcal{J}\omega_X$.

\medskip\noindent
On a $i_*\omega_Y = \SheafHom_{\mathcal{O}_X}(i_*\mathcal{O}_Y, \omega_X)$ d'après le lemme \ref{lemma-closed-immersion-dim-1-CM}. En utilisant de nouveau l'inversibilité de $\omega_X$, on se ramène finalement à montrer que $\SheafHom_{\mathcal{O}_X}(\mathcal{O}_X/\mathcal{I}, \mathcal{O}_X)$ s'envoie isomorphiquement sur $\mathcal{J}$ par évaluation en $1$. Autrement dit, $\mathcal{J}$ doit être l'annulateur de $\mathcal{I}$. Cela résulte de ce qui précède.
\end{proof}







\section{Riemann-Roch}
\label{section-Riemann-Roch}

\noindent
Soit $k$ un corps. Soit $X$ un schéma propre de dimension $\leq 1$ sur $k$. Dans Variétés, section \ref{varieties-section-divisors-curves}, nous avons défini le degré d'un $\mathcal{O}_X$-module localement libre $\mathcal{E}$ de rang constant par la formule \begin{equation}
\label{equation-degree}
\deg(\mathcal{E}) =
\chi(X, \mathcal{E}) - \text{rank}(\mathcal{E})\chi(X, \mathcal{O}_X)
\end{equation}; voir Variétés, définition \ref{varieties-definition-degree-invertible-sheaf}. Dans le chapitre Homologie de Chow, nous avons défini la première classe de Chern de $\mathcal{E}$ comme une opération sur les cycles (Homologie de Chow, section \ref{chow-section-intersecting-chern-classes}) et démontré que \begin{equation}
\label{equation-degree-c1}
\deg(\mathcal{E}) = \deg(c_1(\mathcal{E}) \cap [X]_1)
\end{equation}; voir Homologie de Chow, lemme \ref{chow-lemma-degree-vector-bundle}. En combinant (\ref{equation-degree}) et (\ref{equation-degree-c1}), on obtient une première version de la formule de Riemann-Roch : \begin{equation}
\label{equation-rr}
\chi(X, \mathcal{E}) =
\deg(c_1(\mathcal{E}) \cap [X]_1) +
\text{rank}(\mathcal{E})\chi(X, \mathcal{O}_X)
\end{equation} Si $\mathcal{L}$ est un $\mathcal{O}_X$-module inversible, on peut également considérer le nombre d'intersection $(\mathcal{L} \cdot X)$ défini dans Variétés, définition \ref{varieties-definition-intersection-number}. Cela n'apporte cependant rien de nouveau, puisque \begin{equation}
\label{equation-numerical-degree}
(\mathcal{L} \cdot X) = \deg(\mathcal{L})
\end{equation} d'après Variétés, lemme \ref{varieties-lemma-intersection-numbers-and-degrees-on-curves}. Si $\mathcal{L}$ est ample, cet entier est positif et s'appelle le degré \begin{equation}
\label{equation-degree-X}
\deg_\mathcal{L}(X) = (\mathcal{L} \cdot X) = \deg(\mathcal{L})
\end{equation} de $X$ relativement à $\mathcal{L}$; voir Variétés, définition \ref{varieties-definition-degree}.

\medskip\noindent
{\emergencystretch=3em
Pour obtenir un véritable théorème de Riemann-Roch, nous voudrions écrire $\chi(X, \mathcal{O}_X)$ comme le degré d'un zéro-cycle canonique sur $X$. Nous renvoyons à \cite{F} pour une version entièrement générale. Nous utiliserons la dualité pour obtenir une formule lorsque $X$ est de Gorenstein; c'est toutefois, en un sens, un artifice (notamment parce que cette méthode ne peut pas fonctionner en dimension supérieure).
\par}

\medskip\noindent
Utilisons d'abord les lemmes \ref{lemma-duality-dim-1} et \ref{lemma-duality-dim-1-CM} pour relier la caractéristique d'Euler de $\mathcal{O}_X$ à celle du complexe dualisant ou du module dualisant.

\begin{lemma}
\label{lemma-euler}
Soit $X$ un schéma propre de dimension $\leq 1$ sur un corps $k$. Avec $\omega_X^\bullet$ et $\omega_X$ comme dans le lemme \ref{lemma-duality-dim-1}, on a $$
\chi(X, \mathcal{O}_X) = \chi(X, \omega_X^\bullet)
$$ Si $X$ est de Cohen-Macaulay et équidimensionnel de dimension $1$, alors
$$
\chi(X, \mathcal{O}_X) = - \chi(X, \omega_X)
$$
\end{lemma}

\begin{proof}
Le membre de droite de la première formule est défini comme suit : $$
\chi(X, \omega_X^\bullet) =
\sum\nolimits_{i \in \mathbf{Z}} (-1)^i\dim_k H^i(X, \omega_X^\bullet)
$$ Cette expression est bien définie puisque $\omega_X^\bullet$ appartient à $D^b_{\textit{Coh}}(\mathcal{O}_X)$, mais aussi parce que $$
H^i(X, \omega_X^\bullet) =
\Ext^i(\mathcal{O}_X, \omega_X^\bullet) =
H^{-i}(X, \mathcal{O}_X)
$$, qui est toujours de dimension finie et n'est non nul que si $i = 0, -1$. Cela démontre aussi, bien entendu, la première formule. La seconde en résulte, puisque $\omega_X^\bullet = \omega_X[1]$ dans le cas de Cohen-Macaulay; voir le lemme \ref{lemma-duality-dim-1-CM}.
\end{proof}

\noindent
Nous utiliserons le lemme \ref{lemma-euler} pour obtenir la formule voulue pour $\chi(X, \mathcal{O}_X)$ lorsque $\omega_X$ est inversible, c'est-à-dire lorsque $X$ est de Gorenstein. L'énoncé affirme que le produit de $-1/2$ de la première classe de Chern de $\omega_X$ avec le cycle $[X]_1$ associé à $X$ est un zéro-cycle naturel sur $X$ à coefficients demi-entiers, de degré $\chi(X, \mathcal{O}_X)$. L'apparition de fractions dans l'énoncé de Riemann-Roch est inévitable.

\begin{lemma}[Riemann-Roch] \label{lemma-rr} Soit $X$ un schéma propre sur un corps $k$, de Gorenstein et équidimensionnel de dimension $1$. Soit $\omega_X$ comme dans le lemme \ref{lemma-duality-dim-1}. Alors \begin{enumerate} \item $\omega_X$ est un $\mathcal{O}_X$-module inversible; \item $\deg(\omega_X) = -2\chi(X, \mathcal{O}_X)$; \item pour un $\mathcal{O}_X$-module localement libre $\mathcal{E}$ de rang constant, on a $$
\chi(X, \mathcal{E}) = \deg(\mathcal{E}) -
\textstyle{\frac{1}{2}} \text{rank}(\mathcal{E}) \deg(\omega_X)
$$ et $\dim_k(H^i(X, \mathcal{E})) =
\dim_k(H^{1 - i}(X, \mathcal{E}^\vee \otimes_{\mathcal{O}_X} \omega_X))$ pour tout $i \in \mathbf{Z}$.
\end{enumerate}
\end{lemma}

\noindent
Les courbes non singulières (normales) sont de Gorenstein; voir Dualité pour les schémas, lemme \ref{duality-lemma-regular-gorenstein}.

\begin{proof}
{\emergencystretch=2em
Rappelons que les schémas de Gorenstein sont de Cohen-Macaulay (Dualité pour les schémas, lemme \ref{duality-lemma-gorenstein-CM}); ainsi $\omega_X$ est un module dualisant sur $X$, d'après le lemme \ref{lemma-duality-dim-1-CM}. L'inversibilité du faisceau dualisant résulte essentiellement de la définition de la propriété de Gorenstein; voir Dualité pour les schémas, section \ref{duality-section-gorenstein}. En appliquant (\ref{equation-rr}) à $\omega_X$, on obtient
\par}
$$
\chi(X, \omega_X) = 
\deg(c_1(\omega_X) \cap [X]_1) + \chi(X, \mathcal{O}_X)
$$ Avec le lemme \ref{lemma-euler}, cela donne $$
2\chi(X, \mathcal{O}_X) = - \deg(c_1(\omega_X) \cap [X]_1) = - \deg(\omega_X)
$$, la seconde égalité résultant de (\ref{equation-degree-c1}). En reportant cette expression dans (\ref{equation-rr}) pour $\mathcal{E}$, on obtient la formule affichée du lemme. La symétrie des dimensions est une conséquence de la dualité pour $X$; voir la remarque \ref{remark-rework-duality-locally-free}.
\end{proof}





\section{Quelques résultats d'annulation}
\label{section-vanishing}

\noindent
Cette section rassemble quelques résultats d'annulation très faibles. On trouvera dans la section \ref{section-more-vanishing} d'autres résultats, plus intéressants.

\begin{lemma}
\label{lemma-automatic}
Soit $k$ un corps. Soit $X$ un schéma propre sur $k$ de dimension $1$ tel que $H^0(X, \mathcal{O}_X) = k$. Alors $X$ est connexe, de Cohen-Macaulay et équidimensionnel de dimension $1$.
\end{lemma}

\begin{proof}
Puisque $\Gamma(X, \mathcal{O}_X) = k$ n'a pas d'idempotent non trivial, $X$ est connexe. Cela montre déjà que $X$ est équidimensionnel de dimension $1$ (toute composante irréductible de dimension $0$ serait une composante connexe). Soit $\mathcal{I} \subset \mathcal{O}_X$ le plus grand sous-module cohérent à support dans les points fermés. Alors $\mathcal{I}$ existe (Diviseurs, lemme \ref{divisors-lemma-remove-embedded-points}) et est engendré par ses sections globales (Variétés, lemme \ref{varieties-lemma-chi-tensor-finite}). Puisque $1 \in \Gamma(X, \mathcal{O}_X)$ n'est pas une section de $\mathcal{I}$, on en conclut que $\mathcal{I} = 0$. Ainsi $X$ n'a pas de point immergé (Diviseurs, lemme \ref{divisors-lemma-remove-embedded-points}). Par conséquent, $X$ vérifie $(S_1)$ d'après Diviseurs, lemme \ref{divisors-lemma-S1-no-embedded}. Il s'ensuit que $X$ est de Cohen-Macaulay.
\end{proof}

\noindent
Dans cette section, nous nous plaçons dans la situation suivante.

\begin{situation}
\label{situation-Cohen-Macaulay-curve}
Ici $k$ est un corps et $X$ un schéma propre sur $k$ de dimension $1$ tel que $H^0(X, \mathcal{O}_X) = k$.
\end{situation}

\noindent
D'après le lemme \ref{lemma-automatic}, le schéma $X$ est de Cohen-Macaulay et équidimensionnel de dimension $1$. Le module dualisant $\omega_X$ étudié dans les lemmes \ref{lemma-duality-dim-1} et \ref{lemma-duality-dim-1-CM} a une cohomologie $H^1$ non nulle, puisqu'on a en fait $\dim_k H^1(X, \omega_X) = \dim_k H^0(X, \mathcal{O}_X) = 1$. Il se trouve que tout objet légèrement plus « positif » que $\omega_X$ a une cohomologie $H^1$ nulle.

\begin{lemma}
\label{lemma-vanishing}
Dans la situation \ref{situation-Cohen-Macaulay-curve}, soit une suite exacte $$
\omega_X \to \mathcal{F} \to \mathcal{Q} \to 0
$$ de $\mathcal{O}_X$-modules cohérents telle que $H^1(X, \mathcal{Q}) = 0$ (par exemple si $\dim(\text{Supp}(\mathcal{Q})) = 0$). Alors on a soit $H^1(X, \mathcal{F}) = 0$, soit $\mathcal{F} = \omega_X \oplus \mathcal{Q}$.
\end{lemma}

\begin{proof}
(L'assertion entre parenthèses résulte de Cohomologie des schémas, lemme \ref{coherent-lemma-coherent-support-dimension-0}.) Puisque $H^0(X, \mathcal{O}_X) = k$ est le dual de $H^1(X, \omega_X)$ (voir la section \ref{section-Riemann-Roch}), on a $\dim H^1(X, \omega_X) = 1$. Le faisceau $\omega_X$ représente le foncteur $\mathcal{F} \mapsto \Hom_k(H^1(X, \mathcal{F}), k)$ sur la catégorie des $\mathcal{O}_X$-modules cohérents (Dualité pour les schémas, lemme \ref{duality-lemma-dualizing-module-proper-over-A}). Considérons une suite exacte comme dans l'énoncé et supposons que $H^1(X, \mathcal{F}) \not = 0$. Puisque $H^1(X, \mathcal{Q}) = 0$, l'application $H^1(X, \omega_X) \to H^1(X, \mathcal{F})$ est un isomorphisme. Par la propriété universelle de $\omega_X$ rappelée ci-dessus, il existe une application $\mathcal{F} \to \omega_X$ dont l'action sur $H^1$ est l'inverse de cet isomorphisme. La composée $\omega_X \to \mathcal{F} \to \omega_X$ est l'identité (par la propriété universelle), ce qui démontre le lemme.
\end{proof}

\begin{lemma}
\label{lemma-vanishing-twist}
Dans la situation \ref{situation-Cohen-Macaulay-curve}, soit $\mathcal{L}$ un $\mathcal{O}_X$-module inversible engendré par ses sections globales et non isomorphe à $\mathcal{O}_X$. Alors $H^1(X, \omega_X \otimes \mathcal{L}) = 0$.
\end{lemma}

\begin{proof}
Par la dualité étudiée dans la section \ref{section-Riemann-Roch}, il faut montrer que $H^0(X, \mathcal{L}^{\otimes - 1}) = 0$. Sinon, on peut choisir une section globale $t$ de $\mathcal{L}^{\otimes - 1}$ et une section globale $s$ de $\mathcal{L}$ telles que $st \not = 0$. Mais alors $st$ est un multiple constant de $1$, puisque nous avons supposé que $H^0(X, \mathcal{O}_X) = k$. Il s'ensuit que $\mathcal{L} \cong \mathcal{O}_X$, d'où une contradiction.
\end{proof}

\begin{lemma}
\label{lemma-globally-generated}
Dans la situation \ref{situation-Cohen-Macaulay-curve}, soit une suite exacte $$
\omega_X \to \mathcal{F} \to \mathcal{Q} \to 0
$$ de $\mathcal{O}_X$-modules cohérents telle que $\dim(\text{Supp}(\mathcal{Q})) = 0$ et $\dim_k H^0(X, \mathcal{Q}) \geq 2$, et telle qu'il n'existe aucun sous-module non nul $\mathcal{Q}' \subset \mathcal{F}$ pour lequel $\mathcal{Q}' \to \mathcal{Q}$ soit injective. Alors le sous-module de $\mathcal{F}$ engendré par les sections globales s'envoie surjectivement sur $\mathcal{Q}$.
\end{lemma}

\begin{proof}
Soit $\mathcal{F}' \subset \mathcal{F}$ le sous-module engendré par les sections globales et par l'image de $\omega_X \to \mathcal{F}$. Puisque $\dim_k H^0(X, \mathcal{Q}) \geq 2$ et $\dim_k H^1(X, \omega_X) = \dim_k H^0(X, \mathcal{O}_X) = 1$, on voit que $\mathcal{F}' \to \mathcal{Q}$ n'est pas nul et que $\omega_X \to \mathcal{F}'$ n'est pas un isomorphisme. Ainsi $H^1(X, \mathcal{F}') = 0$ d'après le lemme \ref{lemma-vanishing} et l'hypothèse portant sur $\mathcal{F}$. Considérons la suite exacte courte $$
0 \to \mathcal{F}' \to \mathcal{F} \to
\mathcal{Q}/\Im(\mathcal{F}' \to \mathcal{Q}) \to 0
$$ Si le quotient de droite était non nul, on obtiendrait une contradiction, car $H^0(X, \mathcal{F})$ serait alors strictement supérieur à $H^0(X, \mathcal{F}')$.
\end{proof}

\noindent
Voici un exemple de résultat d'engendrement par les sections globales.

\begin{lemma}
\label{lemma-globally-generated-curve}
Dans la situation \ref{situation-Cohen-Macaulay-curve}, supposons $X$ intègre. Soit $0 \to \omega_X \to \mathcal{F} \to \mathcal{Q} \to 0$ une suite exacte courte de $\mathcal{O}_X$-modules cohérents telle que $\mathcal{F}$ soit sans torsion, $\dim(\text{Supp}(\mathcal{Q})) = 0$ et $\dim_k H^0(X, \mathcal{Q}) \geq 2$. Alors $\mathcal{F}$ est engendré par ses sections globales.
\end{lemma}

\begin{proof}
Considérons le sous-module $\mathcal{F}'$ engendré par les sections globales. D'après le lemme \ref{lemma-globally-generated}, l'application $\mathcal{F}' \to \mathcal{Q}$ est surjective; en particulier, $\mathcal{F}' \not = 0$. Puisque $X$ est une courbe, $\mathcal{F}' \subset \mathcal{F}$ est une inclusion de faisceaux de rang $1$; ainsi $\mathcal{Q}' = \mathcal{F}/\mathcal{F}'$ est à support dans un nombre fini de points. Raisonnons par l'absurde en supposant $\mathcal{Q}'$ non nul. On a alors $H^1(X, \mathcal{F}') \not = 0$. La propriété universelle donne donc une application non nulle $\mathcal{F}' \to \omega_X$ (Dualité pour les schémas, lemme \ref{duality-lemma-dualizing-module-proper-over-A}). L'image de la composée $\mathcal{F}' \to \omega_X \to \mathcal{F}$ est engendrée par des sections globales, donc contenue dans $\mathcal{F}'$. On obtient ainsi un endomorphisme non nul $\mathcal{F}' \to \mathcal{F}'$. Puisque $\mathcal{F}'$ est sans torsion de rang $1$ sur une courbe propre, cet endomorphisme est nécessairement un automorphisme (nous omettons les détails). Mais cela entraîne que $\mathcal{F}'$ est contenu dans $\omega_X \subset \mathcal{F}$, contrairement à la surjectivité de $\mathcal{F}' \to \mathcal{Q}$.
\end{proof}

\begin{lemma}
\label{lemma-tensor-omega-with-globally-generated-invertible}
Dans la situation \ref{situation-Cohen-Macaulay-curve}, soit $\mathcal{L}$ un $\mathcal{O}_X$-module inversible très ample tel que $\deg(\mathcal{L}) \geq 2$. Alors $\omega_X \otimes_{\mathcal{O}_X} \mathcal{L}$ est engendré par ses sections globales.
\end{lemma}

\begin{proof}
Supposons $k$ algébriquement clos. Soit $x \in X$ un point fermé. Soient $C_i \subset X$ les composantes irréductibles et, pour chaque $i$, soit $x_i \in C_i$ le point générique. D'après Variétés, lemme \ref{varieties-lemma-very-ample-vanish-at-point}, on peut choisir une section $s \in H^0(X, \mathcal{L})$ telle que $s$ s'annule en $x$ mais pas en $x_i$, pour tout $i$. L'application de modules correspondante $s : \mathcal{O}_X \to \mathcal{L}$ est injective, de conoyau $\mathcal{Q}$ à support dans un nombre fini de points et tel que $H^0(X, \mathcal{Q}) \geq 2$. Considérons la suite exacte correspondante $$
0 \to \omega_X \to \omega_X \otimes \mathcal{L} \to
\omega_X \otimes \mathcal{Q} \to 0
$$ D'après le lemme \ref{lemma-globally-generated}, le module engendré par les sections globales s'envoie surjectivement sur $\omega_X \otimes \mathcal{Q}$. Puisque $x$ était arbitraire, cela démontre le lemme. Nous omettons quelques détails.

\medskip\noindent
Ramenons le cas où $k$ n'est pas algébriquement clos au cas algébriquement clos. Le lecteur peut omettre la suite de la démonstration. Choisissons une clôture algébrique $\overline{k}$ de $k$ et considérons le changement de base $X_{\overline{k}}$. Vérifions que $X_{\overline{k}} \to \Spec(\overline{k})$ relève de la situation \ref{situation-Cohen-Macaulay-curve}. Par changement de base plat (Cohomologie des schémas, lemme \ref{coherent-lemma-flat-base-change-cohomology}), on a $H^0(X_{\overline{k}}, \mathcal{O}) = \overline{k}$. Le schéma $X_{\overline{k}}$ est propre sur $\overline{k}$ (Morphismes, lemme \ref{morphisms-lemma-base-change-proper}) et équidimensionnel de dimension $1$ (Morphismes, lemme \ref{morphisms-lemma-dimension-fibre-after-base-change}). L'image réciproque de $\omega_X$ sur $X_{\overline{k}}$ est le module dualisant de $X_{\overline{k}}$ d'après le lemme \ref{lemma-sanity-check-duality}. L'image réciproque de $\mathcal{L}$ sur $X_{\overline{k}}$ est très ample (Morphismes, lemme \ref{morphisms-lemma-very-ample-base-change}). Le degré de l'image réciproque de $\mathcal{L}$ sur $X_{\overline{k}}$ est égal au degré de $\mathcal{L}$ sur $X$ (Variétés, lemme \ref{varieties-lemma-degree-base-change}). Enfin, $\omega_X \otimes \mathcal{L}$ est engendré par ses sections globales si et seulement si son image par changement de base l'est (Variétés, lemme \ref{varieties-lemma-globally-generated-base-change}). Le résultat découle donc du cas où le corps de base est algébriquement clos.
\end{proof}




\section{Faisceaux inversibles très amples}
\label{section-very-ample}

\noindent
Un critère usuel pour qu'un module inversible $\mathcal{L}$ sur un schéma $X$ de type fini sur un corps algébriquement clos soit très ample est que les sections de $\mathcal{L}$ séparent les points et les vecteurs tangents (Variétés, section \ref{varieties-section-separating-points-tangent-vectors}). Voici un autre critère pour les courbes; on le comparera à Variétés, sous-section \ref{varieties-subsection-regularity}.

\begin{lemma}
\label{lemma-criterion-very-ample}
Soit $k$ un corps. Soit $X$ un schéma propre sur $k$ de dimension $1$ tel que $H^0(X, \mathcal{O}_X) = k$. Soit $\mathcal{L}$ un $\mathcal{O}_X$-module inversible. Supposons que \begin{enumerate} \item $\mathcal{L}$ possède une section globale régulière, \item $H^1(X, \mathcal{L}) = 0$, et \item $\mathcal{L}$ soit ample. \end{enumerate} Alors $\mathcal{L}^{\otimes 6}$ est très ample sur $X$ relativement à $k$.
\end{lemma}

\begin{proof}
Soit $s$ une section globale régulière de $\mathcal{L}$. Soit $i : Z = Z(s) \to X$ le schéma des zéros de $s$; voir Diviseurs, section \ref{divisors-section-effective-Cartier-invertible}. La condition (3) entraîne $Z \not = \emptyset$ (nous omettons un petit détail). Considérons la suite exacte courte $$
0 \to \mathcal{O}_X \xrightarrow{s} \mathcal{L} \to
i_*(\mathcal{L}|_Z) \to 0
$$ En tensorisant par $\mathcal{L}$, on obtient $$
0 \to \mathcal{L} \to \mathcal{L}^{\otimes 2} \to
i_*(\mathcal{L}^{\otimes 2}|_Z) \to 0
$$ Remarquons que $Z$ est de dimension $0$ (Diviseurs, lemme \ref{divisors-lemma-effective-Cartier-makes-dimension-drop}), donc est le spectre d'un anneau artinien (Variétés, lemme \ref{varieties-lemma-algebraic-scheme-dim-0}); ainsi $\mathcal{L}|_Z \cong \mathcal{O}_Z$ (Algèbre, lemme \ref{algebra-lemma-locally-free-semi-local-free}). La suite exacte courte montre aussi que $H^1(X, \mathcal{L}^{\otimes 2}) = 0$ (on peut par exemple utiliser Variétés, lemme \ref{varieties-lemma-chi-tensor-finite} pour obtenir l'annulation du terme de droite). Par récurrence sur $n \geq 1$, à l'aide de la suite $$
0 \to \mathcal{L}^{\otimes n} \xrightarrow{s}
\mathcal{L}^{\otimes n + 1} \to
i_*(\mathcal{L}^{\otimes n + 1}|_Z) \to 0
$$, on voit que $H^1(X, \mathcal{L}^{\otimes n}) = 0$ pour $n > 0$ et qu'il existe une section globale $t_{n + 1}$ de $\mathcal{L}^{\otimes n + 1}$ qui fournit une trivialisation de $\mathcal{L}^{\otimes n + 1}|_Z \cong \mathcal{O}_Z$.

\medskip\noindent
Considérons l'application de multiplication $$
\mu_n :
H^0(X, \mathcal{L}) \otimes_k H^0(X, \mathcal{L}^{\otimes n})
\oplus
H^0(X, \mathcal{L}^{\otimes 2}) \otimes_k H^0(X, \mathcal{L}^{\otimes n - 1})
\longrightarrow
H^0(X, \mathcal{L}^{\otimes n + 1})
$$ Nous affirmons qu'elle est surjective pour $n \geq 3$. Pour le voir, considérons la suite exacte courte $$
0 \to \mathcal{L}^{\otimes n} \xrightarrow{s}
\mathcal{L}^{\otimes n + 1} \to i_*(\mathcal{L}^{\otimes n + 1}|_Z) \to 0
$$ Les sections de $\mathcal{L}^{\otimes n + 1}$ provenant du terme de gauche sont dans l'image de $\mu_n$. D'autre part, puisque $H^0(\mathcal{L}^{\otimes 2}) \to H^0(\mathcal{L}^{\otimes 2}|_Z)$ est surjective (voir ci-dessus) et que $t_{n - 1}$ a pour image une trivialisation de $\mathcal{L}^{\otimes n - 1}|_Z$, $\mu_n(H^0(X, \mathcal{L}^{\otimes 2}) \otimes t_{n - 1})$ fournit un sous-espace de $H^0(X, \mathcal{L}^{\otimes n + 1})$ qui se projette surjectivement sur les sections globales de $\mathcal{L}^{\otimes n + 1}|_Z$. Cela démontre l'assertion.

\medskip\noindent
L'assertion du paragraphe précédent montre que la $k$-algèbre graduée $$
S = \bigoplus\nolimits_{n \geq 0} H^0(X, \mathcal{L}^{\otimes n})
$$ est engendrée en degrés $0, 1, 2, 3$ sur $k$. Rappelons que $X = \text{Proj}(S)$; voir Morphismes, lemme \ref{morphisms-lemma-proper-ample-is-proj}. Ainsi $S^{(6)} = \bigoplus_{n} S_{6n}$ est engendrée en degré $1$. Cela signifie que $\mathcal{L}^{\otimes 6}$ est très ample, comme voulu.
\end{proof}

\begin{lemma}
\label{lemma-criterion-very-ample-bis}
Soit $k$ un corps. Soit $X$ un schéma propre sur $k$ de dimension $1$ tel que $H^0(X, \mathcal{O}_X) = k$. Soit $\mathcal{L}$ un $\mathcal{O}_X$-module inversible. Supposons que \begin{enumerate} \item $\mathcal{L}$ soit engendré par ses sections globales, \item $H^1(X, \mathcal{L}) = 0$, et \item $\mathcal{L}$ soit ample. \end{enumerate} Alors $\mathcal{L}^{\otimes 2}$ est très ample sur $X$ relativement à $k$.
\end{lemma}

\begin{proof}
Choisissons une base $s_0, \ldots, s_n$ de $H^0(X, \mathcal{L}^{\otimes 2})$ sur $k$. La propriété (1) montre que $\mathcal{L}^{\otimes 2}$ est engendré par ses sections globales et donne un morphisme $$
\varphi_{\mathcal{L}^{\otimes 2}, (s_0, \ldots, s_n)} :
X \longrightarrow \mathbf{P}^n_k
$$ Voir Constructions, section \ref{constructions-section-projective-space}. Le lemme affirme que ce morphisme est une immersion fermée. Pour le vérifier, on peut remplacer $k$ par sa clôture algébrique; voir Descente, lemme \ref{descent-lemma-descending-property-closed-immersion}. On peut donc supposer $k$ algébriquement clos.

\medskip\noindent
Supposons $k$ algébriquement clos. Pour chaque point générique $\eta_i \in X$, soit $V_i \subset H^0(X, \mathcal{L})$ le sous-espace vectoriel sur $k$ des sections s'annulant en $\eta_i$. Puisque $\mathcal{L}$ est engendré par ses sections globales, on a $V_i \not = H^0(X, \mathcal{L})$. Puisque $X$ n'a qu'un nombre fini de composantes irréductibles et que $k$ est infini, on peut trouver $s \in H^0(X, \mathcal{L})$ ne s'annulant en $\eta_i$ pour aucun $i$. Alors $s$ est une section régulière de $\mathcal{L}$ (car $X$ est de Cohen-Macaulay d'après le lemme \ref{lemma-automatic}, et $\mathcal{L}$ n'a donc pas de point associé immergé).

\medskip\noindent
En particulier, toutes les assertions de la démonstration du lemme \ref{lemma-criterion-very-ample} sont valables pour cette section $s$. De plus, puisque $\mathcal{L}$ est engendré par ses sections globales, on peut trouver une section globale $t \in H^0(X, \mathcal{L})$ telle que $t|_Z$ ne s'annule pas (on raisonne comme ci-dessus en utilisant la finitude de l'ensemble des points de $Z$). Dans la démonstration du lemme \ref{lemma-criterion-very-ample}, on peut alors utiliser $t$ pour voir que l'application de multiplication $$
\mu_n :
H^0(X, \mathcal{L}) \otimes_k H^0(X, \mathcal{L}^{\otimes 2})
\longrightarrow
H^0(X, \mathcal{L}^{\otimes 3})
$$ est elle aussi surjective. Ainsi $$
S = \bigoplus\nolimits_{n \geq 0} H^0(X, \mathcal{L}^{\otimes n})
$$ est engendrée en degrés $0, 1, 2$ sur $k$. En raisonnant comme dans la démonstration du lemme \ref{lemma-criterion-very-ample}, on voit que $S^{(2)} = \bigoplus_{n} S_{2n}$ est engendrée en degré $1$. Cela signifie que $\mathcal{L}^{\otimes 2}$ est très ample, comme voulu. Nous omettons quelques détails.
\end{proof}









\section{Genre d'une courbe}
\label{section-genus}

\noindent
Si $X$ est une courbe projective lisse géométriquement irréductible sur un corps $k$, nous avons déjà défini le genre de $X$ comme la dimension de $H^1(X, \mathcal{O}_X)$; voir Schémas de Picard des courbes, définition \ref{pic-definition-genus}. Remarquons que $H^0(X, \mathcal{O}_X) = k$ dans ce cas; voir Variétés, lemme \ref{varieties-lemma-regular-functions-proper-variety}. Généralisons cette définition comme suit.

\begin{definition}
\label{definition-genus}
Soit $k$ un corps. Soit $X$ un schéma propre sur $k$ de dimension $1$ tel que $H^0(X, \mathcal{O}_X) = k$. Le {\it genre} de $X$ est alors $g = \dim_k H^1(X, \mathcal{O}_X)$.
\end{definition}

\noindent
On l'appelle parfois le {\it genre arithmétique} de $X$. Dans la littérature, le genre arithmétique d'une courbe propre $X$ sur $k$ est parfois défini par $$
p_a(X) = 1 - \chi(X, \mathcal{O}_X) =
1 - \dim_k H^0(X, \mathcal{O}_X) + \dim_k H^1(X, \mathcal{O}_X)
$$ Cette définition coïncide avec la nôtre lorsque celle-ci s'applique, puisque nous supposons $H^0(X, \mathcal{O}_X) = k$. Il faut toutefois noter que \begin{enumerate} \item $p_a(X)$ peut être négatif, et \item $p_a(X)$ dépend du corps de base $k$ et devrait être noté $p_a(X/k)$. \end{enumerate} Par exemple, si $k = \mathbf{Q}$ et $X = \mathbf{P}^1_{\mathbf{Q}(i)}$, alors $p_a(X/\mathbf{Q}) = -1$ et $p_a(X/\mathbf{Q}(i)) = 0$.

\medskip\noindent
L'hypothèse $H^0(X, \mathcal{O}_X) = k$ dans notre définition a deux conséquences. D'une part, elle écarte toute ambiguïté sur le corps de base. D'autre part, elle entraîne que le schéma $X$ est de Cohen-Macaulay et équidimensionnel de dimension $1$ (lemme \ref{lemma-automatic}). Si $\omega_X$ désigne le module dualisant comme dans les lemmes \ref{lemma-duality-dim-1} et \ref{lemma-duality-dim-1-CM}, on a \begin{equation}
\label{equation-genus}
g = \dim_k H^1(X, \mathcal{O}_X) = \dim_k H^0(X, \omega_X)
\end{equation} par dualité; voir la remarque \ref{remark-rework-duality-locally-free}.

\medskip\noindent
Si $X$ est propre sur $k$ de dimension $\leq 1$ et si $H^0(X, \mathcal{O}_X)$ n'est pas égal au corps de base $k$, au lieu du genre arithmétique $p_a(X)$ donné par la formule ci-dessus, nous utiliserons l'invariant $\chi(X, \mathcal{O}_X)$. En fait, il est préconisé dans \cite[page 276]{FAC} et \cite[Introduction]{Hirzebruch} d'appeler $\chi(X, \mathcal{O}_X)$ le genre arithmétique.

\begin{lemma}
\label{lemma-genus-base-change}
Soit $k'/k$ une extension de corps. Soit $X$ un schéma propre sur $k$ de dimension $1$ tel que $H^0(X, \mathcal{O}_X) = k$. Alors $X_{k'}$ est un schéma propre sur $k'$ de dimension $1$ tel que $H^0(X_{k'}, \mathcal{O}_{X_{k'}}) = k'$. De plus, le genre de $X_{k'}$ est égal au genre de $X$.
\end{lemma}

\begin{proof}
Le schéma $X_{k'}$ est de dimension $1$, par exemple d'après Morphismes, lemme \ref{morphisms-lemma-dimension-fibre-after-base-change}. Le morphisme $X_{k'} \to \Spec(k')$ est propre d'après Morphismes, lemme \ref{morphisms-lemma-base-change-proper}. L'égalité $H^0(X_{k'}, \mathcal{O}_{X_{k'}}) = k'$ résulte de Cohomologie des schémas, lemme \ref{coherent-lemma-flat-base-change-cohomology}. L'égalité des genres résulte du même lemme.
\end{proof}

\begin{lemma}
\label{lemma-genus-gorenstein}
Soit $k$ un corps. Soit $X$ un schéma propre sur $k$ de dimension $1$ tel que $H^0(X, \mathcal{O}_X) = k$. Si $X$ est de Gorenstein, alors $$
\deg(\omega_X) = 2g - 2
$$, où $g$ est le genre de $X$ et $\omega_X$ est celui du lemme \ref{lemma-duality-dim-1}.
\end{lemma}

\begin{proof}
Cela résulte immédiatement du lemme \ref{lemma-rr}.
\end{proof}

\begin{lemma}
\label{lemma-genus-smooth}
Soit $X$ une courbe propre lisse sur un corps $k$ telle que $H^0(X, \mathcal{O}_X) = k$. Alors $$
\dim_k H^0(X, \Omega_{X/k}) = g
\quad\text{et}\quad
\deg(\Omega_{X/k}) = 2g - 2
$$, où $g$ est le genre de $X$.
\end{lemma}

\begin{proof}
D'après le lemme \ref{lemma-duality-dim-1}, on a $\Omega_{X/k} = \omega_X$. Les formules résultent donc de (\ref{equation-genus}) et du lemme \ref{lemma-genus-gorenstein}.
\end{proof}






\section{Courbes planes}
\label{section-plane-curves}

\noindent
Soit $k$ un corps. Une {\it courbe plane} sera une courbe $X$ isomorphe à un sous-schéma fermé de $\mathbf{P}^2_k$. Le plongement $X \to \mathbf{P}^2_k$ sera souvent considéré comme donné. D'après Diviseurs, exemple \ref{divisors-example-closed-subscheme-of-proj}, une courbe est déterminée par l'idéal homogène correspondant $$
I(X) =
\Ker\left(
k[T_0, T_2, T_2] \longrightarrow \bigoplus \Gamma(X, \mathcal{O}_X(n))
\right)
$$ Rappelons que, dans cette situation, on a $$
X = \text{Proj}(k[T_0, T_2, T_2]/I)
$$ en tant que sous-schémas fermés de $\mathbf{P}^2_k$. Pour des renseignements plus généraux sur ces constructions, nous renvoyons à Diviseurs, exemple \ref{divisors-example-closed-subscheme-of-proj} et aux références qui y figurent. On a en fait $I(X) = (F)$ pour un polynôme homogène $F \in k[T_0, T_1, T_2]$; voir le lemme \ref{lemma-equation-plane-curve}. Puisque $X$ est irréductible, $F$ est irréductible; voir le lemme \ref{lemma-plane-curve}. De plus, en considérant la suite exacte courte $$
0 \to \mathcal{O}_{\mathbf{P}^2_k}(-d) \xrightarrow{F}
\mathcal{O}_{\mathbf{P}^2_k} \to \mathcal{O}_X \to 0
$$, où $d = \deg(F)$, on trouve que $H^0(X, \mathcal{O}_X) = k$ et que $X$ est de genre $(d - 1)(d - 2)/2$; voir la démonstration du lemme \ref{lemma-genus-plane-curve}.

\medskip\noindent
Pour trouver des courbes planes lisses, le plus simple est d'en écrire des équations explicites. Notons $p$ la caractéristique de $k$. Si $p$ ne divise pas $d$, on peut prendre $$
F = T_0^d + T_1^d + T_2^d
$$ La courbe correspondante $X = V_+(F)$ s'appelle la {\it courbe de Fermat} de degré $d$. Elle est lisse, car, sur chaque ouvert affine standard $D_+(T_i)$, on obtient une courbe isomorphe à la courbe affine $$
\Spec(k[x, y]/(x^d + y^d + 1))
$$ L'homomorphisme d'anneaux $k \to k[x, y]/(x^d + y^d + 1)$ est lisse d'après Algèbre, lemme \ref{algebra-lemma-relative-global-complete-intersection-smooth}, puisque $d x^{d - 1}$ et $d y^{d - 1}$ engendrent l'idéal unité de $k[x, y]/(x^d + y^d + 1)$. Si $p | d$ mais $p \not = 3$, on peut utiliser l'équation $$
F = T_0^{d - 1}T_1 + T_1^{d - 1}T_2 + T_2^{d - 1}T_0
$$ En effet, sur les ouverts affines, on obtient $x + x^{d - 1}y + y^{d - 1}$, dont les dérivées sont $1 - x^{d - 2}y$ et $x^{d - 1} - y^{d - 2}$; le lieu des zéros communs de ces trois expressions est vide\footnote{En effet, puisque $x^{d - 1} = y^{d - 2}$, on a $0 = x + x^{d - 1}y + y^{d - 1} =
x + 2 x^{d - 1} y$. Puisque $x \not = 0$ du fait que $1 = x^{d - 2}y$, on obtient $0 = 1 + 2x^{d - 2}y = 3$, ce qui est absurde sauf si $3 = 0$.}. Nous laissons au lecteur le soin de construire des exemples en caractéristique $3$.

\medskip\noindent
Plus généralement, pour tout corps $k$ et tous $n$ et $d$, il existe une hypersurface lisse de degré $d$ dans $\mathbf{P}^n_k$; voir par exemple \cite{Poonen}.

\medskip\noindent
Bien entendu, cette méthode ne fournit que des courbes lisses dont le genre est un nombre triangulaire. Pour obtenir des courbes lisses de genre arbitraire, on peut chercher des courbes lisses tracées sur $\mathbf{P}^1 \times \mathbf{P}^1$ (insérer ici une référence ultérieure).

\begin{lemma}
\label{lemma-equation-plane-curve}
Soit $Z \subset \mathbf{P}^2_k$ un sous-schéma fermé équidimensionnel de dimension $1$ et sans point immergé (ce qui équivaut à dire que $Z$ est de Cohen-Macaulay). Alors l'idéal $I(Z) \subset k[T_0, T_1, T_2]$ correspondant à $Z$ est principal.
\end{lemma}

\begin{proof}
C'est un cas particulier de Diviseurs, lemme \ref{divisors-lemma-equation-codim-1-in-projective-space} (voir aussi Variétés, lemme \ref{varieties-lemma-equation-codim-1-in-projective-space}). L'assertion entre parenthèses résulte du fait qu'un schéma noethérien de dimension $1$ est de Cohen-Macaulay si et seulement s'il n'a pas de point immergé; voir Diviseurs, lemme \ref{divisors-lemma-noetherian-dim-1-CM-no-embedded-points}.
\end{proof}

\begin{lemma}
\label{lemma-plane-curve}
Soit $Z \subset \mathbf{P}^2_k$ comme dans le lemme \ref{lemma-equation-plane-curve} et soit $I(Z) = (F)$ pour un certain $F \in k[T_0, T_1, T_2]$. Alors $Z$ est une courbe si et seulement si $F$ est irréductible.
\end{lemma}

\begin{proof}
Si $F$ est réductible, disons $F = F' F''$, soit $Z'$ le sous-schéma fermé de $\mathbf{P}^2_k$ défini par $F'$. Il est clair que $Z' \subset Z$ et que $Z' \not = Z$. Puisque $Z'$ est lui aussi de dimension $1$, on en conclut que soit $Z$ n'est pas réduit, soit $Z$ n'est pas irréductible. Réciproquement, écrivons $Z = \sum a_i D_i$, où $D_i$ sont les composantes irréductibles de $Z$; voir Diviseurs, lemmes \ref{divisors-lemma-codim-1-part} et \ref{divisors-lemma-codimension-1-is-effective-Cartier}. Soit $F_i \in k[T_0, T_1, T_2]$ le polynôme homogène engendrant l'idéal de $D_i$. Il est alors clair que $F$ et $\prod F_i^{a_i}$ définissent le même sous-schéma fermé de $\mathbf{P}^2_k$. Ainsi $F = \lambda \prod F_i^{a_i}$ pour un certain $\lambda \in k^*$, puisque tous deux engendrent l'idéal de $Z$. On voit donc que, si $F$ est irréductible, $Z$ est un diviseur premier, c'est-à-dire une courbe.
\end{proof}

\begin{lemma}
\label{lemma-genus-plane-curve}
Soit $Z \subset \mathbf{P}^2_k$ comme dans le lemme \ref{lemma-equation-plane-curve} et soit $I(Z) = (F)$ pour un certain $F \in k[T_0, T_1, T_2]$. Alors $H^0(Z, \mathcal{O}_Z) = k$ et le genre de $Z$ est $(d - 1)(d - 2)/2$, où $d = \deg(F)$.
\end{lemma}

\begin{proof}
Posons $S = k[T_0, T_1, T_2]$. On a une suite exacte de modules gradués $$
0 \to S(-d) \xrightarrow{F} S \to S/(F) \to 0
$$ Notons $i : Z \to \mathbf{P}^2_k$ l'immersion fermée donnée. En appliquant le foncteur exact $\widetilde{\ }$ (Constructions, lemme \ref{constructions-lemma-proj-sheaves}), on obtient $$
0 \to \mathcal{O}_{\mathbf{P}^2_k}(-d) \to
\mathcal{O}_{\mathbf{P}^2_k} \to i_*\mathcal{O}_Z \to 0
$$, car $F$ engendre l'idéal de $Z$. Les groupes de cohomologie de $\mathcal{O}_{\mathbf{P}^2_k}(-d)$ et $\mathcal{O}_{\mathbf{P}^2_k}$ sont donnés dans Cohomologie des schémas, lemme \ref{coherent-lemma-cohomology-projective-space-over-ring}. D'autre part, on a $H^q(Z, \mathcal{O}_Z) = H^q(\mathbf{P}^2_k, i_*\mathcal{O}_Z)$ d'après Cohomologie des schémas, lemme \ref{coherent-lemma-relative-affine-cohomology}. La suite exacte longue de cohomologie montre d'abord que $$
k = H^0(\mathbf{P}^2_k, \mathcal{O}_{\mathbf{P}^2_k}) \longrightarrow
H^0(Z, \mathcal{O}_Z)
$$ est un isomorphisme, puis que l'homomorphisme de bord $$
H^1(Z, \mathcal{O}_Z) \longrightarrow
H^2(\mathbf{P}^2_k, \mathcal{O}_{\mathbf{P}^2_k}(-d)) \cong
k[T_0, T_1, T_2]_{d - 3}
$$ est un isomorphisme. Comme on voit facilement que sa dimension est $(d - 1)(d - 2)/2$, la démonstration est achevée.
\end{proof}

\begin{lemma}
\label{lemma-smooth-plane-curve-point-over-separable}
Soit $Z \subset \mathbf{P}^2_k$ comme dans le lemme \ref{lemma-equation-plane-curve} et soit $I(Z) = (F)$ pour un certain $F \in k[T_0, T_1, T_2]$. Si $Z \to \Spec(k)$ est lisse en au moins un point et si $k$ est infini, il existe un point fermé $z \in Z$ dans le lieu lisse tel que $\kappa(z)/k$ soit une extension finie séparable de degré au plus $d$.
\end{lemma}

\begin{proof}
Supposons que $z' \in Z$ soit un point où $Z \to \Spec(k)$ est lisse. Quitte à renuméroter les coordonnées, on peut supposer $z'$ contenu dans $D_+(T_0)$. Posons $f = F(1, x, y) \in k[x, y]$. Alors $Z \cap D_+(X_0)$ est isomorphe au spectre de $k[x, y]/(f)$. Soient $f_x, f_y$ les dérivées partielles de $f$ par rapport à $x, y$. Puisque $z'$ est un point lisse de $Z/k$, l'un des éléments $f_x$ ou $f_y$ est non nul dans $z'$ (voir la discussion dans Algèbre, section \ref{algebra-section-smooth}). Après renumérotation des coordonnées, on peut supposer $f_y$ non nul en $z'$. Il existe donc un sous-schéma ouvert non vide $V \subset Z \cap D_{+}(X_0)$ tel que la projection $$
p : V \longrightarrow \Spec(k[x])
$$ soit \'etale. Puisque le degré de $f$ comme polynôme en $y$ est au plus $d$, les degrés des fibres de la projection $p$ sont au plus $d$ (voir la discussion dans Morphismes, section \ref{morphisms-section-universally-bounded}). De plus, puisque $p$ est \'etale, l'image de $p$ est un ouvert $U \subset \Spec(k[x])$. Enfin, puisque $k$ est infini, l'ensemble des points $k$-rationnels $U(k)$ de $U$ est infini, donc non vide. Choisissons un point $t \in U(k)$ et soit $z \in V$ un point d'image $t$. Alors $z$ convient.
\end{proof}





\section{Courbes de genre zéro}
\label{section-genus-zero}

\noindent
Nous aurons besoin plus loin de connaître la structure d'une courbe propre de genre zéro. Une courbe propre de Gorenstein de genre zéro est une courbe plane de degré $2$, c'est-à-dire une conique; voir le lemme \ref{lemma-genus-zero}. Une courbe propre quelconque de genre zéro s'obtient à partir d'une courbe non singulière (sur un corps plus grand) par une construction de somme amalgamée; voir le lemme \ref{lemma-genus-zero-singular}. Puisqu'une courbe non singulière est de Gorenstein, ces deux résultats couvrent tous les cas possibles.

\begin{lemma}
\label{lemma-genus-zero-pic}
Soit $X$ une courbe propre sur un corps $k$ telle que $H^0(X, \mathcal{O}_X) = k$. Si $X$ est de genre $0$, tout $\mathcal{O}_X$-module inversible $\mathcal{L}$ de degré $0$ est trivial.
\end{lemma}

\begin{proof}
En effet, on a $\dim_k H^0(X, \mathcal{L}) \geq 0 + 1 - 0 = 1$ par Riemann-Roch (lemme \ref{lemma-rr}); ainsi $\mathcal{L}$ possède une section non nulle et l'on a donc $\mathcal{L} \cong \mathcal{O}_X$ d'après Variétés, lemme \ref{varieties-lemma-check-invertible-sheaf-trivial}.
\end{proof}

\begin{lemma}
\label{lemma-genus-zero-positive-degree}
Soit $X$ une courbe propre sur un corps $k$ telle que $H^0(X, \mathcal{O}_X) = k$. Supposons $X$ de genre $0$. Soit $\mathcal{L}$ un $\mathcal{O}_X$-module inversible de degré $d > 0$. Alors \begin{enumerate} \item $\dim_k H^0(X, \mathcal{L}) = d + 1$ et $\dim_k H^1(X, \mathcal{L}) = 0$; \item $\mathcal{L}$ est très ample et définit une immersion fermée dans $\mathbf{P}^d_k$.
\end{enumerate}
\end{lemma}

\begin{proof}
Par définition du degré et du genre, on a $$
\dim_k H^0(X, \mathcal{L}) - \dim_k H^1(X, \mathcal{L}) = d + 1
$$ Soit $s$ une section non nulle de $\mathcal{L}$. Le schéma des zéros de $s$ est alors un diviseur de Cartier effectif $D \subset X$, on a $\mathcal{L} = \mathcal{O}_X(D)$ et l'on dispose d'une suite exacte courte $$
0 \to \mathcal{O}_X \to \mathcal{L} \to \mathcal{L}|_D \to 0
$$; voir Diviseurs, lemme \ref{divisors-lemma-characterize-OD} et remarque \ref{divisors-remark-ses-regular-section}. Puisque $H^1(X, \mathcal{O}_X) = 0$ par hypothèse, $H^0(X, \mathcal{L}) \to H^0(X, \mathcal{L}|_D)$ est surjective. Comme $\mathcal{L}|_D$ est engendré par ses sections globales (car $\dim(D) = 0$; voir Variétés, lemme \ref{varieties-lemma-chi-tensor-finite}), le module inversible $\mathcal{L}$ est engendré par ses sections globales. En fait, puisque $D$ est un schéma artinien, on a $\mathcal{L}|_D \cong \mathcal{O}_D$\footnote{Dans notre cas, cela résulte de Diviseurs, lemme \ref{divisors-lemma-finite-trivialize-invertible-upstairs}, car $D \to \Spec(k)$ est fini.}; on peut donc trouver une section $t$ de $\mathcal{L}$ dont la restriction à $D$ engendre $\mathcal{L}|_D$. La suite exacte courte montre aussi que $H^1(X, \mathcal{L}) = 0$.

\medskip\noindent
Pour $n \geq 1$, considérons l'application de multiplication $$
\mu_n :
H^0(X, \mathcal{L}) \otimes_k H^0(X, \mathcal{L}^{\otimes n})
\longrightarrow
H^0(X, \mathcal{L}^{\otimes n + 1})
$$ Nous affirmons qu'elle est surjective. Pour le voir, considérons la suite exacte courte $$
0 \to \mathcal{L}^{\otimes n} \xrightarrow{s}
\mathcal{L}^{\otimes n + 1} \to \mathcal{L}^{\otimes n + 1}|_D \to 0
$$ Les sections de $\mathcal{L}^{\otimes n + 1}$ provenant du terme de gauche sont dans l'image de $\mu_n$. D'autre part, puisque $H^0(\mathcal{L}) \to H^0(\mathcal{L}|_D)$ est surjective et que $t^n$ a pour image une trivialisation de $\mathcal{L}^{\otimes n}|_D$, on voit que $\mu_n(H^0(X, \mathcal{L}) \otimes t^n)$ fournit un sous-espace de $H^0(X, \mathcal{L}^{\otimes n + 1})$ qui se projette surjectivement sur les sections globales de $\mathcal{L}^{\otimes n + 1}|_D$. Cela démontre l'assertion.

\medskip\noindent
Remarquons que $\mathcal{L}$ est ample d'après Variétés, lemme \ref{varieties-lemma-ample-curve}. Ainsi Morphismes, lemme \ref{morphisms-lemma-proper-ample-is-proj} fournit un isomorphisme $$
X \longrightarrow
\text{Proj}\left(
\bigoplus\nolimits_{n \geq 0} H^0(X, \mathcal{L}^{\otimes n})\right)
$$ Puisque les applications $\mu_n$ sont surjectives pour tout $n \geq 1$, l'algèbre graduée du membre de droite est un quotient de l'algèbre symétrique sur $H^0(X, \mathcal{L})$. En choisissant une base sur $k$, notée $s_0, \ldots, s_d$, de $H^0(X, \mathcal{L})$, on voit qu'il s'agit d'un quotient d'une algèbre de polynômes en $d + 1$ variables. Puisque les quotients d'anneaux gradués correspondent à des immersions fermées des $\text{Proj}$ (Constructions, lemme \ref{constructions-lemma-surjective-graded-rings-generated-degree-1-map-proj}), on obtient une immersion fermée $X \to \mathbf{P}^d_k$. Nous omettons de vérifier que ce morphisme est celui de Constructions, lemme \ref{constructions-lemma-projective-space} associé aux sections $s_0, \ldots, s_d$ de $\mathcal{L}$.
\end{proof}

\begin{lemma}
\label{lemma-genus-zero}
Soit $X$ une courbe propre sur un corps $k$ telle que $H^0(X, \mathcal{O}_X) = k$. Si $X$ est de Gorenstein et de genre $0$, alors $X$ est isomorphe à une courbe plane de degré $2$.
\end{lemma}

\begin{proof}
Considérons le faisceau inversible $\mathcal{L} = \omega_X^{\otimes -1}$, où $\omega_X$ est celui du lemme \ref{lemma-duality-dim-1}. On a $\deg(\omega_X) = -2$ d'après le lemme \ref{lemma-genus-gorenstein}, et donc $\deg(\mathcal{L}) = 2$. D'après le lemme \ref{lemma-genus-zero-positive-degree}, le choix d'une base $s_0, s_1, s_2$ de l'espace vectoriel sur $k$ des sections globales de $\mathcal{L}$ donne une immersion fermée $$
\varphi_{(\mathcal{L}, (s_0, s_1, s_2))} :
X \longrightarrow \mathbf{P}^2_k
$$ Ainsi $X$ est une courbe plane d'un certain degré $d$. Soit $F \in k[T_0, T_1, T_2]_d$ son équation (lemme \ref{lemma-equation-plane-curve}). Puisque le genre de $X$ est $0$, on voit que $d$ vaut $1$ ou $2$ (lemme \ref{lemma-genus-plane-curve}). Remarquons que $F$ se restreint à la section nulle sur $\varphi(X)$ et que, par conséquent, $F(s_0, s_1, s_2)$ est la section nulle de $\mathcal{L}^{\otimes 2}$. Puisque $s_0, s_1, s_2$ sont linéairement indépendantes, $F$ ne peut pas être linéaire, c'est-à-dire que $d = \deg(F) \geq 2$. Ainsi $d = 2$, ce qui achève la démonstration.
\end{proof}

\begin{proposition}[Caractérisation de la droite projective] \label{proposition-projective-line} Soit $k$ un corps. Soit $X$ une courbe propre sur $k$. Les conditions suivantes sont équivalentes : \begin{enumerate} \item $X \cong \mathbf{P}^1_k$; \item $X$ est lisse et géométriquement irréductible sur $k$, $X$ est de genre $0$ et $X$ possède un module inversible de degré impair; \item $X$ est géométriquement intègre sur $k$, $X$ est de genre $0$, $X$ est de Gorenstein et $X$ possède un faisceau inversible de degré impair; \item $H^0(X, \mathcal{O}_X) = k$, $X$ est de genre $0$, $X$ est de Gorenstein et $X$ possède un faisceau inversible de degré impair; \item $X$ est géométriquement intègre sur $k$, $X$ est de genre $0$ et $X$ possède un $\mathcal{O}_X$-module inversible de degré $1$; \item $H^0(X, \mathcal{O}_X) = k$, $X$ est de genre $0$ et $X$ possède un $\mathcal{O}_X$-module inversible de degré $1$; \item $H^1(X, \mathcal{O}_X) = 0$ et $X$ possède un $\mathcal{O}_X$-module inversible de degré $1$; \item $H^1(X, \mathcal{O}_X) = 0$ et $X$ possède des points fermés $x_1, \ldots, x_n$ tels que $\mathcal{O}_{X, x_i}$ soit normal et $\gcd([\kappa(x_i) : k]) = 1$; \item ajouter d'autres conditions ici.
\end{enumerate}
\end{proposition}

\begin{proof}
Nous montrerons que chacune des conditions (2) -- (8) entraîne (1), et nous omettons de vérifier que (1) entraîne (2) -- (8).

\medskip\noindent
Supposons (2). Un schéma lisse sur $k$ est géométriquement réduit (Variétés, lemme \ref{varieties-lemma-smooth-geometrically-normal}) et régulier (Variétés, lemme \ref{varieties-lemma-smooth-regular}). Ainsi $X$ est de Gorenstein (Dualité pour les schémas, lemme \ref{duality-lemma-regular-gorenstein}). On est donc ramené à (3).

\medskip\noindent
Supposons (3). Puisque $X$ est géométriquement intègre sur $k$, on a $H^0(X, \mathcal{O}_X) = k$ d'après Variétés, lemme \ref{varieties-lemma-regular-functions-proper-variety}, et l'on est ramené à (4).

\medskip\noindent
Supposons (4). Puisque $X$ est de Gorenstein, le module dualisant $\omega_X$ du lemme \ref{lemma-duality-dim-1} est de degré $\deg(\omega_X) = -2$ d'après le lemme \ref{lemma-genus-gorenstein}. Avec l'existence supposée d'un module inversible de degré impair, cela entraîne l'existence d'un module inversible de degré $1$. On est ainsi ramené à (6).

\medskip\noindent
Supposons (5). Puisque $X$ est géométriquement intègre sur $k$, on a $H^0(X, \mathcal{O}_X) = k$ d'après Variétés, lemme \ref{varieties-lemma-regular-functions-proper-variety}, et l'on est ramené à (6).

\medskip\noindent
Supposons (6). Alors $X \cong \mathbf{P}^1_k$ d'après le lemme \ref{lemma-genus-zero-positive-degree}.

\medskip\noindent
Supposons (7). Remarquons que $\kappa = H^0(X, \mathcal{O}_X)$ est un corps fini sur $k$ d'après Variétés, lemme \ref{varieties-lemma-regular-functions-proper-variety}. Si $d = [\kappa : k] > 1$, le degré de tout faisceau inversible est divisible par $d$ et il ne peut donc exister de faisceau inversible de degré $1$. Ainsi $d = 1$ et l'on est ramené au cas (6).

\medskip\noindent
Supposons (8). Remarquons que $\kappa = H^0(X, \mathcal{O}_X)$ est un corps fini sur $k$ d'après Variétés, lemme \ref{varieties-lemma-regular-functions-proper-variety}. Puisque $\kappa \subset \kappa(x_i)$, l'hypothèse sur le pgcd des degrés entraîne $k = \kappa$. La même condition permet de trouver des entiers $a_i$ tels que $1 = \sum a_i[\kappa(x_i) : k]$. Puisque $x_i$ définit un diviseur de Cartier effectif sur $X$ d'après Variétés, lemme \ref{varieties-lemma-regular-point-on-curve}, on peut considérer le module inversible $\mathcal{O}_X(\sum a_i x_i)$. Par le choix de $a_i$, le degré de $\mathcal{L}$ est $1$. Ainsi $X \cong \mathbf{P}^1_k$ d'après le lemme \ref{lemma-genus-zero-positive-degree}.
\end{proof}

\begin{lemma}
\label{lemma-genus-zero-singular}
Soit $X$ une courbe propre sur un corps $k$ telle que $H^0(X, \mathcal{O}_X) = k$. Supposons $X$ singulière et de genre $0$. Il existe alors un diagramme $$
\xymatrix{
x' \ar[d] \ar[r] & X' \ar[d]^\nu \ar[r] & \Spec(k') \ar[d] \\
x \ar[r] & X \ar[r] & \Spec(k)
}
$$ dans lequel \begin{enumerate} \item $k'/k$ est une extension finie non triviale; \item $X' \cong \mathbf{P}^1_{k'}$; \item $x'$ est un point $k'$-rationnel de $X'$; \item $x$ est un point $k$-rationnel de $X$; \item $X' \setminus \{x'\} \to X \setminus \{x\}$ est un isomorphisme; \item $0 \to \mathcal{O}_X \to \nu_*\mathcal{O}_{X'} \to k'/k \to 0$ est une suite exacte courte, où $k'/k = \kappa(x')/\kappa(x)$ désigne le faisceau gratte-ciel au point $x$.
\end{enumerate}
\end{lemma}

\begin{proof}
Soit $\nu : X' \to X$ la normalisation de $X$; voir Variétés, sections \ref{varieties-section-normalization} et \ref{varieties-section-normalization-one-dimensional}. Puisque $X$ est singulière, $\nu$ n'est pas un isomorphisme. Alors $k' = H^0(X', \mathcal{O}_{X'})$ est une extension finie de $k$ (Variétés, lemme \ref{varieties-lemma-regular-functions-proper-variety}). La suite exacte courte $$
0 \to \mathcal{O}_X \to \nu_*\mathcal{O}_{X'} \to \mathcal{Q} \to 0
$$ et le fait que $\mathcal{Q}$ est à support dans un nombre fini de points fermés donnent \begin{enumerate} \item $H^1(X', \mathcal{O}_{X'}) = 0$, c'est-à-dire que $X'$ est de genre $0$ comme courbe sur $k'$; \item une suite exacte courte $0 \to k \to k' \to H^0(X, \mathcal{Q}) \to 0$. \end{enumerate} En particulier, $k'/k$ est une extension non triviale.

\medskip\noindent
Considérons ensuite ce qu'on appelle souvent l'{\it idéal conducteur} $$
\mathcal{I} = \SheafHom_{\mathcal{O}_X}(\nu_*\mathcal{O}_{X'}, \mathcal{O}_X)
$$ C'est un $\mathcal{O}_X$-module quasi-cohérent. On considère $\mathcal{I}$ comme un idéal de $\mathcal{O}_X$ par l'application $\varphi \mapsto \varphi(1)$. Ainsi $\mathcal{I}(U)$ est l'ensemble des éléments $f \in \mathcal{O}_X(U)$ tels que $f \left(\nu_*\mathcal{O}_{X'}(U)\right) \subset \mathcal{O}_X(U)$. Autrement dit, la condition est que $f$ annule $\mathcal{Q}$. On a donc une suite exacte qui le définit : $$
0 \to \mathcal{I} \to \mathcal{O}_X \to
\SheafHom_{\mathcal{O}_X}(\mathcal{Q}, \mathcal{Q})
$$ Soit $U \subset X$ un ouvert affine contenant le support de $\mathcal{Q}$. Alors $V = \mathcal{Q}(U) = H^0(X, \mathcal{Q})$ est un espace vectoriel sur $k$ de dimension $n - 1$. L'image de $\mathcal{O}_X(U) \to \Hom_k(V, V)$ est une sous-algèbre commutative, donc de dimension $\leq n - 1$ sur $k$ (c'est une propriété des sous-algèbres commutatives des algèbres de matrices; nous omettons les détails). On en déduit une suite exacte courte $$
0 \to \mathcal{I} \to \mathcal{O}_X \to \mathcal{A} \to 0
$$, où $\text{Supp}(\mathcal{A}) = \text{Supp}(\mathcal{Q})$ et $\dim_k H^0(X, \mathcal{A}) \leq n - 1$. D'autre part, la description $\mathcal{I} = \SheafHom_{\mathcal{O}_X}(\nu_*\mathcal{O}_{X'}, \mathcal{O}_X)$ munit $\mathcal{I}$ d'une structure de $\nu_*\mathcal{O}_{X'}$-module telle que l'inclusion $\mathcal{I} \to \nu_*\mathcal{O}_{X'}$ soit une application de $\nu_*\mathcal{O}_{X'}$-modules. On en conclut que $\mathcal{I} = \nu_*\mathcal{I}'$ pour un faisceau quasi-cohérent d'idéaux $\mathcal{I}' \subset \mathcal{O}_{X'}$; voir Morphismes, lemme \ref{morphisms-lemma-affine-equivalence-modules}. Définissons $\mathcal{A}'$ comme le conoyau : $$
0 \to \mathcal{I}' \to \mathcal{O}_{X'} \to \mathcal{A}' \to 0
$$ En combinant les suites exactes précédentes, on obtient une suite exacte courte $0 \to \mathcal{A} \to \nu_*\mathcal{A}' \to \mathcal{Q} \to 0$. L'estimation ci-dessus, jointe à $\dim_k H^0(X, \mathcal{Q}) = n - 1$, donne $$
\dim_k H^0(X', \mathcal{A}') =
\dim_k H^0(X, \mathcal{A}) + \dim_k H^0(X, \mathcal{Q}) \leq 2 n - 2
$$ Cependant, puisque $X'$ est une courbe sur $k'$, le membre de gauche est divisible par $n$ (Variétés, lemme \ref{varieties-lemma-divisible}). Comme $\mathcal{A}$ et $\mathcal{A}'$ ne peuvent être nuls, on en déduit que $\dim_k H^0(X', \mathcal{A}') = n$, ce qui signifie que $\mathcal{I}'$ est le faisceau d'idéaux d'un point $k'$-rationnel, noté $x'$. La proposition \ref{proposition-projective-line} donne $X' \cong \mathbf{P}^1_{k'}$. En revenant aux égalités précédentes, on en conclut que $\dim_k H^0(X, \mathcal{A}) = 1$. Cela signifie que $\mathcal{I}$ est le faisceau d'idéaux d'un point $k$-rationnel, noté $x$. Alors $\mathcal{A} = \kappa(x) = k$ et $\mathcal{A}' = \kappa(x') = k'$ en tant que faisceaux gratte-ciel. En comparant les suites exactes données ci-dessus, on obtient immédiatement l'assertion du lemme relative aux faisceaux structuraux.
\end{proof}

\begin{example}
\label{example-squish-on-P1}
En fait, la situation décrite dans le lemme \ref{lemma-genus-zero-singular} se présente pour toute extension finie non triviale $k'/k$. En effet, on peut considérer $$
A = \{f \in k'[x] \mid f(0) \in k \}
$$ Le spectre de $A$ est une courbe affine que l'on peut recoller au spectre de $B = k'[y]$ au moyen de l'isomorphisme $A_x \cong B_y$ envoyant $x^{-1}$ sur $y$. On obtient une courbe propre $X$ telle que $H^0(X, \mathcal{O}_X) = k$ et dont le point singulier $x$ correspond à l'idéal maximal $A \cap (x)$. La normalisation de $X$ est $\mathbf{P}^1_{k'}$, exactement comme dans le lemme.
\end{example}








\section{Genre géométrique}
\label{section-geometric-genus}

\noindent
Si $X$ est une courbe propre et {\bf lisse} sur $k$ telle que $H^0(X, \mathcal{O}_X) = k$, on appelle $$
p_g(X) = \dim_k H^0(X, \Omega_{X/k})
$$ le {\it genre géométrique} de $X$. D'après le lemme \ref{lemma-genus-smooth}, le genre géométrique de $X$ coïncide avec son genre (arithmétique). En dimension supérieure, en revanche, le genre géométrique diffère du genre arithmétique; voir la remarque \ref{remark-genus-higher-dimension}.

\medskip\noindent
Pour les courbes singulières, nous définirons le genre géométrique comme suit.

\begin{definition}
\label{definition-geometric-genus}
Soit $k$ un corps. Soit $X$ une courbe géométriquement irréductible sur $k$. Le {\it genre géométrique} de $X$ est le genre d'un modèle projectif lisse de $X$, éventuellement défini sur une extension de $k$ comme dans le lemme \ref{lemma-smooth-models}.
\end{definition}

\noindent
Si $k$ est parfait, le modèle projectif non singulier $Y$ de $X$ est lisse (lemme \ref{lemma-nonsingular-model-smooth}) et le genre géométrique de $X$ est simplement le genre de $Y$. Cela peut être faux si $k$ n'est pas parfait. Dans ce cas, on choisit une extension $K/k$ telle que le modèle projectif non singulier $Y_K$ de $(X_K)_{red}$ soit une courbe projective lisse et l'on définit le genre géométrique de $X$ comme le genre de $Y_K$. Les lemmes \ref{lemma-smooth-models} et \ref{lemma-genus-base-change} assurent que cette définition a un sens.

\begin{remark}
\label{remark-genus-higher-dimension}
Supposons que $X$ soit une variété propre lisse de dimension $d$ sur un corps algébriquement clos $k$. On définit alors souvent le {\it genre arithmétique} par $p_a(X) = (-1)^d(\chi(X, \mathcal{O}_X) - 1)$ et le {\it genre géométrique} par $p_g(X) = \dim_k H^0(X, \Omega^d_{X/k})$. Dans cette situation, le genre arithmétique et le genre géométrique ne coïncident plus, bien que l'on ait toujours $\omega_X \cong \Omega_{X/k}^d$. Par exemple, si $d = 2$, on a \begin{align*}
p_a(X) - p_g(X) & =
h^0(X, \mathcal{O}_X) - h^1(X, \mathcal{O}_X) + h^2(X, \mathcal{O}_X) - 1
- h^0(X, \Omega^2_{X/k}) \\
& =
- h^1(X, \mathcal{O}_X) + h^2(X, \mathcal{O}_X) - h^0(X, \omega_X) \\
& =
- h^1(X, \mathcal{O}_X)
\end{align*}, où $h^i(X, \mathcal{F}) = \dim_k H^i(X, \mathcal{F})$ et où la dernière égalité résulte de la dualité. Ainsi, pour une surface, la différence $p_g(X) - p_a(X)$ est toujours positive ou nulle; on l'appelle parfois l'irrégularité de la surface. Si $X = C_1 \times C_2$ est un produit de courbes projectives lisses de genres $g_1$ et $g_2$, son irrégularité est $g_1 + g_2$.
\end{remark}






\section{Riemann-Hurwitz}
\label{section-riemann-hurewitz}

\noindent
Soit $k$ un corps. Soit $f : X \to Y$ un morphisme de courbes lisses sur $k$. On obtient une suite exacte canonique $$
f^*\Omega_{Y/k} \xrightarrow{\text{d}f} \Omega_{X/k}
\longrightarrow \Omega_{X/Y} \longrightarrow 0
$$ d'après Morphismes, lemme \ref{morphisms-lemma-triangle-differentials}. Puisque $X$ et $Y$ sont lisses, les faisceaux $\Omega_{X/k}$ et $\Omega_{Y/k}$ sont des modules inversibles; voir Morphismes, lemme \ref{morphisms-lemma-smooth-omega-finite-locally-free}. Supposons la première application non nulle, c'est-à-dire $f$ génériquement \'etale; voir le lemme \ref{lemma-generically-etale}. Soit $R \subset X$ le sous-schéma fermé défini par la différente $\mathfrak{D}_f$ de $f$. D'après Discriminants, lemme \ref{discriminant-lemma-discriminant-quasi-finite-morphism-smooth}, c'est aussi le lieu d'annulation de $\text{d}f$; il s'agit d'un diviseur de Cartier effectif, et l'on obtient $$
f^*\Omega_{Y/k} \otimes_{\mathcal{O}_X} \mathcal{O}_X(R) = \Omega_{X/k}
$$ En particulier, si $X$, $Y$ sont projectives avec $k = H^0(Y, \mathcal{O}_Y) = H^0(X, \mathcal{O}_X)$ et si $X$, $Y$ sont de genres $g_X$, $g_Y$, on obtient la formule de Riemann-Hurwitz \begin{align*}
2g_X - 2 & =
\deg(\Omega_{X/k}) \\
& =
\deg(f^*\Omega_{Y/k} \otimes_{\mathcal{O}_X} \mathcal{O}_X(R)) \\
& =
\deg(f) \deg(\Omega_{Y/k}) + \deg(R) \\
& =
\deg(f) (2g_Y - 2) + \deg(R)
\end{align*} La première et la dernière égalité résultent du lemme \ref{lemma-genus-smooth}; la deuxième, de l'isomorphisme de faisceaux inversibles ci-dessus; la troisième, de l'additivité des degrés (Variétés, lemme \ref{varieties-lemma-degree-tensor-product}), de la formule du degré d'une image réciproque (Variétés, lemme \ref{varieties-lemma-degree-pullback-map-proper-curves}) et, enfin, de la formule du degré de $\mathcal{O}_X(R)$ (Variétés, lemme \ref{varieties-lemma-degree-effective-Cartier-divisor}).

\medskip\noindent
{\emergencystretch=3em
Pour utiliser la formule de Riemann-Hurwitz, il faut calculer $\deg(R) = \dim_k \Gamma(R, \mathcal{O}_R)$. La structure des schémas de dimension zéro sur $k$ (voir par exemple Variétés, lemme \ref{varieties-lemma-algebraic-scheme-dim-0}) montre que $R$ est une réunion disjointe finie de spectres d'anneaux locaux artiniens $R = \coprod_{x \in R} \Spec(\mathcal{O}_{R, x})$, chacun des $\mathcal{O}_{R, x}$ étant de dimension finie sur $k$. Ainsi
\par}
$$
\deg(R) = \sum\nolimits_{x \in R} \dim_k \mathcal{O}_{R, x} =
\sum\nolimits_{x \in R} d_x [\kappa(x) : k]
$$, où $$
d_x = \text{length}_{\mathcal{O}_{R, x}} \mathcal{O}_{R, x} =
\text{length}_{\mathcal{O}_{X, x}} \mathcal{O}_{R, x}
$$ est la multiplicité de $x$ dans $R$ (voir Algèbre, lemme \ref{algebra-lemma-pushdown-module}). Soit $x \in X$ un point fermé d'image $y \in Y$. En passant aux fibres, on obtient une suite exacte $$
\Omega_{Y/k, y} \to \Omega_{X/k, x} \to \Omega_{X/Y, x} \to 0
$$ En choisissant des générateurs locaux $\eta_x$ et $\eta_y$ des modules (libres de rang $1$) $\Omega_{X/k, x}$ et $\Omega_{Y/k, y}$, on voit que $
\eta_y \mapsto h \eta_x
$ pour un élément non nul $h \in \mathcal{O}_{X, x}$. La suite exacte montre que $\Omega_{X/Y, x} \cong \mathcal{O}_{X, x}/h\mathcal{O}_{X, x}$ en tant que $\mathcal{O}_{X, x}$-modules. Puisque le diviseur $R$ est défini par $h$ (voir ci-dessus), on a $\mathcal{O}_{R, x} = \mathcal{O}_{X, x}/h\mathcal{O}_{X, x}$. On obtient donc les égalités suivantes : \begin{align*}
d_x
& =
\text{length}_{\mathcal{O}_{X, x}}(\mathcal{O}_{R, x}) \\
& =
\text{length}_{\mathcal{O}_{X, x}}(\mathcal{O}_{X, x}/h\mathcal{O}_{X, x}) \\
& =
\text{length}_{\mathcal{O}_{X, x}}(\Omega_{X/Y, x}) \\
& =
\text{ord}_{\mathcal{O}_{X, x}}(h) \\
& =
\text{ord}_{\mathcal{O}_{X, x}}(``\eta_y/\eta_x")
\end{align*} La première résulte de la définition de $d_x$. Nous avons établi la deuxième et la troisième ci-dessus. La quatrième égalité est la définition de $\text{ord}$; voir Algèbre, définition \ref{algebra-definition-ord}. Puisque $\mathcal{O}_{X, x}$ est un anneau de valuation discrète, l'entier $\text{ord}_{\mathcal{O}_{X, x}}(h)$ est simplement la valuation de $h$. La cinquième égalité sert de moyen mnémotechnique.

\medskip\noindent
Voici un cas où l'on peut « calculer » la multiplicité $d_x$ en fonction d'autres invariants. Si $\kappa(x)$ est séparable sur $k$, on peut choisir $\eta_x = \text{d}s$ et $\eta_y = \text{d}t$, où $s$ et $t$ sont des uniformisantes de $\mathcal{O}_{X, x}$ et $\mathcal{O}_{Y, y}$ (lemme \ref{lemma-uniformizer-works}). Alors $t \mapsto u s^{e_x}$ pour une unité $u \in \mathcal{O}_{X, x}$, où $e_x$ est l'indice de ramification de l'extension $\mathcal{O}_{Y, y} \subset \mathcal{O}_{X, x}$. On obtient donc $$
\eta_y = \text{d}t = \text{d}(u s^{e_x}) =
e s^{e_x - 1} u \text{d}s + s^{e_x} \text{d}u
$$ En écrivant $\text{d}u = w \text{d}s$ pour un certain $w \in \mathcal{O}_{X, x}$, on voit que $$
``\eta_y/\eta_x" = e s^{e_x - 1} u + s^{e_x} w = (e_x u + s w)s^{e_x - 1}
$$ Son ordre d'annulation est donc $e_x - 1$, sauf si la caractéristique de $\kappa(x)$ est $p > 0$ et si $p$ divise $e_x$, auquel cas l'ordre d'annulation est $> e_x - 1$.

\medskip\noindent
En combinant tout ce qui précède, on voit que, si $k$ est de caractéristique zéro, alors $$
2g_X - 2 = (2g_Y - 2)\deg(f) +
\sum\nolimits_{x \in X} (e_x - 1)[\kappa(x) : k]
$$, où $e_x$ est l'indice de ramification de $\mathcal{O}_{X, x}$ au-dessus de $\mathcal{O}_{Y, f(x)}$. Cette formule précise est valable si et seulement si toute la ramification est modérée, c'est-à-dire lorsque les extensions résiduelles $\kappa(x)/\kappa(y)$ sont séparables et que $e_x$ est premier à la caractéristique de $k$; les arguments précédents ne suffisent cependant pas à le démontrer. Nous renvoyons au lemme \ref{lemma-rhe} et à sa démonstration.

\begin{lemma}
\label{lemma-generically-etale}
\begin{slogan}
Un morphisme de courbes lisses est séparable si et seulement s'il est étale presque partout
\end{slogan}
Soit $k$ un corps. Soit $f : X \to Y$ un morphisme de courbes lisses sur $k$. Les conditions suivantes sont équivalentes : \begin{enumerate} \item $\text{d}f : f^*\Omega_{Y/k} \to \Omega_{X/k}$ est non nulle; \item $\Omega_{X/Y}$ est à support dans un fermé propre de $X$; \item il existe un ouvert non vide $U \subset X$ tel que $f|_U : U \to Y$ soit non ramifié; \item il existe un ouvert non vide $U \subset X$ tel que $f|_U : U \to Y$ soit \'etale; \item l'extension $k(X)/k(Y)$ de corps de fonctions est finie séparable.
\end{enumerate}
\end{lemma}

\begin{proof}
Puisque $X$ et $Y$ sont lisses, les faisceaux $\Omega_{X/k}$ et $\Omega_{Y/k}$ sont des modules inversibles; voir Morphismes, lemme \ref{morphisms-lemma-smooth-omega-finite-locally-free}. La suite exacte $$
f^*\Omega_{Y/k} \longrightarrow \Omega_{X/k}
\longrightarrow \Omega_{X/Y} \longrightarrow 0
$$ de Morphismes, lemme \ref{morphisms-lemma-triangle-differentials} montre que (1) et (2) sont équivalentes entre elles et que chacune équivaut à la non-nullité de $f^*\Omega_{Y/k} \to \Omega_{X/k}$ au point générique. L'équivalence de (2) et (3) résulte de Morphismes, lemme \ref{morphisms-lemma-unramified-omega-zero}. L'équivalence de (3) et (4) résulte de Morphismes, lemme \ref{morphisms-lemma-flat-unramified-etale} et de la platitude automatique (lemme \ref{lemma-flat}). Pour l'équivalence de (5) et (4), on utilise Algèbre, lemme \ref{algebra-lemma-smooth-at-generic-point}. Nous omettons quelques détails.
\end{proof}

\begin{lemma}
\label{lemma-rh}
Soit $f : X \to Y$ un morphisme de courbes propres lisses sur un corps $k$ satisfaisant aux conditions équivalentes du lemme \ref{lemma-generically-etale}. Si $k = H^0(Y, \mathcal{O}_Y) = H^0(X, \mathcal{O}_X)$ et si $X$ et $Y$ sont de genres $g_X$ et $g_Y$, alors $$
2g_X - 2 = (2g_Y - 2) \deg(f) + \deg(R)
$$, où $R \subset X$ est le diviseur de Cartier effectif défini par la différente de $f$.
\end{lemma}

\begin{proof}
Voir la discussion ci-dessus; nous avons utilisé Discriminants, lemme \ref{discriminant-lemma-discriminant-quasi-finite-morphism-smooth}, le lemme \ref{lemma-genus-smooth} et Variétés, lemmes \ref{varieties-lemma-degree-tensor-product} et \ref{varieties-lemma-degree-pullback-map-proper-curves}.
\end{proof}

\begin{lemma}
\label{lemma-uniformizer-works}
Soit $X \to \Spec(k)$ lisse de dimension relative $1$ en un point fermé $x \in X$. Si $\kappa(x)$ est séparable sur $k$, alors, pour toute uniformisante $s$ de l'anneau de valuation discrète $\mathcal{O}_{X, x}$, l'élément $\text{d}s$ engendre librement $\Omega_{X/k, x}$ sur $\mathcal{O}_{X, x}$.
\end{lemma}

\begin{proof}
L'anneau $\mathcal{O}_{X, x}$ est un anneau de valuation discrète d'après Algèbre, lemme \ref{algebra-lemma-characterize-smooth-over-field}. Puisque $x$ est fermé, $\kappa(x)$ est fini sur $k$. Ainsi, si $\kappa(x)/k$ est séparable, toute uniformisante $s$ a une image non nulle dans $\Omega_{X/k, x} \otimes_{\mathcal{O}_{X, x}} \kappa(x)$ d'après Algèbre, lemme \ref{algebra-lemma-computation-differential}. Puisque $\Omega_{X/k, x}$ est libre de rang $1$ sur $\mathcal{O}_{X, x}$, le résultat s'ensuit.
\end{proof}

\begin{lemma}
\label{lemma-rhe}
Reprenons les notations et les hypothèses du lemme \ref{lemma-rh}. Pour un point fermé $x \in X$, soit $d_x$ la multiplicité de $x$ dans $R$. Alors $$
2g_X - 2 = (2g_Y - 2) \deg(f) + \sum\nolimits d_x [\kappa(x) : k]
$$ De plus, on a les résultats suivants : \begin{enumerate} \item $d_x = \text{length}_{\mathcal{O}_{X, x}}(\Omega_{X/Y, x})$; \item $d_x \geq e_x - 1$, où $e_x$ est l'indice de ramification de $\mathcal{O}_{X, x}$ au-dessus de $\mathcal{O}_{Y, y}$; \item $d_x = e_x - 1$ si et seulement si $\mathcal{O}_{X, x}$ est modérément ramifié au-dessus de $\mathcal{O}_{Y, y}$.
\end{enumerate}
\end{lemma}

\begin{proof}
D'après le lemme \ref{lemma-rh} et la discussion ci-dessus (où l'on a utilisé Variétés, lemme \ref{varieties-lemma-algebraic-scheme-dim-0} et Algèbre, lemme \ref{algebra-lemma-pushdown-module}), il suffit de démontrer les assertions concernant la multiplicité $d_x$ de $x$ dans $R$. L'assertion (1) a été établie dans cette discussion. Nous n'y avons démontré (2) et (3) que lorsque $\kappa(x)$ est séparable sur $k$. Dans la suite, nous traiterons (2) et (3) de façon uniforme à l'aide des résultats sur les différentes des morphismes quasi-finis de Gorenstein.

\medskip\noindent
Remarquons d'abord que $f$ est un morphisme quasi-fini de Gorenstein. Cela résulte, par exemple, du fait que $f$ est un morphisme plat quasi-fini et que $X$ est de Gorenstein (voir Dualité pour les schémas, lemme \ref{duality-lemma-flat-morphism-from-gorenstein-scheme}); cela a aussi été démontré dans la preuve de Discriminants, lemme \ref{discriminant-lemma-discriminant-quasi-finite-morphism-smooth} (utilisé ci-dessus). Ainsi $\omega_{X/Y}$ est inversible d'après Discriminants, lemme \ref{discriminant-lemma-gorenstein-quasi-finite}, et il le reste après remplacement de $X$ par des ouverts et après changement de base par un morphisme $Y' \to Y$. Nous utiliserons ce fait sans autre mention.

\medskip\noindent
Choisissons des ouverts affines $U \subset X$ et $V \subset Y$ tels que $x \in U$, $y \in V$, $f(U) \subset V$, et tels que $x$ soit le seul point de $U$ au-dessus de $y$. Écrivons $U = \Spec(A)$ et $V = \Spec(B)$. Alors $R \cap U$ est la différente de $f|_U : U \to V$. D'après Discriminants, lemme \ref{discriminant-lemma-base-change-different}, sa formation commute dans notre cas à tout changement de base. Par le choix de $U$ et $V$, on a $$
A \otimes_B \kappa(y) =
\mathcal{O}_{X, x} \otimes_{\mathcal{O}_{Y, y}} \kappa(y) =
\mathcal{O}_{X, x}/(s^{e_x})
$$, où $e_x$ est l'indice de ramification de l'énoncé. Soit $C = \mathcal{O}_{X, x}/(s^{e_x})$, considéré comme une algèbre finie sur $\kappa(y)$. Soit $\mathfrak{D}_{C/\kappa(y)}$ la différente de $C$ sur $\kappa(y)$ au sens de Discriminants, définition \ref{discriminant-definition-different}. Il suffit de montrer que $\mathfrak{D}_{C/\kappa(y)}$ est non nulle si et seulement si l'extension $\mathcal{O}_{Y, y} \subset \mathcal{O}_{X, x}$ est modérément ramifiée et que, dans ce cas, $\mathfrak{D}_{C/\kappa(y)}$ est l'idéal engendré par $s^{e_x - 1}$ dans $C$. Rappelons que la ramification modérée signifie exactement que $\kappa(x)/\kappa(y)$ est séparable et que la caractéristique de $\kappa(y)$ ne divise pas $e_x$. D'autre part, la différente de $C/\kappa(y)$ est non nulle si et seulement si $\tau_{C/\kappa(y)} \in \omega_{C/\kappa(y)}$ est non nulle. En effet, puisque $\omega_{C/\kappa(y)}$ est un $C$-module inversible (obtenu par changement de base à partir de $\omega_{A/B}$), il est libre de rang $1$, avec un générateur que nous noterons $\lambda$. Écrivons $\tau_{C/\kappa(y)} = h\lambda$ pour un certain $h \in C$. Alors $\mathfrak{D}_{C/\kappa(y)} = (h) \subset C$, d'où l'assertion. D'après Discriminants, lemme \ref{discriminant-lemma-tau-nonzero}, on a $\tau_{C/\kappa(y)} \not = 0$ si et seulement si $\kappa(x)/\kappa(y)$ est séparable et $e_x$ premier à la caractéristique. Enfin, même si $\tau_{C/\kappa(y)}$ est non nulle, on a encore $s \tau_{C/\kappa(y)} = 0$, car $s\tau_{C/\kappa(y)} : C \to \kappa(y)$ envoie $c$ sur la trace de l'opérateur nilpotent $sc$, laquelle est nulle. Ainsi $sh = 0$, puis $h \in (s^{e_x - 1})$, ce qui prouve que $\mathfrak{D}_{C/\kappa(y)} \subset (s^{e_x - 1})$ dans tous les cas. Puisque $(s^{e_x - 1}) \subset C$ est le plus petit idéal non nul, la dernière assertion est démontrée.
\end{proof}






\section{Applications inséparables}
\label{section-inseparable}

\noindent
Quelques remarques sur le comportement du genre par les applications inséparables.

\begin{lemma}
\label{lemma-dominated-by-smooth}
Soit $k$ un corps. Soit $f : X \to Y$ un morphisme surjectif de courbes sur $k$. Si $X$ est lisse sur $k$ et si $Y$ est normale, alors $Y$ est lisse sur $k$.
\end{lemma}

\begin{proof}
Soit $y \in Y$. Choisissons $x \in X$ d'image $y$. D'après Variétés, lemme \ref{varieties-lemma-flat-under-smooth}, il suffit de montrer que $f$ est plat en $x$. Cela résulte du lemme \ref{lemma-flat}.
\end{proof}

\begin{lemma}
\label{lemma-purely-inseparable}
Soit $k$ un corps de caractéristique $p > 0$. Soit $f : X \to Y$ un morphisme non constant de courbes propres non singulières sur $k$. Si l'extension $k(X)/k(Y)$ de corps de fonctions est purement inséparable, il existe une factorisation $$
X = X_0 \to X_1 \to \ldots \to X_n = Y
$$ telle que chaque $X_i$ soit une courbe propre non singulière et que $X_i \to X_{i + 1}$ soit un morphisme de degré $p$ avec $k(X_{i + 1}) \subset k(X_i)$ inséparable.
\end{lemma}

\begin{proof}
Cela résulte du théorème \ref{theorem-curves-rational-maps} et du fait qu'une extension finie purement inséparable de corps s'obtient toujours comme une succession d'extensions (inséparables) de degré $p$; voir Corps, lemme \ref{fields-lemma-finite-purely-inseparable}.
\end{proof}

\begin{lemma}
\label{lemma-inseparable-deg-p-smooth}
Soit $k$ un corps de caractéristique $p > 0$. Soit $f : X \to Y$ un morphisme non constant de courbes propres non singulières sur $k$. Si $X$ est lisse et si $k(Y) \subset k(X)$ est inséparable de degré $p$, il existe un unique isomorphisme $Y = X^{(p)}$ tel que $f$ soit $F_{X/k}$.
\end{lemma}

\begin{proof}
{\emergencystretch=2em
Le morphisme de Frobenius relatif $F_{X/k} : X \to X^{(p)}$ est construit dans Variétés, section \ref{varieties-section-frobenius}. Remarquons que $X^{(p)}$ est une courbe propre lisse sur $k$, obtenue par changement de base à partir de $X$. Le morphisme $F_{X/k}$ est de degré $p$ d'après Variétés, lemme \ref{varieties-lemma-inseparable-deg-p-smooth}. Ainsi $k(X^{(p)})$ et $k(Y)$ sont deux sous-corps de $k(X)$ tels que $[k(X) : k(Y)] = [k(X) : k(X^{(p)})] = p$. Pour démontrer le lemme, il suffit de montrer que $k(Y) = k(X^{(p)})$ dans $k(X)$. Voir le théorème \ref{theorem-curves-rational-maps}.
\par}

\medskip\noindent
Posons $K = k(X)$. Considérons l'application $\text{d} : K \to \Omega_{K/k}$. Il résulte du lemme \ref{lemma-generically-etale} que $k(Y)$ est contenu dans le noyau de $\text{d}$. D'après Variétés, lemme \ref{varieties-lemma-relative-frobenius-omega}, on voit que $k(X^{(p)})$ est dans le noyau de $\text{d}$. Puisque $X$ est une courbe lisse, $\Omega_{K/k}$ est un espace vectoriel de dimension $1$ sur $K$. Compléments d'algèbre, lemme \ref{more-algebra-lemma-p-basis} entraîne alors que $\Ker(\text{d}) = kK^p$ et que $[K : kK^p] = p$. Ainsi $k(Y) = kK^p = k(X^{(p)})$ pour des raisons de degré.
\end{proof}

\begin{lemma}
\label{lemma-purely-inseparable-smooth}
Soit $k$ un corps de caractéristique $p > 0$. Soit $f : X \to Y$ un morphisme non constant de courbes propres non singulières sur $k$. Si $X$ est lisse et si $k(Y) \subset k(X)$ est purement inséparable, il existe un unique $n \geq 0$ et un unique isomorphisme $Y = X^{(p^n)}$ tels que $f$ soit le Frobenius relatif itéré $n$ fois de $X/k$.
\end{lemma}

\begin{proof}
Le Frobenius relatif itéré $n$ fois de $X/k$ est défini dans Variétés, remarque \ref{varieties-remark-n-fold-relative-frobenius}. Le lemme s'obtient en combinant les lemmes \ref{lemma-inseparable-deg-p-smooth} et \ref{lemma-purely-inseparable}.
\end{proof}

\begin{lemma}
\label{lemma-purely-inseparable-smooth-genus}
Soit $k$ un corps de caractéristique $p > 0$. Soit $f : X \to Y$ un morphisme non constant de courbes propres non singulières sur $k$. Supposons que \begin{enumerate} \item $X$ soit lisse, \item $H^0(X, \mathcal{O}_X) = k$, \item $k(X)/k(Y)$ soit purement inséparable. \end{enumerate} Alors $Y$ est lisse, $H^0(Y, \mathcal{O}_Y) = k$ et le genre de $Y$ est égal au genre de $X$.
\end{lemma}

\begin{proof}
D'après le lemme \ref{lemma-purely-inseparable-smooth}, $Y = X^{(p^n)}$ s'obtient par changement de base à partir de $X$ par $F_{\Spec(k)}^n$. Ainsi $Y$ est lisse, et les assertions sur la cohomologie et le genre résultent du lemme \ref{lemma-genus-base-change}.
\end{proof}

\begin{example}
\label{example-inseparable}
Cet exemple montre que le genre peut changer par un morphisme purement inséparable de courbes projectives non singulières. Soit $k$ un corps de caractéristique $3$. Supposons qu'il existe un élément $a \in k$ qui ne soit pas une puissance $3$-ième. Par exemple, $k = \mathbf{F}_3(a)$ convient. Soit $X$ la courbe plane d'équation homogène $$
F = T_1^2T_0 - T_2^3 + aT_0^3
$$ comme dans la section \ref{section-plane-curves}. Sur l'ouvert affine $D_+(T_0)$, avec les coordonnées $x = T_1/T_0$ et $y = T_2/T_0$, on obtient $x^2 - y^3 + a = 0$, qui définit une courbe affine non singulière. De plus, le point à l'infini $(0 : 1: 0)$ est lisse. Ainsi $X$ est une courbe projective non singulière de genre $1$ (lemme \ref{lemma-genus-plane-curve}). Considérons d'autre part le morphisme $f : X \to \mathbf{P}^1_k$ qui, sur $D_+(T_0)$, envoie $(x, y)$ sur $x \in \mathbf{A}^1_k \subset \mathbf{P}^1_k$. Alors $f$ est un morphisme de courbes propres non singulières sur $k$ induisant une extension inséparable de corps de fonctions de degré $p = 3$, mais le genre de $X$ est $1$ et celui de $\mathbf{P}^1_k$ est $0$.
\end{example}

\begin{proposition}
\label{proposition-unwind-morphism-smooth}
Soit $k$ un corps de caractéristique $p > 0$. Soit $f : X \to Y$ un morphisme non constant de courbes propres lisses sur $k$. On peut factoriser $f$ sous la forme $$
X \longrightarrow X^{(p^n)} \longrightarrow Y
$$, où $X^{(p^n)} \to Y$ est un morphisme non constant de courbes propres lisses induisant une extension séparable $k(X^{(p^n)})/k(Y)$, où $$
X^{(p^n)} = X \times_{\Spec(k), F_{\Spec(k)}^n} \Spec(k),
$$ et où $X \to X^{(p^n)}$ est le Frobenius relatif itéré $n$ fois de $X$.
\end{proposition}

\begin{proof}
D'après Corps, lemme \ref{fields-lemma-separable-first}, il existe une extension intermédiaire $k(X)/E/k(Y)$ telle que $k(X)/E$ soit purement inséparable et $E/k(Y)$ séparable. Par le théorème \ref{theorem-curves-rational-maps}, elle correspond à une factorisation $X \to Z \to Y$ de $f$, où $Z$ est une courbe propre non singulière. On conclut en appliquant le lemme \ref{lemma-purely-inseparable-smooth} au morphisme $X \to Z$.
\end{proof}

\begin{lemma}
\label{lemma-inseparable-linear-system}
Soit $k$ un corps de caractéristique $p > 0$. Soit $X$ une courbe propre lisse sur $k$. Soit $(\mathcal{L}, V)$ une $\mathfrak g^r_d$ avec $r \geq 1$. Alors l'une des deux assertions suivantes est vraie : \begin{enumerate} \item il existe une $\mathfrak g^1_d$ dont le morphisme correspondant $X \to \mathbf{P}^1_k$ (lemme \ref{lemma-linear-series}) est génériquement \'etale (c'est-à-dire satisfait au lemme \ref{lemma-generically-etale}); \item il existe une $\mathfrak g^r_{d'}$ sur $X^{(p)}$, où $d' \leq d/p$.
\end{enumerate}
\end{lemma}

\begin{proof}
Choisissons deux éléments linéairement indépendants sur $k$, notés $s, t \in V$. Alors $f = s/t$ est la fonction rationnelle définissant le morphisme $X \to \mathbf{P}^1_k$ correspondant à la série linéaire $(\mathcal{L}, ks + kt)$. Si ce morphisme n'est pas génériquement \'etale, alors $f \in k(X^{(p)})$ d'après la proposition \ref{proposition-unwind-morphism-smooth}. Choisissons maintenant une base $s_0, \ldots, s_r$ de $V$ et soit $\mathcal{L}' \subset \mathcal{L}$ le faisceau inversible engendré par $s_0, \ldots, s_r$. Posons $f_i = s_i/s_0$ dans $k(X)$. Si, pour chaque couple $(s_0, s_i)$, on a $f_i \in k(X^{(p)})$, le morphisme $$
\varphi = \varphi_{(\mathcal{L}', (s_0, \ldots, s_r)} :
X
\longrightarrow
\mathbf{P}^r_k = \text{Proj}(k[T_0, \ldots, T_r])
$$ se factorise par $X^{(p)}$, car c'est vrai au-dessus de l'ouvert affine $D_+(T_0)$, et le morphisme sur la partie affine se prolonge à la courbe lisse $X^{(p)}$ tout entière d'après le lemme \ref{lemma-extend-over-normal-curve}. Notons cette factorisation $$
X \xrightarrow{F_{X/k}} X^{(p)} \xrightarrow{\psi} \mathbf{P}^r_k
$$ de $\varphi$. Alors $\mathcal{N} = \psi^*\mathcal{O}_{\mathbf{P}^1_k}(1)$ est un $\mathcal{O}_{X^{(p)}}$-module inversible tel que $\mathcal{L}' = F_{X/k}^*\mathcal{N}$ et tel que $\psi^*T_0, \ldots, \psi^*T_r$ soient linéairement indépendantes sur $k$ (leurs images réciproques étant $s_0, \ldots, s_r$ sur $X$). Enfin, on a $$
d = \deg(\mathcal{L}) \geq \deg(\mathcal{L}') =
\deg(F_{X/k}) \deg(\mathcal{N}) = p \deg(\mathcal{N})
$$, comme voulu. Nous avons utilisé ici Variétés, lemmes \ref{varieties-lemma-check-invertible-sheaf-trivial}, \ref{varieties-lemma-degree-pullback-map-proper-curves} et \ref{varieties-lemma-inseparable-deg-p-smooth}.
\end{proof}

\begin{lemma}
\label{lemma-point-over-separable-extension}
Soit $k$ un corps. Soit $X$ une courbe propre lisse sur $k$ telle que $H^0(X, \mathcal{O}_X) = k$ et de genre $g \geq 2$. Il existe alors un point fermé $x \in X$ tel que $\kappa(x)/k$ soit séparable de degré $\leq 2g - 2$.
\end{lemma}

\begin{proof}
Posons $\omega = \Omega_{X/k}$. D'après le lemme \ref{lemma-genus-smooth}, ce module est de degré $2g - 2$ et possède $g$ sections globales. Nous disposons donc d'une $\mathfrak g^{g - 1}_{2g - 2}$. Le lemme élémentaire \ref{lemma-linear-series-trivial-existence} donne une $\mathfrak g^1_{2g - 2}$ et le lemme \ref{lemma-g1d} fournit un morphisme $$
\varphi : X \longrightarrow \mathbf{P}^1_k
$$ d'un certain degré $d \leq 2g - 2$. Puisque $\varphi$ est plat (lemme \ref{lemma-flat}) et fini (lemme \ref{lemma-finite}), il est fini localement libre de degré $d$ (Morphismes, lemme \ref{morphisms-lemma-finite-flat}). Choisissons un point rationnel $t \in \mathbf{P}^1_k$ et un point $x \in X$ tel que $\varphi(x) = t$. Alors $$
d \geq [\kappa(x) : \kappa(t)] = [\kappa(x) : k]
$$, par exemple d'après Morphismes, lemmes \ref{morphisms-lemma-finite-locally-free-universally-bounded} et \ref{morphisms-lemma-characterize-universally-bounded}. Si $k$ est parfait (par exemple de caractéristique zéro ou fini), le lemme est donc démontré. On est ainsi ramené au cas du paragraphe suivant.

\medskip\noindent
Supposons que $k$ soit un corps infini de caractéristique $p > 0$. Comme ci-dessus, nous utiliserons le fait que $X$ possède une $\mathfrak g^{g - 1}_{2g - 2}$. La courbe propre lisse $X^{(p)}$ a le même genre que $X$. Son genre est donc $> 0$. On en conclut que $X^{(p)}$ ne possède de $\mathfrak g^{g - 1}_d$ pour aucun $d \leq g - 1$ d'après le lemme \ref{lemma-grd-inequalities}. En appliquant le lemme \ref{lemma-inseparable-linear-system} à notre $\mathfrak g^{g - 1}_{2g - 2}$ (et en remarquant que $2g - 2/p \leq g - 1$), on conclut que l'éventualité (2) ne se produit pas. On obtient donc un morphisme $$
\varphi : X \longrightarrow \mathbf{P}^1_k
$$ génériquement \'etale (au sens du lemme) et de degré $\leq 2g - 2$. Soit $U \subset X$ le sous-schéma ouvert non vide où $\varphi$ est \'etale. Alors $\varphi(U) \subset \mathbf{P}^1_k$ est un ouvert de Zariski non vide et l'on peut choisir un point $k$-rationnel $t \in \varphi(U)$, puisque $k$ est infini. Soit $u \in U$ un point tel que $\varphi(u) = t$. Alors $\kappa(u)/\kappa(t)$ est séparable (Morphismes, lemme \ref{morphisms-lemma-etale-over-field}), $\kappa(t) = k$ et $[\kappa(u) : k] \leq 2g - 2$ comme précédemment.
\end{proof}

\noindent
Le lemme suivant n'appartient pas vraiment à cette section, mais nous ne lui connaissons pas de meilleure place.

\begin{lemma}
\label{lemma-ramification-to-algebraic-closure}
Soit $X$ une courbe lisse sur un corps $k$. Soit $\overline{x} \in X_{\overline{k}}$ un point fermé d'image $x \in X$. L'indice de ramification de $\mathcal{O}_{X, x} \subset \mathcal{O}_{X_{\overline{k}}, \overline{x}}$ est le degré d'inséparabilité de $\kappa(x)/k$.
\end{lemma}

\begin{proof}
Quitte à restreindre $X$, on peut supposer qu'il existe un morphisme \'etale $\pi : X \to \mathbf{A}^1_k$; voir Morphismes, lemme \ref{morphisms-lemma-smooth-etale-over-affine-space}. On peut alors considérer le diagramme d'anneaux locaux $$
\xymatrix{
\mathcal{O}_{X_{\overline{k}}, \overline{x}} &
\mathcal{O}_{\mathbf{A}^1_{\overline{k}}, \pi(\overline{x})} \ar[l] \\
\mathcal{O}_{X, x} \ar[u] &
\mathcal{O}_{\mathbf{A}^1_k, \pi(x)} \ar[l] \ar[u]
}
$$ Les flèches horizontales sont d'indice de ramification $1$, car elles correspondent à des morphismes \'etales. De plus, l'extension $\kappa(x)/\kappa(\pi(x))$ est séparable; ainsi $\kappa(x)$ et $\kappa(\pi(x))$ ont le même degré d'inséparabilité sur $k$. Par multiplicativité des indices de ramification, il suffit de démontrer le résultat lorsque $x$ est un point de la droite affine.

\medskip\noindent
Supposons $X = \mathbf{A}^1_k$. Dans ce cas, l'anneau local de $X$ en $x$ est de la forme $$
\mathcal{O}_{X, x} = k[t]_{(P)}
$$, où $P$ est un polynôme unitaire irréductible sur $k$. Alors $P(t) = Q(t^q)$ pour un polynôme séparable $Q \in k[t]$; voir Corps, lemme \ref{fields-lemma-irreducible-polynomials}. Remarquons que $\kappa(x) = k[t]/(P)$ est de degré d'inséparabilité $q$ sur $k$. D'autre part, sur $\overline{k}$, on peut factoriser $Q(t) = \prod (t - \alpha_i)$ avec les $\alpha_i$ deux à deux distincts. Écrivons $\alpha_i = \beta_i^q$ pour un unique $\beta_i \in \overline{k}$. Notre point $\overline{x}$ correspond alors à l'un des $\beta_i$; on conclut, car l'indice de ramification de $$
k[t]_{(P)} \longrightarrow \overline{k}[t]_{(t - \beta_i)}
$$ est bien égal à $q$, puisque l'uniformisante $P$ s'envoie sur $(t - \beta_i)^q$ multiplié par une unité.
\end{proof}





\section{Sommes amalgamées}
\label{section-pushouts}

\noindent
Soit $k$ un corps. Considérons un diagramme en traits pleins $$
\xymatrix{
Z' \ar[d] \ar[r]_{i'} & X' \ar@{..>}[d]^a \\
Z \ar@{..>}[r]^i & X
}
$$ de schémas sur $k$ satisfaisant aux conditions suivantes : \begin{enumerate} \item[(a)] $X'$ est séparé de type fini sur $k$, de dimension $\leq 1$; \item[(b)] $i : Z' \to X'$ est une immersion fermée; \item[(c)] $Z'$ et $Z$ sont finis sur $\Spec(k)$; \item[(d)] $Z' \to Z$ est surjectif. \end{enumerate} Dans cette situation, toute partie finie de $X'$ est contenue dans un ouvert affine; voir Variétés, proposition \ref{varieties-proposition-finite-set-of-points-of-codim-1-in-affine}. Les hypothèses de Compléments sur les morphismes, proposition \ref{more-morphisms-proposition-pushout-along-closed-immersion-and-integral} sont donc satisfaites, et l'on obtient les propriétés suivantes : \begin{enumerate} \item la somme amalgamée $X = Z \amalg_{Z'} X'$ existe dans la catégorie des schémas; \item $i : Z \to X$ est une immersion fermée; \item $a : X' \to X$ est entier et surjectif; \item $X \to \Spec(k)$ est séparé d'après Compléments sur les morphismes, lemme \ref{more-morphisms-lemma-pushout-separated}; \item $X \to \Spec(k)$ est de type fini d'après Compléments sur les morphismes, lemmes \ref{more-morphisms-lemma-pushout-finite-type}; \item ainsi $a : X' \to X$ est fini d'après Morphismes, lemmes \ref{morphisms-lemma-finite-integral} et \ref{morphisms-lemma-permanence-finite-type}; \item si $X' \to \Spec(k)$ est propre, alors $X \to \Spec(k)$ est propre d'après Morphismes, lemme \ref{morphisms-lemma-image-proper-is-proper}.
\end{enumerate}

\noindent
Le lemme suivant admet des généralisations considérables.

\begin{lemma}
\label{lemma-complete-local-ring-pushout}
Dans la situation ci-dessus, supposons $Z = \Spec(k')$, où $k'$ est un corps, et $Z' = \Spec(k'_1 \times \ldots \times k'_n)$, où les $k'_i/k'$ sont des extensions finies de corps. Soient $x \in X$ l'image de $Z \to X$ et $x'_i \in X'$ l'image de $\Spec(k'_i) \to X'$. On a alors un diagramme de produit fibré $$
\xymatrix{
\prod\nolimits_{i = 1, \ldots, n} k'_i &
\prod\nolimits_{i = 1, \ldots, n} \mathcal{O}_{X', x'_i}^\wedge \ar[l] \\
k' \ar[u] &
\mathcal{O}_{X, x}^\wedge \ar[u] \ar[l]
}
$$ dans lequel les flèches horizontales sont les applications vers les corps résiduels.
\end{lemma}

\begin{proof}
Choisissons un voisinage ouvert affine $\Spec(A)$ de $x$ dans $X$. Soit $\Spec(A') \subset X'$ son image réciproque. Par construction, on a un diagramme de produit fibré $$
\xymatrix{
\prod\nolimits_{i = 1, \ldots, n} k'_i &
A' \ar[l] \\
k' \ar[u] &
A \ar[u] \ar[l]
}
$$ Puisque tous les objets sont finis sur $A$, ce diagramme reste un diagramme de produit fibré après complétion par rapport à l'idéal maximal $\mathfrak m \subset A$ correspondant à $x$ (Algèbre, lemme \ref{algebra-lemma-completion-flat}). Enfin, on applique Algèbre, lemme \ref{algebra-lemma-completion-finite-extension} pour identifier le complété de $A'$.
\end{proof}




\section{Identification de points et écrasement de vecteurs tangents}
\label{section-glueing-squishing}

\noindent
Dans la suite, nous noterons $k[\epsilon]$ l'algèbre des nombres duaux sur $k$ définie dans Variétés, définition \ref{varieties-definition-dual-numbers}.

\begin{lemma}
\label{lemma-no-in-between-over-k}
Soit $k$ un corps algébriquement clos. Soit $k \subset A$ une extension d'anneaux telle que $A$ possède exactement deux sous-algèbres sur $k$. Alors on a soit $A = k \times k$, soit $A = k[\epsilon]$.
\end{lemma}

\begin{proof}
L'hypothèse signifie que $k \not = A$ et que tout sous-anneau $k \subset C \subset A$ est égal à $k$ ou à $A$. Soit $t \in A$ avec $t \not \in k$. Alors $A$ est engendré par $t$ sur $k$. Ainsi $A = k[x]/I$ pour un idéal $I$. Si $I = (0)$, on dispose de la sous-algèbre $k[x^2]$, ce qui est exclu. Sinon, $I$ est engendré par un polynôme unitaire $P$. Écrivons $P = \prod_{i = 1}^d (t - a_i)$. Si $d > 2$, la sous-algèbre engendrée par $(t - a_1)(t - a_2)$ donne une contradiction. Ainsi $d = 2$. Si $a_1 \not = a_2$, alors $A = k \times k$; si $a_1 = a_2$, alors $A = k[\epsilon]$.
\end{proof}

\begin{example}[Identification de points] \label{example-glue-points} Soit $k$ un corps algébriquement clos. Soit $f : X' \to X$ un morphisme de schémas algébriques sur $k$. On dit que $X$ est obtenu en identifiant $a$ et $b$ dans $X'$ si les conditions suivantes sont satisfaites : \begin{enumerate} \item $a, b \in X'(k)$ sont des points distincts ayant pour image un même point $x \in X(k)$; \item $f$ est fini et $f^{-1}(X \setminus \{x\}) \to X \setminus \{x\}$ est un isomorphisme; \item on a une suite exacte courte $$
0 \to \mathcal{O}_X \to f_*\mathcal{O}_{X'} \xrightarrow{a - b} x_*k \to 0
$$ dont la flèche de droite envoie une section locale $h$ de $f_*\mathcal{O}_{X'}$ sur la différence $h(a) - h(b) \in k$. \end{enumerate} Dans ce cas, on a aussi une suite exacte courte $$
0 \to \mathcal{O}_X^* \to f_*\mathcal{O}_{X'}^*
\xrightarrow{ab^{-1}} x_*k^* \to 0
$$ dont la flèche de droite envoie une section locale $h$ de $f_*\mathcal{O}_{X'}^*$ sur la différence multiplicative $h(a)h(b)^{-1} \in k^*$.
\end{example}

\begin{example}[Écrasement d'un vecteur tangent] \label{example-squish-tangent-vector} Soit $k$ un corps algébriquement clos. Soit $f : X' \to X$ un morphisme de schémas algébriques sur $k$. On dit que $X$ est obtenu en écrasant le vecteur tangent $\vartheta$ dans $X'$ si les conditions suivantes sont satisfaites : \begin{enumerate} \item $\vartheta : \Spec(k[\epsilon]) \to X'$ est une immersion fermée sur $k$ telle que $f \circ \vartheta$ se factorise par un point $x \in X(k)$; \item $f$ est fini et $f^{-1}(X \setminus \{x\}) \to X \setminus \{x\}$ est un isomorphisme; \item on a une suite exacte courte $$
0 \to \mathcal{O}_X \to f_*\mathcal{O}_{X'} \xrightarrow{\vartheta} x_*k \to 0
$$ dont la flèche de droite envoie une section locale $h$ de $f_*\mathcal{O}_{X'}$ sur le coefficient de $\epsilon$ dans $\vartheta^\sharp(h) \in k[\epsilon]$. \end{enumerate} Dans ce cas, on a aussi une suite exacte courte $$
0 \to \mathcal{O}_X^* \to f_*\mathcal{O}_{X'}^*
\xrightarrow{\vartheta} x_*k \to 0
$$ dont la flèche de droite envoie une section locale $h$ de $f_*\mathcal{O}_{X'}^*$ sur $\text{d}\log(\vartheta^\sharp(h))$, où $\text{d}\log : k[\epsilon]^* \to k$ est l'homomorphisme de groupes abéliens envoyant $a + b\epsilon$ sur $b/a \in k$.
\end{example}

\begin{lemma}
\label{lemma-factor-almost-isomorphism}
Soit $k$ un corps algébriquement clos. Soit $f : X' \to X$ un morphisme fini de schémas algébriques sur $k$ tel que $\mathcal{O}_X \subset f_*\mathcal{O}_{X'}$ et tel que $f$ soit un isomorphisme hors d'un ensemble fini de points. Il existe alors une factorisation $$
X' = X_n \to X_{n - 1} \to \ldots \to X_1 \to X_0 = X
$$ telle que chaque $X_i \to X_{i - 1}$ soit soit une identification de deux points, soit l'écrasement d'un vecteur tangent (voir les exemples \ref{example-glue-points} et \ref{example-squish-tangent-vector}).
\end{lemma}

\begin{proof}
Soit $U \subset X$ le plus grand ouvert au-dessus duquel $f$ est un isomorphisme. Alors $X \setminus U = \{x_1, \ldots, x_n\}$ avec $x_i \in X(k)$. Nous considérerons les factorisations $X' \to Y \to X$ de $f$ telles que les deux morphismes soient finis et que $$
\mathcal{O}_X \subset g_*\mathcal{O}_Y \subset f_*\mathcal{O}_{X'}
$$, où $g : Y \to X$ est le morphisme donné. Par hypothèse, $\mathcal{O}_{X, x} \to (f_*\mathcal{O}_{X'})_x$ est un isomorphisme sauf si $x = x_i$ pour un certain $i$. Le conoyau $$
f_*\mathcal{O}_{X'}/\mathcal{O}_X = \bigoplus \mathcal{Q}_i
$$ est donc une somme directe de faisceaux gratte-ciel $\mathcal{Q}_i$ de supports $x_1, \ldots, x_n$. Puisque le quotient affiché est un $\mathcal{O}_X$-module cohérent, $\mathcal{Q}_i$ est de longueur finie sur $\mathcal{O}_{X, x_i}$. On peut donc raisonner par récurrence sur la somme de ces longueurs, c'est-à-dire sur la longueur du conoyau tout entier.

\medskip\noindent
Si $n > 1$, on peut définir une sous-algèbre sur $\mathcal{O}_X$, notée $\mathcal{A} \subset f_*\mathcal{O}_{X'}$, en prenant l'image réciproque de $\mathcal{Q}_1$. On obtient ainsi une factorisation non triviale et l'on conclut par récurrence.

\medskip\noindent
Supposons $n = 1$. Abrégeons la notation en posant $x = x_1$. Considérons l'extension finie de $k$-algèbres $$
A = \mathcal{O}_{X, x} \subset (f_*\mathcal{O}_{X'})_x = B
$$ Remarquons que $\mathcal{Q} = \mathcal{Q}_1$ est le faisceau gratte-ciel de valeur $B/A$. On dispose d'une sous-algèbre sur $k$, à savoir $A \subset A + \mathfrak m_A B \subset B$. Si les deux inclusions sont strictes, on obtient une factorisation non triviale et l'on conclut par récurrence comme ci-dessus. Si $A + \mathfrak m_A B = B$, alors $A = B$ par Nakayama; ainsi $f$ est un isomorphisme et il n'y a rien à démontrer. On peut donc supposer $B = A + \mathfrak m_A B$. Posons $C = B/\mathfrak m_A B$. Si $C$ possède plus de $2$ sous-algèbres sur $k$, on obtient une sous-algèbre strictement intermédiaire entre $A$ et $B$ en prenant l'image réciproque dans $B$. On peut donc supposer que $C$ possède exactement $2$ sous-algèbres sur $k$. Alors $C = k \times k$ ou $C = k[\epsilon]$ d'après le lemme \ref{lemma-no-in-between-over-k}. Dans ces deux cas, $f$ est respectivement l'identification de deux points ou l'écrasement d'un vecteur tangent.
\end{proof}

\begin{lemma}
\label{lemma-glue-points}
Soit $k$ un corps algébriquement clos. Si $f : X' \to X$ est l'identification de deux points $a, b$ comme dans l'exemple \ref{example-glue-points}, on a une suite exacte $$
k^* \to \Pic(X) \to \Pic(X') \to 0
$$ La première application est nulle si $a$ et $b$ appartiennent à des composantes connexes distinctes de $X'$, et injective si $X'$ est propre et si $a$ et $b$ appartiennent à la même composante connexe de $X'$.
\end{lemma}

\begin{proof}
L'application $\Pic(X) \to \Pic(X')$ est surjective d'après Variétés, lemme \ref{varieties-lemma-surjective-pic-birational-finite}. La suite exacte courte $$
0 \to \mathcal{O}_X^* \to f_*\mathcal{O}_{X'}^*
\xrightarrow{ab^{-1}} x_*k^* \to 0
$$ donne $$
H^0(X', \mathcal{O}_{X'}^*) \xrightarrow{ab^{-1}} k^* \to
H^1(X, \mathcal{O}_X^*) \to H^1(X, f_*\mathcal{O}_{X'}^*)
$$ On a $H^1(X, f_*\mathcal{O}_{X'}^*) \subset H^1(X', \mathcal{O}_{X'}^*)$ (par exemple par la suite spectrale de Leray; voir Cohomologie, lemme \ref{cohomology-lemma-Leray}). Le noyau de $\Pic(X) \to \Pic(X')$ est donc le conoyau de $ab^{-1} : H^0(X', \mathcal{O}_{X'}^*) \to k^*$. Si $a$ et $b$ appartiennent à des composantes connexes distinctes de $X'$, alors $ab^{-1}$ est surjective. Puisque $k$ est algébriquement clos, toute fonction régulière sur un schéma réduit connexe propre sur $k$ provient d'un élément de $k$; voir Variétés, lemme \ref{varieties-lemma-proper-geometrically-reduced-global-sections}. Ainsi $ab^{-1}$ est nulle si $X'$ est propre et si $a$ et $b$ appartiennent à la même composante connexe.
\end{proof}

\begin{lemma}
\label{lemma-squish-tangent-vector}
Soit $k$ un corps algébriquement clos. Si $f : X' \to X$ est l'écrasement d'un vecteur tangent $\vartheta$ comme dans l'exemple \ref{example-squish-tangent-vector}, on a une suite exacte $$
(k, +) \to \Pic(X) \to \Pic(X') \to 0
$$ dont la première application est injective si $X'$ est propre et réduit.
\end{lemma}

\begin{proof}
L'application $\Pic(X) \to \Pic(X')$ est surjective d'après Variétés, lemme \ref{varieties-lemma-surjective-pic-birational-finite}. La suite exacte courte $$
0 \to \mathcal{O}_X^* \to f_*\mathcal{O}_{X'}^*
\xrightarrow{\vartheta} x_*k \to 0
$$ de l'exemple \ref{example-squish-tangent-vector} donne $$
H^0(X', \mathcal{O}_{X'}^*) \xrightarrow{\vartheta} k \to
H^1(X, \mathcal{O}_X^*) \to H^1(X, f_*\mathcal{O}_{X'}^*)
$$ On a $H^1(X, f_*\mathcal{O}_{X'}^*) \subset H^1(X', \mathcal{O}_{X'}^*)$ (par exemple par la suite spectrale de Leray; voir Cohomologie, lemme \ref{cohomology-lemma-Leray}). Le noyau de $\Pic(X) \to \Pic(X')$ est donc le conoyau de l'application $\vartheta : H^0(X', \mathcal{O}_{X'}^*) \to k$. Puisque $k$ est algébriquement clos, toute fonction régulière sur un schéma réduit connexe propre sur $k$ provient d'un élément de $k$; voir Variétés, lemme \ref{varieties-lemma-proper-geometrically-reduced-global-sections}. D'où la dernière assertion du lemme.
\end{proof}


\section{Singularités en croisement multiple et singularités nodales}
\label{section-multicross}

\noindent
Dans cette section, nous étudions les singularités de courbes les plus simples possibles.

\medskip\noindent
Soit $k$ un corps. Considérons la $k$-algèbre locale complète \begin{equation}
\label{equation-multicross}
A = \{(f_1, \ldots, f_n) \in k[[t]] \times \ldots \times k[[t]] \mid
f_1(0) = \ldots = f_n(0)\}
\end{equation} Dans la terminologie introduite dans Variétés, définition \ref{varieties-definition-wedge}, $A$ est un bouquet de $n$ copies de l'anneau des séries formelles en $1$ variable sur $k$. Remarquons que $k[[t]] \times \ldots \times k[[t]]$ est la clôture intégrale de $A$ dans son anneau total des fractions. L'invariant $\delta$ de $A$ est donc $n - 1$. On a un isomorphisme $$
k[[x_1, \ldots, x_n]]/(\{x_ix_j\}_{i \not = j}) \longrightarrow A
$$ obtenu en envoyant $x_i$ sur $(0, \ldots, 0, t, 0, \ldots, 0)$ dans $A$. Il s'ensuit que $\dim(A) = 1$ et $\dim_k \mathfrak m/\mathfrak m^2 = n$. En particulier, $A$ est régulier si et seulement si $n = 1$.

\begin{lemma}
\label{lemma-multicross-algebra}
Soit $k$ un corps séparablement clos. Soit $A$ une algèbre locale de Nagata réduite de dimension $1$ sur $k$, de corps résiduel $k$. Alors $$
\text{invariant }\delta\text{ de }A \geq \text{nombre de branches de }A - 1
$$ En cas d'égalité, $A^\wedge$ est de la forme (\ref{equation-multicross}).
\end{lemma}

\begin{proof}
Puisque le corps résiduel de $A$ est séparablement clos, le nombre de branches de $A$ est égal au nombre de branches géométriques de $A$; voir Compléments d'algèbre, définition \ref{more-algebra-definition-number-of-branches}. L'inégalité résulte de Variétés, lemme \ref{varieties-lemma-delta-number-branches-inequality}. Supposons qu'il y ait égalité. On peut remplacer $A$ par le complété de $A$; cela ne modifie ni le nombre de branches ni l'invariant $\delta$, d'après Compléments d'algèbre, lemme \ref{more-algebra-lemma-one-dimensional-number-of-branches} et Variétés, lemme \ref{varieties-lemma-delta-same-after-completion}. Alors $A$ est strictement hensélien; voir Algèbre, lemme \ref{algebra-lemma-complete-henselian}. D'après Variétés, lemme \ref{varieties-lemma-delta-number-branches-inequality-sh}, $A$ est un bouquet d'anneaux de valuation discrète complets. Chacun d'eux est isomorphe à $k[[t]]$ d'après Algèbre, lemme \ref{algebra-lemma-regular-complete-containing-coefficient-field}. Ainsi $A$ est de la forme (\ref{equation-multicross}).
\end{proof}

\begin{definition}
\label{definition-multicross}
Soit $k$ un corps algébriquement clos. Soit $X$ un schéma algébrique de dimension $1$ sur $k$. Soit $x \in X$ un point fermé. On dit que $x$ définit une {\it singularité en croisement multiple} si le complété $\mathcal{O}_{X, x}^\wedge$ est isomorphe à (\ref{equation-multicross}) pour un certain $n \geq 2$. On dit que $x$ est un {\it nœud}, ou un {\it point double ordinaire}, ou {\it définit une singularité nodale}, si $n = 2$.
\end{definition}

\noindent
Ces singularités sont, en un sens, les singularités les plus simples que puisse présenter une courbe sur un corps algébriquement clos.

\begin{lemma}
\label{lemma-multicross}
Soit $k$ un corps algébriquement clos. Soit $X$ un schéma algébrique réduit de dimension $1$ sur $k$. Soit $x \in X$. Les conditions suivantes sont équivalentes : \begin{enumerate} \item $x$ définit une singularité en croisement multiple; \item l'invariant $\delta$ de $X$ en $x$ est égal au nombre de branches de $X$ en $x$ diminué de $1$; \item il existe une suite de morphismes $U_n \to U_{n - 1} \to \ldots \to U_0 = U \subset X$, où $U$ est un voisinage ouvert de $x$, où $U_n$ est non singulier et où chaque $U_i \to U_{i - 1}$ est une identification de deux points comme dans l'exemple \ref{example-glue-points}.
\end{enumerate}
\end{lemma}

\begin{proof}
L'équivalence de (1) et (2) est le lemme \ref{lemma-multicross-algebra}.

\medskip\noindent
Supposons (3). Montrons par récurrence descendante sur $i$ que toutes les singularités de $U_i$ sont des croisements multiples. C'est vrai pour $U_n$, puisque $U_n$ n'a pas de point singulier. Si $U_i$ s'obtient à partir de $U_{i + 1}$ en identifiant $a, b \in U_{i + 1}$ en un point $c \in U_i$, alors $$
\mathcal{O}_{U_i, c}^\wedge \subset
\mathcal{O}_{U_{i + 1}, a}^\wedge \times \mathcal{O}_{U_{i + 1}, b}^\wedge
$$ est l'ensemble des éléments ayant les mêmes classes résiduelles dans $k$. Le nombre de branches en $c$ est donc la somme des nombres de branches en $a$ et $b$, et l'invariant $\delta$ en $c$ est la somme des invariants $\delta$ en $a$ et $b$ augmentée de $1$ (car l'inclusion affichée est de codimension $1$). Cela établit (2), comme voulu.

\medskip\noindent
Supposons les conditions équivalentes (1) et (2). On peut choisir un ouvert $U \subset X$ tel que $x$ soit le seul point singulier de $U$. Appliquons alors le lemme \ref{lemma-factor-almost-isomorphism} au morphisme de normalisation $$
U^\nu = U_n \to U_{n - 1} \to \ldots \to U_1 \to U_0 = U
$$ Il reste seulement à montrer qu'aucune étape n'écrase un vecteur tangent. Supposons que $U_{i + 1} \to U_i$ soit le premier morphisme, pour le plus petit $i$, qui écrase un vecteur tangent $\theta$ en $u' \in U_{i + 1}$ au-dessus de $u \in U_i$. En raisonnant comme ci-dessus, on voit que $u_i$ est une singularité en croisement multiple (car les morphismes $U_i \to \ldots \to U_0$ identifient des paires de points). Mais le nombre de branches en $u'$ et en $u$ est le même, tandis que l'invariant $\delta$ de $U_i$ en $u$ dépasse de $1$ l'invariant $\delta$ de $U_{i + 1}$ en $u'$. D'après le lemme \ref{lemma-multicross-algebra}, $u$ ne peut donc pas être une singularité en croisement multiple, d'où une contradiction.
\end{proof}

\begin{lemma}
\label{lemma-multicross-gorenstein-is-nodal}
Soit $k$ un corps algébriquement clos. Soit $X$ un schéma algébrique réduit de dimension $1$ sur $k$. Soit $x \in X$ une singularité en croisement multiple (définition \ref{definition-multicross}). Si $X$ est de Gorenstein, alors $x$ est un nœud.
\end{lemma}

\begin{proof}
L'application $\mathcal{O}_{X, x} \to \mathcal{O}_{X, x}^\wedge$ est plate et non ramifiée au sens où $\kappa(x) = \mathcal{O}_{X, x}^\wedge/\mathfrak m_x \mathcal{O}_{X, x}^\wedge$. (Voir Compléments d'algèbre, section \ref{more-algebra-section-permanence-completion}.) Ainsi, si $X$ est de Gorenstein, $\mathcal{O}_{X, x}$ est de Gorenstein, puis $\mathcal{O}_{X, x}^\wedge$ est de Gorenstein d'après Complexes dualisants, lemme \ref{dualizing-lemma-flat-under-gorenstein}. Il suffit donc de montrer que l'anneau $A$ de (\ref{equation-multicross}), avec $n \geq 2$, est de Gorenstein si et seulement si $n = 2$.

\medskip\noindent
Si $n = 2$, alors $A = k[[x, y]]/(xy)$ est une intersection complète, donc de Gorenstein. Cela résulte par exemple de Dualité pour les schémas, lemme \ref{duality-lemma-gorenstein-lci} appliqué à $k[[x, y]] \to A$, et du fait que l'anneau local régulier $k[[x, y]]$ est de Gorenstein d'après Complexes dualisants, lemme \ref{dualizing-lemma-regular-gorenstein}.

\medskip\noindent
Supposons $n > 2$. Si $A$ était de Gorenstein, $A$ serait un complexe dualisant sur $A$ (Dualité pour les schémas, définition \ref{duality-definition-gorenstein}). Alors $R\Hom(k, A)$ serait égal à $k[n]$ pour un certain $n \in \mathbf{Z}$; voir Complexes dualisants, lemme \ref{dualizing-lemma-find-function}. Il en résulterait que $\Ext^1_A(k, A) \cong k$ ou $\Ext^1_A(k, A) = 0$ (suivant la valeur de $n$; en fait, $n$ doit valoir $-1$, mais cela n'importe pas ici). La suite exacte $$
0 \to \mathfrak m_A \to A \to k \to 0
$$ donne $$
\Ext^1_A(k, A) = \Hom_A(\mathfrak m_A, A)/A
$$, où $A \to \Hom_A(\mathfrak m_A, A)$ est défini par $a \mapsto (a' \mapsto aa')$. Soit $e_i \in \Hom_A(\mathfrak m_A, A)$ l'élément qui envoie $(f_1, \ldots, f_n) \in \mathfrak m_A$ sur $(0, \ldots, 0, f_i, 0, \ldots, 0)$. Le lecteur vérifie aisément que $e_1, \ldots, e_{n - 1}$ sont linéairement indépendants sur $k$ dans $\Hom_A(\mathfrak m_A, A)/A$. Ainsi $\dim_k \Ext^1_A(k, A) \geq n - 1 \geq 2$, ce qui achève la démonstration. (Remarquons que $e_1 + \ldots + e_n$ est l'image de $1$ par l'application $A \to \Hom_A(\mathfrak m_A, A)$.)
\end{proof}




\section{Torsion du groupe de Picard}
\label{section-torsion-in-pic}

\noindent
Dans cette section, nous majorons la torsion du groupe de Picard d'un schéma propre de dimension $1$ sur un corps. Nous utiliserons ce résultat dans l'étude de la réduction semi-stable des courbes.

\medskip\noindent
Il ne semble pas exister de démonstration élémentaire du résultat du lemme \ref{lemma-torsion-picard-smooth-projective}. L'analyse de sa preuve fait apparaître deux ingrédients essentiels : (1) l'existence d'une variété abélienne classifiant les faisceaux inversibles de degré zéro sur une courbe projective lisse et (2) la détermination de la structure des points de torsion d'une variété abélienne.

\begin{lemma}
\label{lemma-torsion-picard-smooth-projective}
Soit $k$ un corps algébriquement clos. Soit $X$ une courbe projective lisse de genre $g$ sur $k$. \begin{enumerate} \item Si $n \geq 1$ est inversible dans $k$, alors $\Pic(X)[n] \cong (\mathbf{Z}/n\mathbf{Z})^{\oplus 2g}$. \item Si la caractéristique de $k$ est $p > 0$, il existe un entier $0 \leq f \leq g$ tel que $\Pic(X)[p^m] \cong (\mathbf{Z}/p^m\mathbf{Z})^{\oplus f}$ pour tout $m \geq 1$.
\end{enumerate}
\end{lemma}

\begin{proof}
Notons $\Pic^0(X) \subset \Pic(X)$ le sous-groupe des faisceaux inversibles de degré $0$. Autrement dit, on a une suite exacte courte $$
0 \to \Pic^0(X) \to \Pic(X) \xrightarrow{\deg} \mathbf{Z} \to 0.
$$ Le groupe $\Pic^0(X)$ est le groupe des points à valeurs dans $k$ du schéma en groupes $\underline{\Picardfunctor}^0_{X/k}$; voir Schémas de Picard des courbes, lemme \ref{pic-lemma-picard-pieces}. Le même lemme affirme que $\underline{\Picardfunctor}^0_{X/k}$ est une variété abélienne de dimension $g$ sur $k$ au sens de Groupoïdes, définition \ref{groupoids-definition-abelian-variety}. On conclut donc par Groupoïdes, proposition \ref{groupoids-proposition-review-abelian-varieties}.
\end{proof}

\begin{lemma}
\label{lemma-torsion-picard-becomes-visible}
Soit $k$ un corps. Soit $n$ premier à la caractéristique de $k$. Soit $X$ une courbe propre lisse sur $k$ telle que $H^0(X, \mathcal{O}_X) = k$ et de genre $g$. \begin{enumerate} \item Si $g = 1$, il existe une extension finie séparable $k'/k$ telle que $X_{k'}$ possède un point $k'$-rationnel et que $\Pic(X_{k'})[n] \cong (\mathbf{Z}/n\mathbf{Z})^{\oplus 2}$. \item Si $g \geq 2$, il existe une extension finie séparable $k'/k$ avec $[k' : k] \leq (2g - 2)(n^{2g})!$ telle que $X_{k'}$ possède un point $k'$-rationnel et que $\Pic(X_{k'})[n] \cong (\mathbf{Z}/n\mathbf{Z})^{\oplus 2g}$.
\end{enumerate}
\end{lemma}

\begin{proof}
Supposons $g \geq 2$. On peut d'abord choisir une extension finie séparable de degré au plus $2g - 2$ sur laquelle $X$ acquiert un point rationnel; voir le lemme \ref{lemma-point-over-separable-extension}. On peut donc supposer que $X$ possède un point $k$-rationnel, noté $x \in X(k)$, mais il faut désormais démontrer le lemme avec $[k' : k] \leq (n^{2g})!$. Soit $k \subset k^{sep} \subset \overline{k}$ une clôture séparable contenue dans une clôture algébrique. D'après le lemme \ref{lemma-torsion-picard-smooth-projective}, on a $$
\Pic(X_{\overline{k}})[n] \cong (\mathbf{Z}/n\mathbf{Z})^{\oplus 2g}
$$ D'après Schémas de Picard des courbes, lemme \ref{pic-lemma-torsion-descends}, on en déduit $$
\Pic(X_{k^{sep}})[n] \cong (\mathbf{Z}/n\mathbf{Z})^{\oplus 2g}
$$ D'après Schémas de Picard des courbes, lemme \ref{pic-lemma-torsion-descends}, il existe une action continue $$
\text{Gal}(k^{sep}/k)
\longrightarrow
\text{Aut}(\Pic(X_{k^{sep}})[n]
$$, et le lemme est vrai pour le corps fixe $k'$ du noyau de cette application. Ce noyau est ouvert par continuité de l'action, ce qui entraîne que $k'/k$ est finie. Par la théorie de Galois, $\text{Gal}(k'/k)$ est l'image de la flèche affichée. Puisque le groupe des permutations d'un ensemble de cardinal $n^{2g}$ a pour cardinal $(n^{2g})!$, la théorie de Galois donne $[k' : k] \leq (n^{2g})!$. (Bien entendu, cela démontre le lemme avec la borne $|\text{GL}_{2g}(\mathbf{Z}/n\mathbf{Z})|$, mais seule l'existence d'une borne nous importe ici.)

\medskip\noindent
Si le genre est $1$, il n'y a pas de majoration du degré d'une extension finie séparable sur laquelle $X$ acquiert un point rationnel (nous omettons les détails). Une telle extension existe néanmoins, par exemple d'après Variétés, lemme \ref{varieties-lemma-smooth-separable-closed-points-dense}. La suite de la démonstration est identique au cas $g \geq 2$.
\end{proof}

\begin{proposition}
\label{proposition-torsion-picard-reduced-proper}
Soit $k$ un corps algébriquement clos. Soit $X$ un schéma propre sur $k$, réduit, connexe et de dimension $1$. Soit $g$ le genre de $X$ et soit $g_{geom}$ la somme des genres géométriques des composantes irréductibles de $X$. Pour tout nombre premier $\ell$ différent de la caractéristique de $k$, on a $$
\dim_{\mathbf{F}_\ell} \Pic(X)[\ell]
\leq g + g_{geom}
$$, avec égalité si et seulement si toutes les singularités de $X$ sont des croisements multiples.
\end{proposition}

\begingroup\emergencystretch=4em
\begin{proof}
Soit $\nu : X^\nu \to X$ la normalisation (Variétés, lemme \ref{varieties-lemma-prepare-delta-invariant}). Choisissons une factorisation $$
X^\nu = X_n \to X_{n - 1} \to \ldots \to X_1 \to X_0 = X
$$ comme dans le lemme \ref{lemma-factor-almost-isomorphism}. Notons $h^0_i = \dim_k H^0(X_i, \mathcal{O}_{X_i})$ et $h^1_i = \dim_k H^1(X_i, \mathcal{O}_{X_i})$. D'après les lemmes \ref{lemma-glue-points} et \ref{lemma-squish-tangent-vector}, pour chaque $n > i \geq 0$, on est dans l'un des trois cas suivants : \begin{enumerate} \item $X_i$ s'obtient en identifiant les points $a, b \in X_{i + 1}$, situés sur des composantes connexes distinctes : dans ce cas, $\Pic(X_i) = \Pic(X_{i + 1})$, $h^0_{i + 1} = h^0_i + 1$, $h^1_{i + 1} = h^1_i$; \item $X_i$ s'obtient en identifiant les points $a, b \in X_{i + 1}$, situés sur une même composante connexe : dans ce cas, on a une suite exacte courte $$
0 \to k^* \to \Pic(X_i) \to \Pic(X_{i + 1}) \to 0,
$$ et $h^0_{i + 1} = h^0_i$, $h^1_{i + 1} = h^1_i - 1$; \item $X_i$ s'obtient en écrasant un vecteur tangent dans $X_{i + 1}$ : dans ce cas, on a une suite exacte courte $$
0 \to (k, +) \to \Pic(X_i) \to \Pic(X_{i + 1}) \to 0,
$$ et $h^0_{i + 1} = h^0_i$, $h^1_{i + 1} = h^1_i - 1$. \end{enumerate} Pour les assertions sur les dimensions des groupes de cohomologie du faisceau structural, on utilise les suites exactes des exemples \ref{example-glue-points} et \ref{example-squish-tangent-vector}. Puisque $k$ est algébriquement clos de caractéristique première à $\ell$, les groupes $(k, +)$ et $k^*$ sont divisibles par $\ell$ et leur torsion d'ordre $\ell$ est respectivement $(k, +)[\ell] = 0$ et $k^*[\ell] \cong \mathbf{F}_\ell$. Ainsi $$
\dim_{\mathbf{F}_\ell} \Pic(X_{i + 1})[\ell] -
\dim_{\mathbf{F}_\ell}\Pic(X_i)[\ell]
$$ est nul, sauf dans le cas (2), où il vaut $-1$. À l'issue de ce processus, on obtient la normalisation $X^\nu = X_n$, qui est une réunion disjointe de courbes projectives lisses sur $k$. On a donc \begin{enumerate} \item $h^1_n = g_{geom}$ et \item $\dim_{\mathbf{F}_\ell} \Pic(X_n)[\ell] = 2g_{geom}$. \end{enumerate} La dernière égalité résulte du lemme \ref{lemma-torsion-picard-smooth-projective}. Puisque $g = h^1_0$, le nombre d'étapes de types (2) et (3) est au plus $h^1_0 - h^1_n = g - g_{geom}$. Le calcul des différences de rangs donne $$
\dim_{\mathbf{F}_\ell} \Pic(X)[\ell] \leq
g - g_{geom} + 2g_{geom} = g + g_{geom}
$$, avec égalité si et seulement si aucune étape de type (3) n'intervient. Cela signifie bien que toutes les singularités de $X$ sont des croisements multiples, d'après le lemme \ref{lemma-multicross}. Réciproquement, si toutes les singularités sont des croisements multiples, le lemme \ref{lemma-multicross} assure l'existence d'une suite $X^\nu = X_n \to \ldots \to X_0 = X$ comme ci-dessus sans étape de type (3), et l'on obtient l'égalité dans le lemme (il suffit de recoller les suites locales en chaque point pour en obtenir une qui convienne à tous les points singuliers de $x$; nous omettons quelques détails).
\end{proof}
\endgroup





\section{Comparaison du genre et du genre géométrique}
\label{section-genus-geometric-genus}

\noindent
Soit $k$ un corps, et notons $\overline{k}$ une clôture algébrique de ce corps. Soit $X$ un schéma propre de dimension $\leq 1$ sur $k$. Nous définissons $g_{geom}(X/k)$ comme la somme des genres géométriques des composantes irréductibles de $X_{\overline{k}}$ de dimension $1$.

\begin{lemma}
\label{lemma-bound-geometric-genus}
Soit $k$ un corps. Soit $X$ un schéma propre de dimension $\leq 1$ sur $k$. Alors $$
g_{geom}(X/k) = \sum\nolimits_{C \subset X} g_{geom}(C/k)
$$, la somme portant sur les composantes irréductibles $C \subset X$ de dimension $1$.
\end{lemma}

\begin{proof}
Cela résulte immédiatement de la définition et du fait qu'une composante irréductible $\overline{Z}$ de $X_{\overline{k}}$ se projette surjectivement sur une composante irréductible $Z$ de $X$ (Variétés, lemme \ref{varieties-lemma-image-irreducible}) de même dimension (Morphismes, lemme \ref{morphisms-lemma-dimension-fibre-after-base-change} appliqué au point générique de $\overline{Z}$).
\end{proof}

\begingroup\emergencystretch=4em
\begin{lemma}
\label{lemma-geometric-genus-normalization}
Soit $k$ un corps. Soit $X$ un schéma propre de dimension $\leq 1$ sur $k$. Alors \begin{enumerate} \item On a $g_{geom}(X/k) = g_{geom}(X_{red}/k)$. \item Si $X' \to X$ est un morphisme propre birationnel, alors $g_{geom}(X'/k) = g_{geom}(X/k)$. \item Si $X^\nu \to X$ est le morphisme de normalisation, alors $g_{geom}(X^\nu/k) = g_{geom}(X/k)$.
\end{enumerate}
\end{lemma}
\endgroup

\begin{proof}
L'assertion (1) résulte immédiatement du lemme \ref{lemma-bound-geometric-genus}. Si $X' \to X$ est propre birationnel, il est fini et c'est un isomorphisme au-dessus d'un ouvert dense (voir Variétés, lemmes \ref{varieties-lemma-finite-in-codim-1} et \ref{varieties-lemma-modification-normal-iso-over-codimension-1}). Ainsi $X'_{\overline{k}} \to X_{\overline{k}}$ est un isomorphisme au-dessus d'un ouvert dense. Les composantes irréductibles de $X'_{\overline{k}}$ et $X_{\overline{k}}$ sont donc en bijection, et les composantes correspondantes ont des corps de fonctions isomorphes. Leurs modèles projectifs non singuliers sont donc isomorphes, et leurs genres géométriques égaux. Cela prouve (2). L'assertion (3) résulte de (1), de (2) et du fait que $X^\nu \to X_{red}$ est birationnel (Morphismes, lemme \ref{morphisms-lemma-normalization-birational}).
\end{proof}

\begin{lemma}
\label{lemma-genus-goes-down}
Soit $k$ un corps. Soit $X$ un schéma propre de dimension $\leq 1$ sur $k$. Soit $f : Y \to X$ un morphisme fini tel qu'il existe un ouvert dense $U \subset X$ au-dessus duquel $f$ soit une immersion fermée. Alors
$$
\dim_k H^1(X, \mathcal{O}_X) \geq \dim_k H^1(Y, \mathcal{O}_Y)
$$
\end{lemma}

\begin{proof}
Considérons la suite exacte $$
0 \to \mathcal{G} \to \mathcal{O}_X \to f_*\mathcal{O}_Y \to \mathcal{F} \to 0
$$ de faisceaux cohérents sur $X$. Par hypothèse, $\mathcal{F}$ est à support dans un nombre fini de points fermés, et sa cohomologie supérieure est donc nulle (Variétés, lemme \ref{varieties-lemma-chi-tensor-finite}). D'autre part, on a $H^2(X, \mathcal{G}) = 0$ d'après Cohomologie, proposition \ref{cohomology-proposition-vanishing-Noetherian}. Il s'ensuit formellement que l'application induite $H^1(X, \mathcal{O}_X) \to H^1(X, f_*\mathcal{O}_Y)$ est surjective. Puisque $H^1(X, f_*\mathcal{O}_Y) = H^1(Y, \mathcal{O}_Y)$ (Cohomologie des schémas, lemme \ref{coherent-lemma-relative-affine-cohomology}), le lemme en résulte.
\end{proof}

\begin{lemma}
\label{lemma-genus-normalization}
Soit $k$ un corps. Soit $X$ un schéma propre de dimension $\leq 1$ sur $k$. Si $X' \to X$ est un morphisme propre birationnel, alors $$
\dim_k H^1(X, \mathcal{O}_X) \geq \dim_k H^1(X', \mathcal{O}_{X'})
$$ Si $X$ est réduit, si $H^0(X, \mathcal{O}_X) \to H^0(X', \mathcal{O}_{X'})$ est surjective et s'il y a égalité, alors $X' = X$.
\end{lemma}

\begin{proof}
Si $f : X' \to X$ est propre birationnel, il est fini et c'est un isomorphisme au-dessus d'un ouvert dense (voir Variétés, lemmes \ref{varieties-lemma-finite-in-codim-1} et \ref{varieties-lemma-modification-normal-iso-over-codimension-1}). L'inégalité résulte donc du lemme \ref{lemma-genus-goes-down}. Supposons $X$ réduit. Alors $\mathcal{O}_X \to f_*\mathcal{O}_{X'}$ est injective et l'on obtient une suite exacte courte $$
0 \to \mathcal{O}_X \to f_*\mathcal{O}_{X'} \to \mathcal{F} \to 0
$$ Sous les hypothèses de la seconde assertion, la suite exacte longue de cohomologie donne $H^0(X, \mathcal{F}) = 0$. Alors $\mathcal{F} = 0$, car $\mathcal{F}$ est engendré par ses sections globales (Variétés, lemme \ref{varieties-lemma-chi-tensor-finite}), et $\mathcal{O}_X = f_*\mathcal{O}_{X'}$. Puisque $f$ est affine, cela entraîne $X = X'$.
\end{proof}

\begin{lemma}
\label{lemma-bound-geometric-genus-curve}
Soit $k$ un corps. Soit $C$ une courbe propre sur $k$. Posons $\kappa = H^0(C, \mathcal{O}_C)$. Alors
$$
[\kappa : k]_s \dim_\kappa H^1(C, \mathcal{O}_C) \geq g_{geom}(C/k)
$$
\end{lemma}

\begin{proof}
D'après Variétés, lemme \ref{varieties-lemma-regular-functions-proper-variety}, $\kappa$ est un corps et une extension finie de $k$. Par Corps, lemme \ref{fields-lemma-separable-degree}, on a $[\kappa : k]_s = |\Mor_k(\kappa, \overline{k})|$; ainsi $\Spec(\kappa \otimes_k \overline{k})$ possède $[\kappa : k]_s$ points ayant chacun pour corps résiduel $\overline{k}$. On a donc $$
C_{\overline{k}} =
\bigcup\nolimits_{t \in \Spec(\kappa \otimes_k \overline{k})} C_t
$$ (réunion ensembliste). Ici $C_t = C \times_{\Spec(\kappa), t} \Spec(\overline{k})$, où l'on considère $t$ comme un homomorphisme de $k$-algèbres $t : \kappa \to \overline{k}$. Il en résulte que $g_{geom}(C/k) = \sum_t g_{geom}(C_t/\overline{k})$, la somme étant indexée par un ensemble de cardinal $[\kappa : k]_s$. On a $$
H^0(C_t, \mathcal{O}_{C_t}) = \overline{k}
\quad\text{et}\quad
\dim_{\overline{k}} H^1(C_t, \mathcal{O}_{C_t}) =
\dim_\kappa H^1(C, \mathcal{O}_C)
$$ par cohomologie et changement de base (Cohomologie des schémas, lemme \ref{coherent-lemma-flat-base-change-cohomology}). Remarquons que la normalisation $C_t^\nu$ est la réunion disjointe des modèles projectifs non singuliers des composantes irréductibles de $C_t$ (Morphismes, lemme \ref{morphisms-lemma-normalization-in-terms-of-components}). Ainsi $\dim_{\overline{k}} H^1(C_t^\nu, \mathcal{O}_{C_t^\nu})$ est égal à $g_{geom}(C_t/\overline{k})$. D'après le lemme \ref{lemma-genus-goes-down}, on a $$
\dim_{\overline{k}} H^1(C_t, \mathcal{O}_{C_t}) \geq
\dim_{\overline{k}} H^1(C_t^\nu, \mathcal{O}_{C_t^\nu})
$$, ce qui achève la démonstration.
\end{proof}

\begin{lemma}
\label{lemma-bound-torsion-simple}
Soit $k$ un corps. Soit $X$ un schéma propre de dimension $\leq 1$ sur $k$. Soit $\ell$ un nombre premier inversible dans $k$. Alors $$
\dim_{\mathbf{F}_\ell} \Pic(X)[\ell] \leq
\dim_k H^1(X, \mathcal{O}_X) + g_{geom}(X/k)
$$, où $g_{geom}(X/k)$ est défini ci-dessus.
\end{lemma}

\begingroup\emergencystretch=4em
\begin{proof}
L'application $\Pic(X) \to \Pic(X_{\overline{k}})$ est injective d'après Variétés, lemme \ref{varieties-lemma-change-fields-pic}. D'après Cohomologie des schémas, lemme \ref{coherent-lemma-flat-base-change-cohomology}, $\dim_k H^1(X, \mathcal{O}_X)$ est égal à $\dim_{\overline{k}} H^1(X_{\overline{k}}, \mathcal{O}_{X_{\overline{k}}})$. On peut donc supposer $k$ algébriquement clos.

\medskip\noindent
Soit $X_{red}$ le réduit de $X$. La surjection $\mathcal{O}_X \to \mathcal{O}_{X_{red}}$ induit une surjection $H^1(X, \mathcal{O}_X) \to H^1(X, \mathcal{O}_{X_{red}})$, car la cohomologie des faisceaux quasi-cohérents s'annule en degrés $\geq 2$ d'après Cohomologie, proposition \ref{cohomology-proposition-vanishing-Noetherian}. Puisque $X_{red} \to X$ induit une bijection sur les composantes irréductibles au-dessus de $\overline{k}$ et un isomorphisme sur la torsion d'ordre $\ell$ des groupes de Picard (Schémas de Picard des courbes, lemme \ref{pic-lemma-torsion-descends}), on peut remplacer $X$ par $X_{red}$. On est ainsi ramené à la proposition \ref{proposition-torsion-picard-reduced-proper}.
\end{proof}
\endgroup






\section{Courbes nodales}
\label{section-nodal}

\noindent
Nous avons déjà défini les points doubles ordinaires sur un corps algébriquement clos; voir la définition \ref{definition-multicross}. Si $x \in X$ est un point fermé d'un schéma algébrique de dimension $1$ sur un corps algébriquement clos $k$, alors $x$ est un point double ordinaire si et seulement si $$
\mathcal{O}_{X, x}^\wedge \cong k[[x, y]]/(xy)
$$ Voir la discussion qui suit (\ref{equation-multicross}) dans la section \ref{section-multicross}.

\begin{definition}
\label{definition-nodal}
Soit $k$ un corps. Soit $X$ un schéma localement algébrique de dimension $1$ sur $k$. \begin{enumerate} \item Un point fermé $x \in X$ est appelé un {\it nœud}, ou un {\it point double ordinaire}, ou est dit {\it définir une singularité nodale}, s'il existe un point double ordinaire $\overline{x} \in X_{\overline{k}}$ d'image $x$. \item On dit que {\it les singularités de $X$ sont au plus nodales} si tout point fermé de $X$ appartient au lieu lisse du morphisme structural $X \to \Spec(k)$ ou est un point double ordinaire.
\end{enumerate}
\end{definition}

\noindent
On appelle souvent {\it courbe nodale} un schéma algébrique de dimension $1$, noté $X$, dont les singularités $X$ sont au plus nodales. On exige parfois aussi qu'une courbe nodale soit propre. Puisqu'une courbe nodale ainsi définie n'est pas nécessairement irréductible, cette terminologie entre en conflit avec notre définition antérieure d'une courbe comme variété de dimension $1$.

\begin{lemma}
\label{lemma-reduced-quotient-regular-ring-dim-2}
Soit $(A, \mathfrak m)$ un anneau local régulier de dimension $2$. Soit $I \subset \mathfrak m$ un idéal. \begin{enumerate} \item Si $A/I$ est réduit, alors $I = (0)$, $I = \mathfrak m$ ou $I = (f)$ pour un élément non nul $f \in \mathfrak m$. \item Si $A/I$ est de profondeur $1$, alors $I = (f)$ pour un élément non nul $f \in \mathfrak m$. 
\end{enumerate}
\end{lemma}

\begin{proof}
Supposons $I \not = 0$. Écrivons $I = (f_1, \ldots, f_r)$. Puisque $A$ est factoriel (Compléments d'algèbre, lemme \ref{more-algebra-lemma-regular-local-UFD}), on peut écrire $f_i = fg_i$, où $f$ est le pgcd des $f_1, \ldots, f_r$. Le pgcd des $g_1, \ldots, g_r$ est donc $1$, ce qui signifie qu'il n'y a pas d'idéal premier de hauteur $1$ contenant $g_1, \ldots, g_r$. Alors soit $(g_1, \ldots, g_r) = A$, ce qui entraîne $I = (f)$, soit, dans le cas contraire, $\dim(A) = 2$ entraîne $V(g_1, \ldots, g_r) = \{\mathfrak m\}$, c'est-à-dire $\mathfrak m = \sqrt{(g_1, \ldots, g_r)}$.

\medskip\noindent
Supposons $A/I$ réduit, c'est-à-dire $I$ radical. Si $f$ est une unité, puisque $I$ est radical, on a $I = \mathfrak m$. Si $f \in \mathfrak m$, on voit que $f^n$ s'envoie sur zéro dans $A/I$. Ainsi $f \in I$ par réduction, et l'on en conclut que $I = (f)$.

\medskip\noindent
Supposons $A/I$ de profondeur $1$. Alors $\mathfrak m$ n'est pas un idéal premier associé à $A/I$. Puisque la classe de $f$ modulo $I$ est annulée par $g_1, \ldots, g_r$, la classe de $f$ est nulle dans $A/I$. Ainsi $I = (f)$, comme voulu.
\end{proof}

\noindent
Soit $\kappa$ un corps et soit $V$ un espace vectoriel sur $\kappa$. On dira que $q \in \text{Sym}^2_\kappa(V)$ est {\it non dégénérée} si l'application linéaire sur $\kappa$ induite $V^\vee \to V$ est un isomorphisme. Si $q = \sum_{i \leq j} a_{ij} x_i x_j$ pour une base sur $\kappa$, notée $x_1, \ldots, x_n$, de $V$, cela signifie que le déterminant de la matrice $$
\left(
\begin{matrix}
2a_{11} & a_{12} & \ldots \\
a_{12} & 2a_{22} & \ldots \\
\ldots & \ldots & \ldots
\end{matrix}
\right)
$$ est non nul. Cela équivaut à la condition que les dérivées partielles de $q$ par rapport aux $x_i$ définissent $0$ au sens schématique.

\begin{lemma}
\label{lemma-nodal-algebraic}
Soit $k$ un corps. Soit $(A, \mathfrak m, \kappa)$ une $k$-algèbre locale noethérienne. Les conditions suivantes sont équivalentes : \begin{enumerate} \item $\kappa/k$ est séparable, $A$ est réduit, $\dim_\kappa(\mathfrak m/\mathfrak m^2) = 2$, et il existe un élément non dégénéré $q \in \text{Sym}^2_\kappa(\mathfrak m/\mathfrak m^2)$ qui s'envoie sur zéro dans $\mathfrak m^2/\mathfrak m^3$; \item $\kappa/k$ est séparable, $\text{depth}(A) = 1$, $\dim_\kappa(\mathfrak m/\mathfrak m^2) = 2$, et il existe un élément non dégénéré $q \in \text{Sym}^2_\kappa(\mathfrak m/\mathfrak m^2)$ qui s'envoie sur zéro dans $\mathfrak m^2/\mathfrak m^3$; \item $\kappa/k$ est séparable, $A^\wedge \cong \kappa[[x, y]]/(ax^2 + bxy + cy^2)$ en tant que $k$-algèbre, où $ax^2 + bxy + cy^2$ est une forme quadratique non dégénérée sur $\kappa$.
\end{enumerate}
\end{lemma}

\begin{proof}
Supposons (3). Alors $A^\wedge$ est réduit, car $ax^2 + bxy + cy^2$ est soit irréductible, soit produit de deux éléments premiers non associés. Ainsi $A \subset A^\wedge$ est réduit. Il s'ensuit que (1) est vraie.

\medskip\noindent
Supposons (1). Alors $A$ ne peut pas être artinien, car il ne serait pas réduit puisque $\mathfrak m \not = (0)$. Ainsi $\dim(A) \geq 1$, puis $\text{depth}(A) \geq 1$ d'après Algèbre, lemme \ref{algebra-lemma-criterion-reduced}. D'autre part, $\dim(A) = 2$ entraîne que $A$ est régulier, ce qui contredit l'existence de $q$ d'après Algèbre, lemme \ref{algebra-lemma-regular-graded}. Ainsi $\dim(A) \leq 1$, puis $\text{depth}(A) = 1$ d'après Algèbre, lemme \ref{algebra-lemma-bound-depth}. Il s'ensuit que (2) est vraie.

\medskip\noindent
Supposons (2). Puisque la profondeur de $A$ est égale à celle de $A^\wedge$ (Compléments d'algèbre, lemme \ref{more-algebra-lemma-completion-depth}) et que les autres conditions sont invariantes par complétion, on peut supposer $A$ complet. Choisissons $\kappa \to A$ comme dans Compléments d'algèbre, lemme \ref{more-algebra-lemma-lift-residue-field}. Puisque $\dim_\kappa(\mathfrak m/\mathfrak m^2) = 2$, on peut choisir $x_0, y_0 \in \mathfrak m$ dont les images forment une base. On obtient un homomorphisme continu de $\kappa$-algèbres $$
\kappa[[x, y]] \longrightarrow A
$$ par les règles $x \mapsto x_0$ et $y \mapsto y_0$. Soit $q$ la classe de $ax_0^2 + bx_0y_0 + cy_0^2$ dans $\text{Sym}^2_\kappa(\mathfrak m/\mathfrak m^2)$. Écrivons $Q(x, y) = ax^2 + bxy + cy^2$, considéré comme un polynôme en deux variables. On voit alors que $$
Q(x_0, y_0) = ax_0^2 + bx_0y_0 + cy_0^2 =
\sum\nolimits_{i + j = 3} a_{ij} x_0^iy_0^j
$$ pour certains $a_{ij}$ dans $A$. Nous voulons augmenter l'ordre d'annulation en modifiant le choix de $x_0$, $y_0$. Supposons que $x_1, y_1 \in \mathfrak m^2$. Alors $$
Q(x_0 + x_1, y_0 + y_1) = Q(x_0, y_0) +
(2ax_0 + by_0)x_1 + (bx_0 + 2cy_0)y_1 \bmod \mathfrak m^4
$$ La non-dégénérescence de $Q$ signifie précisément que $2ax_0 + by_0$ et $bx_0 + 2cy_0$ forment une base sur $\kappa$ de $\mathfrak m/\mathfrak m^2$; voir la discussion précédant le lemme. On peut donc choisir $x_1, y_1 \in \mathfrak m^2$ tels que $Q(x_0 + x_1, y_0 + y_1) \in \mathfrak m^4$. En poursuivant par récurrence, on trouve $x_i, y_i \in \mathfrak m^{i + 1}$ tels que $$
Q(x_0 + x_1 + \ldots + x_n, y_0 + y_1 + \ldots + y_n) \in \mathfrak m^{n + 3}
$$ Puisque $A$ est complet, on peut poser $x_\infty = \sum x_i$ et $y_\infty = \sum y_i$ et considérer l'application $\kappa[[x, y]] \longrightarrow A$ envoyant $x$ sur $x_\infty$ et $y$ sur $y_\infty$. Elle induit une surjection $\kappa[[x, y]]/(Q) \longrightarrow A$ d'après Algèbre, lemme \ref{algebra-lemma-completion-generalities}. D'après le lemme \ref{lemma-reduced-quotient-regular-ring-dim-2}, le noyau de $k[[x, y]] \to A$ est principal. Mais ce noyau ne peut contenir de diviseur propre de $Q$, car un tel diviseur serait de degré $1$ dans $x, y$, ce qui contredirait $\dim(\mathfrak m/\mathfrak m^2) = 2$. Ainsi $Q$ engendre le noyau, comme voulu.
\end{proof}

\begin{lemma}
\label{lemma-2-branches-delta-1}
Soit $k$ un corps. Soit $(A, \mathfrak m, \kappa)$ une $k$-algèbre locale de Nagata. Les conditions suivantes sont équivalentes : \begin{enumerate} \item $k \to A$ est comme dans le lemme \ref{lemma-nodal-algebraic}; \item $\kappa/k$ est séparable, $A$ est réduit de dimension $1$, l'invariant $\delta$ de $A$ vaut $1$, et $A$ possède $2$ branches géométriques. \end{enumerate} Si ces conditions sont satisfaites, la clôture intégrale $A'$ de $A$ dans son anneau total des fractions possède soit $1$, soit $2$ idéaux maximaux $\mathfrak m'$, et les extensions $\kappa(\mathfrak m')/k$ sont séparables.
\end{lemma}

\begin{proof}
Dans les deux cas, $A$ et $A^\wedge$ sont réduits. Dans le cas (2), cela tient au fait que le complété d'un anneau local de Nagata réduit est réduit (Compléments d'algèbre, lemme \ref{more-algebra-lemma-completion-reduced}). Dans les deux cas, $A$ et $A^\wedge$ sont de dimension $1$ (Compléments d'algèbre, lemme \ref{more-algebra-lemma-completion-dimension}). L'invariant $\delta$ et le nombre de branches géométriques de $A$ et $A^\wedge$ coïncident d'après Variétés, lemme \ref{varieties-lemma-delta-same-after-completion} et Compléments d'algèbre, lemme \ref{more-algebra-lemma-one-dimensional-number-of-branches}. Soit $A'$ la clôture intégrale de $A$ dans son anneau total des fractions, comme dans Variétés, lemme \ref{varieties-lemma-pre-delta-invariant}. D'après Variétés, lemme \ref{varieties-lemma-normalization-same-after-completion}, $A' \otimes_A A^\wedge$ joue le même rôle pour $A^\wedge$. On peut donc remplacer $A$ par $A^\wedge$ et supposer $A$ complet.

\medskip\noindent
Supposons (1). Il suffit de montrer que $A$ possède deux branches géométriques et que son invariant $\delta$ vaut $1$. On peut supposer $A = \kappa[[x, y]]/(ax^2 + bxy + cy^2)$ avec $q = ax^2 + bxy + cy^2$ non dégénérée. Il y a deux cas.

\medskip\noindent
Cas I : $q$ se décompose sur $\kappa$. Après changement de coordonnées, on peut alors supposer $q = xy$. On voit que
$$
A' = \kappa[[x, y]]/(x) \times \kappa[[x, y]]/(y)
$$

\medskip\noindent
Cas II : $q$ ne se décompose pas. Dans ce cas, $c \not = 0$ et la non-dégénérescence signifie que $b^2 - 4ac \not = 0$. Ainsi $\kappa' = \kappa[t]/(a + bt + ct^2)$ est une extension séparable de degré $2$ de $\kappa$. Alors $t = y/x$ est entier sur $A$, et l'on en conclut que $$
A' = \kappa'[[x]]
$$, où $y$ s'envoie sur $tx$ dans le membre de droite.

\medskip\noindent
Dans les deux cas, on vérifie directement que l'invariant $\delta$ vaut $1$ et que le nombre de branches géométriques est $2$. Cela montre que (1) entraîne (2). La dernière assertion du lemme est également établie.

\medskip\noindent
Supposons (2). Compléments d'algèbre, lemme \ref{more-algebra-lemma-number-of-branches-1} montre que $A'$ possède soit deux idéaux maximaux, soit un seul idéal maximal pour $A'$ avec $[\kappa(\mathfrak m') : \kappa]_s = 2$.

\medskip\noindent
Cas I : $A'$ possède deux idéaux maximaux $\mathfrak m'_1$, $\mathfrak m'_2$, de corps résiduels $\kappa_1$, $\kappa_2$. Puisque l'invariant $\delta$ est la longueur de $A'/A$ et qu'il existe une surjection $A'/A \to (\kappa_1 \times \kappa_2)/\kappa$, on a $\kappa = \kappa_1 = \kappa_2$. Puisque $A$ est complet (et hensélien d'après Algèbre, lemme \ref{algebra-lemma-complete-henselian}) et que $A'$ est fini sur $A$, on a $A' = A_1 \times A_2$ (par Algèbre, lemme \ref{algebra-lemma-finite-over-henselian}). Puisque $A'$ est un anneau normal, $A_1$ et $A_2$ sont des anneaux de valuation discrète. Ainsi $A_1$ et $A_2$ sont isomorphes à $\kappa[[t]]$ (comme $k$-algèbres) d'après Compléments d'algèbre, lemme \ref{more-algebra-lemma-power-series-over-residue-field}. Puisque l'invariant $\delta$ vaut $1$, $A$ est le bouquet de $A_1$ et $A_2$ (Variétés, définition \ref{varieties-definition-wedge}). Il s'ensuit aisément que $A \cong \kappa[[x, y]]/(xy)$.

\medskip\noindent
Cas II : $A'$ possède un seul idéal maximal $\mathfrak m'$, de corps résiduel $\kappa'$, et $[\kappa' : \kappa]_s = 2$. En raisonnant exactement comme dans le cas I, on voit que $[\kappa' : \kappa] = 2$ et que $\kappa'$ est séparable sur $\kappa$. Puisque $A'$ est normal, $A'$ est isomorphe à $\kappa'[[t]]$ (voir la référence ci-dessus). Puisque $A'/A$ est de longueur $1$, on en conclut que $$
A = \{f \in \kappa'[[t]] \mid f(0) \in \kappa\}
$$ Un calcul simple montre alors que $A$ est comme dans le cas (1).
\end{proof}

\begin{lemma}
\label{lemma-fitting-ideal-well-defined}
Soit $k$ un corps. Soit $A = k[[x_1, \ldots, x_n]]$. Soit $I = (f_1, \ldots, f_m) \subset A$ un idéal. Pour tout $r \geq 0$, l'idéal de $A/I$ engendré par les mineurs d'ordre $r \times r$ de la matrice $(\partial f_j/\partial x_i)$ est indépendant du choix des générateurs de $I$ et du système régulier de paramètres $x_1, \ldots, x_n$ de $A$.
\end{lemma}

\begin{proof}
La démonstration « naturelle » consiste à montrer que cet idéal est le $(n - r)$-ième idéal de Fitting d'un module de différentielles continues de $A/I$ sur $k$. Voici une démonstration directe. Si $g_1, \ldots g_l$ est un second système de générateurs de $I$, on peut écrire $g_s = \sum a_{sj}f_j$ et l'on a l'égalité de matrices $$
(\partial g_s/\partial x_i) = (a_{sj}) (\partial f_j/\partial x_i)
+ (\partial a_{sj}/\partial x_i f_j)
$$ Le dernier terme est nul dans $A/I$. Par la formule de Cauchy-Binet, l'idéal des mineurs associé aux $g_s$ est contenu dans celui associé aux $f_j$. Par symétrie, ces idéaux sont égaux. Si $y_1, \ldots, y_n \in \mathfrak m_A$ est un second système régulier de paramètres, la matrice $(\partial y_j/\partial x_i)$ est inversible et l'on peut utiliser la règle de dérivation des fonctions composées. Nous omettons quelques détails.
\end{proof}

\begin{lemma}
\label{lemma-fitting-ideal}
Soit $k$ un corps. Soit $A = k[[x_1, \ldots, x_n]]$. Soit $I = (f_1, \ldots, f_m) \subset \mathfrak m_A$ un idéal. Les conditions suivantes sont équivalentes : \begin{enumerate} \item $k \to A/I$ est comme dans le lemme \ref{lemma-nodal-algebraic}; \item $A/I$ est réduit et les mineurs d'ordre $(n - 1) \times (n - 1)$ de la matrice $(\partial f_j/\partial x_i)$ engendrent $I + \mathfrak m_A$; \item $\text{depth}(A/I) = 1$ et les mineurs d'ordre $(n - 1) \times (n - 1)$ de la matrice $(\partial f_j/\partial x_i)$ engendrent $I + \mathfrak m_A$.
\end{enumerate}
\end{lemma}

\begin{proof}
D'après le lemme \ref{lemma-fitting-ideal-well-defined}, nous pouvons changer le système de coordonnées et le choix des générateurs au cours de la démonstration.

\medskip\noindent
Si (1) est satisfaite, nous pouvons changer de coordonnées de sorte que $x_1, \ldots, x_{n - 2}$ s'envoient sur zéro dans $A/I$ et que $A/I = k[[x_{n - 1}, x_n]]/(a x_{n - 1}^2 + b x_{n - 1}x_n + c x_n^2)$ pour une certaine forme quadratique non dégénérée $a x_{n - 1}^2 + b x_{n - 1}x_n + c x_n^2$. Un calcul explicite montre alors que (2) et (3) sont vraies.

\medskip\noindent
Supposons que les mineurs d'ordre $(n - 1) \times (n - 1)$ de la matrice $(\partial f_j/\partial x_i)$ engendrent $I + \mathfrak m_A$. Supposons que, pour certains $i$ et $j$, la dérivée partielle $\partial f_j/\partial x_i$ soit une unité dans $A$. Nous pouvons alors prendre le système de paramètres $f_j, x_1, \ldots, x_{i - 1}, \hat x_i, x_{i + 1}, \ldots, x_n$ et les générateurs $f_j, f_1, \ldots, f_{j - 1}, \hat f_j, f_{j + 1}, \ldots, f_m$ de $I$. Nous obtenons ainsi un système régulier de paramètres $x_1, \ldots, x_n$ et des générateurs $x_1, f_2, \ldots, f_m$ de $I$. Cherchons ensuite $i \geq 2$ et $j \geq 2$ tels que $\partial f_j/\partial x_i$ soit une unité dans $A$. Si une telle paire existe, nous pouvons effectuer un remplacement comme ci-dessus et supposer que nous disposons d'un système régulier de paramètres $x_1, \ldots, x_n$ et de générateurs $x_1, x_2, f_3, \ldots, f_m$ de $I$. En poursuivant, nous parvenons en un nombre fini d'étapes à un système régulier de paramètres $x_1, \ldots, x_n$ et à des générateurs $x_1, \ldots, x_t, f_{t + 1}, \ldots, f_m$ de $I$ tels que $\partial f_j/\partial x_i \in \mathfrak m_A$ pour tous $i, j \geq t + 1$.

\medskip\noindent
Dans ce cas, la matrice des dérivées partielles a la forme par blocs suivante : $$
\left(
\begin{matrix}
I_{t \times t} & * \\
0 & \mathfrak m_A
\end{matrix}
\right)
$$ Tout mineur d'ordre $(n - 1) \times (n - 1)$ appartient donc à $\mathfrak m_A^{n - 1 - t}$. Remarquons que $I \not = \mathfrak m_A$, sans quoi l'idéal des mineurs contiendrait $1$. Il s'ensuit que $n - 1 - t \leq 1$, car il existe un élément de $\mathfrak m_A \setminus \mathfrak m_A^2 + I$ (sinon $I = \mathfrak m_A$ par le lemme de Nakayama). Ainsi $t \geq n - 2$. Nous avons vu ci-dessus que $t \not = n$; de même, si $t = n - 1$, il existerait un mineur inversible d'ordre $(n - 1) \times (n - 1)$, ce qui est également exclu. Par conséquent, $t = n - 2$. Alors $A/I$ est un quotient de $k[[x_{n - 1}, x_n]]$, et le lemme \ref{lemma-reduced-quotient-regular-ring-dim-2} implique, dans les deux cas (2) et (3), que $I$ est engendré par $x_1, \ldots, x_{n - 2}, f$ pour un certain $f = f(x_{n - 1}, x_n)$. Dans ce cas, la condition sur les mineurs signifie exactement que le terme quadratique de $f$ est non dégénéré, c'est-à-dire que $A/I$ est comme dans le lemme \ref{lemma-nodal-algebraic}.
\end{proof}

\begin{lemma}
\label{lemma-nodal}
Soit $k$ un corps. Soit $X$ un schéma algébrique de dimension $1$ sur $k$. Soit $x \in X$ un point fermé. Les conditions suivantes sont équivalentes : \begin{enumerate} \item $x$ est un nœud; \item $k \to \mathcal{O}_{X, x}$ est comme dans le lemme \ref{lemma-nodal-algebraic}; \item tout $\overline{x} \in X_{\overline{k}}$ s'envoyant sur $x$ définit une singularité nodale; \item $\kappa(x)/k$ est séparable, $\mathcal{O}_{X, x}$ est réduit, et le premier idéal de Fitting de $\Omega_{X/k}$ engendre $\mathfrak m_x$ dans $\mathcal{O}_{X, x}$; \item $\kappa(x)/k$ est séparable, $\text{depth}(\mathcal{O}_{X, x}) = 1$, et le premier idéal de Fitting de $\Omega_{X/k}$ engendre $\mathfrak m_x$ dans $\mathcal{O}_{X, x}$; \item $\kappa(x)/k$ est séparable et $\mathcal{O}_{X, x}$ est réduit, a pour invariant $\delta$ la valeur $1$, et possède $2$ branches géométriques.
\end{enumerate}
\end{lemma}

\begin{proof}
Supposons d'abord que $k$ soit algébriquement clos. Dans ce cas, l'équivalence de (1) et (3) est immédiate. L'équivalence de (1) et (3) avec (2) tient au fait que, à changement de coordonnées près, la seule forme quadratique non dégénérée en deux variables est $xy$. L'équivalence de (1) et (6) est le lemme \ref{lemma-multicross-algebra}. Après avoir remplacé $X$ par un voisinage affine de $x$, nous pouvons supposer qu'il existe une immersion fermée $X \to \mathbf{A}^n_k$ envoyant $x$ sur $0$. Soient $f_1, \ldots, f_m \in k[x_1, \ldots, x_n]$ des générateurs de l'idéal $I$ de $X$ dans $\mathbf{A}^n_k$. Alors $\Omega_{X/k}$ correspond au $R = k[x_1, \ldots, x_n]/I$-module $\Omega_{R/k}$ qui admet une présentation $$
R^{\oplus m} \xrightarrow{(\partial f_j/\partial x_i)}
R^{\oplus n} \to \Omega_{R/k} \to 0
$$ (voir Algèbre, sections \ref{algebra-section-differentials} et \ref{algebra-section-netherlander}). Le premier idéal de Fitting de $\Omega_{R/k}$ est donc l'idéal engendré par les mineurs d'ordre $(n - 1) \times (n - 1)$ de la matrice $(\partial f_j/\partial x_i)$. Le lemme \ref{lemma-fitting-ideal}, appliqué au complété de $k[x_1, \ldots, x_n] \to R$ en l'idéal maximal $(x_1, \ldots, x_n)$, montre ainsi que (2), (4) et (5) sont équivalentes.

\medskip\noindent
Supposons maintenant que $k$ soit un corps quelconque. Dans les cas (2), (4), (5) et (6), le corps résiduel $\kappa(x)$ est séparable sur $k$. Montrons qu'il en va de même dans les cas (1) et (3). Soit donc $Z \subset X$ le sous-schéma fermé de $X$ défini par le premier idéal de Fitting de $\Omega_{X/k}$. La formation de $Z$ commute aux extensions de corps (Diviseurs, lemme \ref{divisors-lemma-base-change-and-fitting-ideal-omega}). Si (1) ou (3) est vraie, il existe un point $\overline{x}$ de $X_{\overline{k}}$ tel que $\overline{x}$ soit un point isolé de multiplicité $1$ de $Z_{\overline{k}}$ (puisque nous disposons de l'équivalence des conditions du lemme sur $\overline{k}$). En particulier, $Z_{\overline{x}}$ est géométriquement réduit en $\overline{x}$ (car $\overline{k}$ est algébriquement clos). Ainsi $Z$ est géométriquement réduit en $x$ (Variétés, lemme \ref{varieties-lemma-geometrically-reduced-upstairs}). En particulier, $Z$ est réduit en $x$; par conséquent, $Z = \Spec(\kappa(x))$ dans un voisinage de $x$ et $\kappa(x)$ est géométriquement réduit sur $k$. Cela signifie que $\kappa(x)/k$ est séparable (Algèbre, lemme \ref{algebra-lemma-characterize-separable-field-extensions}).

\medskip\noindent
L'argument du paragraphe précédent montre que, si (1) ou (3) est satisfaite, le premier idéal de Fitting de $\Omega_{X/k}$ engendre $\mathfrak m_x$. Puisque $\mathcal{O}_{X, x} \to \mathcal{O}_{X_{\overline{k}}, \overline{x}}$ est plat et que $\mathcal{O}_{X_{\overline{k}}, \overline{x}}$ est réduit et de profondeur $1$, on voit que (4) et (5) sont satisfaites (utiliser Algèbre, lemmes \ref{algebra-lemma-descent-reduced} et \ref{algebra-lemma-apply-grothendieck}). Réciproquement, (4) implique (5) d'après Algèbre, lemme \ref{algebra-lemma-criterion-reduced}. Si (5) est satisfaite, alors $Z$ est géométriquement réduit en $x$ (car $\kappa(x)/k$ est séparable et $Z$ est égal à $x$ dans un voisinage). Ainsi $Z_{\overline{k}}$ est réduit en tout point $\overline{x}$ de $X_{\overline{k}}$ situé au-dessus de $x$. Autrement dit, le premier idéal de Fitting de $\Omega_{X_{\overline{k}}/\overline{k}}$ engendre $\mathfrak m_{\overline{x}}$ dans $\mathcal{O}_{X_{\overline{k}, \overline{x}}}$. De plus, puisque $\mathcal{O}_{X, x} \to \mathcal{O}_{X_{\overline{k}}, \overline{x}}$ est plat, on voit que $\text{depth}(\mathcal{O}_{X_{\overline{k}}, \overline{x}}) = 1$ (voir la référence ci-dessus). La condition (5) est donc satisfaite pour $\overline{x} \in X_{\overline{k}}$, et nous concluons que (3) est satisfaite (par l'équivalence sur les corps algébriquement clos). Nous avons ainsi établi l'équivalence de (1), (3), (4) et (5).

\medskip\noindent
L'équivalence de (2) et (6) résulte du lemme \ref{lemma-2-branches-delta-1}.

\medskip\noindent
Enfin, démontrons l'équivalence de (2) $=$ (6) avec (1) $=$ (3) $=$ (4) $=$ (5). Remarquons d'abord que le nombre géométrique de branches de $X$ en $x$ et celui de $X_{\overline{k}}$ en $\overline{x}$ sont égaux d'après Variétés, lemme \ref{varieties-lemma-geometric-branches-and-change-of-fields}. Les informations dont nous disposons à ce stade montrent que, dans tous les cas, ce nombre est égal à $2$. D'autre part, dans le cas (1), il est clair que $X$ est géométriquement réduit en $x$, et donc que $$
\text{invariant }\delta\text{ de }X\text{ en }x \leq
\text{invariant }\delta\text{ de }X_{\overline{k}}\text{ en }\overline{x}
$$ d'après Variétés, lemme \ref{varieties-lemma-delta-invariant-and-change-of-fields}. Puisque, dans le cas (1), le membre de droite vaut $1$, cela impose que l'invariant $\delta$ de $X$ en $x$ soit égal à $1$ (car, s'il était nul, $\mathcal{O}_{X, x}$ serait un anneau de valuation discrète d'après Variétés, lemme \ref{varieties-lemma-delta-invariant-is-zero}, donc serait unibranche, contradiction). Ainsi (5) est satisfaite. Réciproquement, si (2) $=$ (5) est vraie, alors les hypothèses (a), (b) et (c) de Variétés, lemme \ref{varieties-lemma-geometrically-normal-in-codim-1} sont satisfaites pour $x \in X$ d'après le lemme \ref{lemma-2-branches-delta-1}. Variétés, lemme \ref{varieties-lemma-delta-invariant-and-change-of-fields-better} s'applique donc et donne l'égalité dans l'inégalité affichée ci-dessus. Nous en concluons que (5) est satisfaite pour $\overline{x} \in X_{\overline{k}}$, et l'équivalence des conditions sur un corps algébriquement clos nous ramène au cas (1).
\end{proof}

\begin{remark}[L'extension quadratique associée à un nœud] \label{remark-quadratic-extension} Soit $k$ un corps. Soit $(A, \mathfrak m, \kappa)$ une $k$-algèbre locale noethérienne. Supposons soit que $(A, \mathfrak m, \kappa)$ soit comme dans le lemme \ref{lemma-nodal-algebraic}, soit que $A$ soit de Nagata comme dans le lemme \ref{lemma-2-branches-delta-1}, soit que $A$ soit complet et comme dans le lemme \ref{lemma-fitting-ideal}. Alors $A$ définit canoniquement une algèbre séparable de degré $2$ sur $\kappa$, notée $\kappa'$, comme suit : \begin{enumerate} \item soit $q = ax^2 + bxy + cy^2$ une forme quadratique non dégénérée comme dans le lemme \ref{lemma-nodal-algebraic}, les coordonnées $x, y$ étant choisies de sorte que $a \not = 0$, et posons $\kappa' = \kappa[x]/(ax^2 + bx + c)$; \item soit $A' \supset A$ la clôture intégrale de $A$ dans son anneau total des fractions, et posons $\kappa' = A'/\mathfrak m A'$; \item soit $\kappa'$ la $\kappa$-algèbre telle que $\text{Proj}(\bigoplus_{n \geq 0} \mathfrak m^n/\mathfrak m^{n + 1}) =
\Spec(\kappa')$. \end{enumerate} L'équivalence de (1) et (2) a été établie dans la démonstration du lemme \ref{lemma-2-branches-delta-1}. Nous omettons la démonstration de leur équivalence avec (3). Si $X$ est un $k$-schéma localement noethérien et si $x \in X$ est un point tel que $\mathcal{O}_{X, x} = A$, alors (3) montre que $\Spec(\kappa') = X^\nu \times_X \Spec(\kappa)$, où $\nu : X^\nu \to X$ est le morphisme de normalisation.
\end{remark}

\begin{remark}[Extension quadratique triviale] \label{remark-trivial-quadratic-extension} Soit $k$ un corps. Soit $(A, \mathfrak m, \kappa)$ comme dans la remarque \ref{remark-quadratic-extension}, et soit $\kappa'/\kappa$ l'algèbre séparable de degré $2$ associée. Les conditions suivantes sont équivalentes : \begin{enumerate} \item $\kappa' \cong \kappa \times \kappa$ comme $\kappa$-algèbre; \item la forme $q$ du lemme \ref{lemma-nodal-algebraic} peut être choisie égale à $xy$; \item $A$ possède deux branches; \item l'extension $A'/A$ du lemme \ref{lemma-2-branches-delta-1} possède deux idéaux maximaux; \item $A^\wedge \cong \kappa[[x, y]]/(xy)$ comme $k$-algèbre. \end{enumerate} L'équivalence de ces conditions a été établie dans la démonstration du lemme \ref{lemma-2-branches-delta-1}. Si $X$ est un $k$-schéma localement noethérien et si $x \in X$ est un point tel que $\mathcal{O}_{X, x} = A$, cela signifie exactement qu'il existe deux points $x_1, x_2$ de la normalisation $X^\nu$ au-dessus de $x$ et que $\kappa(x) = \kappa(x_1) = \kappa(x_2)$.
\end{remark}

\begin{definition}
\label{definition-split-node}
Soit $k$ un corps. Soit $X$ un schéma algébrique de dimension $1$ sur $k$. Soit $x \in X$ un point fermé. On dit que $x$ est un {\it nœud déployé} si $x$ est un nœud, si $\kappa(x) = k$, et si les assertions équivalentes de la remarque \ref{remark-trivial-quadratic-extension} sont satisfaites pour $A = \mathcal{O}_{X, x}$.
\end{definition}

\noindent
Formulons le lemme obligé qui rassemble ce que nous savons déjà de cette notion.

\begin{lemma}
\label{lemma-split-node}
Soit $k$ un corps. Soit $X$ un schéma algébrique de dimension $1$ sur $k$. Soit $x \in X$ un point fermé. Les conditions suivantes sont équivalentes : \begin{enumerate} \item $x$ est un nœud déployé; \item $x$ est un nœud et il existe exactement deux points $x_1, x_2$ de la normalisation $X^\nu$ au-dessus de $x$, avec $k = \kappa(x_1) = \kappa(x_2)$; \item $\mathcal{O}_{X, x}^\wedge \cong k[[x, y]]/(xy)$ comme $k$-algèbre; \item ajouter davantage ici.
\end{enumerate}
\end{lemma}

\begin{proof}
Cela résulte de la discussion de la remarque \ref{remark-trivial-quadratic-extension} et du lemme \ref{lemma-nodal}.
\end{proof}

\begin{lemma}
\label{lemma-node-field-extension}
Soit $K/k$ une extension de corps. Soit $X$ un $k$-schéma localement algébrique de dimension $1$. Soit $y \in X_K$ un point d'image $x \in X$. Les conditions suivantes sont équivalentes : \begin{enumerate} \item $x$ est un point fermé de $X$ et un nœud; \item $y$ est un point fermé de $Y$ et un nœud.
\end{enumerate}
\end{lemma}

\begin{proof}
Si $x$ est un point fermé de $X$, alors $y$ l'est aussi (considérer les corps résiduels). La réciproque n'est toutefois pas nécessairement vraie : il peut arriver qu'un point fermé de $Y$ s'envoie sur un point non fermé de $X$. Dans ce cas, cependant, $y$ ne peut pas être un nœud. En effet, $X$ serait alors géométriquement unibranche en $x$ (car $x$ serait un point générique de $X$, tandis que $\mathcal{O}_{X, x}$ serait artinien et que tout anneau local artinien est géométriquement unibranche); par conséquent, $Y$ serait géométriquement unibranche en $y$ (Variétés, lemme \ref{varieties-lemma-geometrically-unibranch-and-change-of-fields}), ce qui signifie que $y$ ne peut pas être un nœud d'après le lemme \ref{lemma-nodal}. Nous pouvons donc supposer, et nous supposerons, que $x$ et $y$ sont tous deux des points fermés.

\medskip\noindent
Choisissons des clôtures algébriques $\overline{k}$, $\overline{K}$ et une application $\overline{k} \to \overline{K}$ prolongeant l'application donnée $k \to K$. En utilisant l'équivalence de (1) et (3) dans le lemme \ref{lemma-nodal}, nous nous ramenons au cas où $k$ et $K$ sont algébriquement clos. Dans ce cas, nous pouvons raisonner comme dans la démonstration du lemme \ref{lemma-nodal} pour voir que le nombre géométrique de branches et les invariants $\delta$ de $X$ en $x$ et de $Y$ en $y$ sont les mêmes. On peut aussi choisir un isomorphisme $k[[x_1, \ldots, x_n]]/(g_1, \ldots, g_m) \to \mathcal{O}_{X, x}^\wedge$ de $k$-algèbres comme dans Variétés, lemme \ref{varieties-lemma-complete-local-ring-structure-as-algebra}. D'après Variétés, lemme \ref{varieties-lemma-base-change-complete-local-ring}, il induit un isomorphisme $K[[x_1, \ldots, x_n]]/(g_1, \ldots, g_m) \to \mathcal{O}_{Y, y}^\wedge$ de $K$-algèbres. Par définition, il faut montrer que $$
k[[x_1, \ldots, x_n]]/(g_1, \ldots, g_m) \cong k[[s, t]]/(st)
$$ si et seulement si $$
K[[x_1, \ldots, x_n]]/(g_1, \ldots, g_m) \cong K[[s, t]]/(st)
$$ Nous invitons le lecteur à le démontrer. Puisque $k$ et $K$ sont des corps algébriquement clos, cela revient à demander que ces anneaux soient comme dans le lemme \ref{lemma-nodal-algebraic}. Par le lemme \ref{lemma-fitting-ideal}, cette condition se traduit en une assertion sur les mineurs d'ordre $(n - 1) \times (n - 1)$ de la matrice $(\partial g_j/\partial x_i)$, manifestement indépendante du corps utilisé. Nous omettons les détails.
\end{proof}

\begin{lemma}
\label{lemma-node-etale-local}
Soit $k$ un corps. Soit $X$ un $k$-schéma localement algébrique de dimension $1$. Soit $Y \to X$ un morphisme étale. Soit $y \in Y$ un point d'image $x \in X$. Les conditions suivantes sont équivalentes : \begin{enumerate} \item $x$ est un point fermé de $X$ et un nœud; \item $y$ est un point fermé de $Y$ et un nœud.
\end{enumerate}
\end{lemma}

\begin{proof}
D'après le lemme \ref{lemma-node-field-extension}, nous pouvons effectuer un changement de base vers la clôture algébrique de $k$. Les corps résiduels de $x$ et $y$ sont alors $k$. L'application $\mathcal{O}_{X, x}^\wedge \to \mathcal{O}_{Y, y}^\wedge$ est donc un isomorphisme (par exemple d'après Morphismes étales, lemme \ref{etale-lemma-characterize-etale-completions} ou Compléments d'algèbre, lemme \ref{more-algebra-lemma-flat-unramified}). Le lemme est dès lors clair.
\end{proof}

\begin{lemma}
\label{lemma-node-over-separable-extension}
Soit $k'/k$ une extension finie séparable de corps. Soit $X$ un $k'$-schéma localement algébrique de dimension $1$. Soit $x \in X$ un point fermé. Les conditions suivantes sont équivalentes : \begin{enumerate} \item $x$ est un nœud; \item $x$ est un nœud lorsque $X$ est considéré comme un $k$-schéma localement algébrique.
\end{enumerate}
\end{lemma}

\begin{proof}
Cela résulte immédiatement de la caractérisation des nœuds du lemme \ref{lemma-nodal}.
\end{proof}

\begin{lemma}
\label{lemma-nodal-lci}
Soit $k$ un corps. Soit $X$ un $k$-schéma localement algébrique équidimensionnel de dimension $1$. Les conditions suivantes sont équivalentes : \begin{enumerate} \item les singularités de $X$ sont au plus nodales; \item $X$ est une intersection complète locale sur $k$ et le sous-schéma fermé $Z \subset X$ découpé par le premier idéal de Fitting de $\Omega_{X/k}$ est non ramifié sur $k$.
\end{enumerate}
\end{lemma}

\begin{proof}
Nous invitons vivement le lecteur à trouver sa propre démonstration de ce lemme; celle qui suit ne fait que réunir des résultats antérieurs et risque de masquer ce qui se passe réellement.

\medskip\noindent
Supposons (2). Puisque $Z \to \Spec(k)$ est quasi-fini (Morphismes, lemme \ref{morphisms-lemma-unramified-quasi-finite}), les corps résiduels des points $x \in Z$ sont finis sur $k$ (et séparables), d'après Morphismes, lemme \ref{morphisms-lemma-residue-field-quasi-finite}. Chaque $x \in Z$ est donc un point fermé de $X$ d'après Morphismes, lemme \ref{morphisms-lemma-algebraic-residue-field-extension-closed-point-fibre}. L'anneau local $\mathcal{O}_{X, x}$ est de Cohen--Macaulay d'après Algèbre, lemme \ref{algebra-lemma-lci-CM}. Comme $\dim(\mathcal{O}_{X, x}) = 1$ par la théorie de la dimension (Variétés, section \ref{varieties-section-algebraic-schemes}), nous en concluons que $\text{depth}(\mathcal{O}_{X, x})) = 1$. Ainsi $x$ est un nœud d'après le lemme \ref{lemma-nodal}. Si $x \in X$, $x \not \in Z$, alors $X \to \Spec(k)$ est lisse en $x$ d'après Diviseurs, lemme \ref{divisors-lemma-d-fitting-ideal-omega-smooth}.

\medskip\noindent
Supposons (1). Sous cette hypothèse, $X$ est géométriquement réduit en tout point fermé (voir Variétés, lemme \ref{varieties-lemma-geometrically-reduced-upstairs}). Ainsi $X \to \Spec(k)$ est lisse sur un ouvert dense d'après Variétés, lemme \ref{varieties-lemma-geometrically-reduced-dense-smooth-open}. Par conséquent, $Z$ est fermé et formé de points fermés. D'après Diviseurs, lemme \ref{divisors-lemma-d-fitting-ideal-omega-smooth}, le morphisme $X \setminus Z \to \Spec(k)$ est lisse. Ainsi $X \setminus Z$ est une intersection complète locale d'après Morphismes, lemme \ref{morphisms-lemma-smooth-syntomic} et la définition d'une intersection complète locale de Morphismes, définition \ref{morphisms-definition-syntomic}. D'après le lemme \ref{lemma-nodal}, pour tout point $x \in Z$, l'anneau local $\mathcal{O}_{Z, x}$ est égal à $\kappa(x)$ et $\kappa(x)$ est séparable sur $k$. Ainsi $Z \to \Spec(k)$ est non ramifié (Morphismes, lemme \ref{morphisms-lemma-unramified-over-field}). Enfin, le lemme \ref{lemma-nodal}, par la partie (3) du lemme \ref{lemma-nodal-algebraic}, montre que $\mathcal{O}_{X, x}$ est une intersection complète au sens d'Algèbre à puissances divisées, définition \ref{dpa-definition-lci}. Or Algèbre à puissances divisées, lemme \ref{dpa-lemma-check-lci-agrees} et Morphismes, lemme \ref{morphisms-lemma-local-complete-intersection} montrent que cette notion coïncide avec celle utilisée pour définir un schéma intersection complète locale sur un corps, ce qui achève la démonstration.
\end{proof}

\begin{lemma}
\label{lemma-facts-about-nodal-curves}
Soit $k$ un corps. Soit $X$ un $k$-schéma localement algébrique équidimensionnel de dimension $1$ dont les singularités sont au plus nodales. Alors $X$ est de Gorenstein et géométriquement réduit.
\end{lemma}

\begin{proof}
L'assertion de Gorenstein résulte du lemme \ref{lemma-nodal-lci} et de Dualité pour les schémas, lemme \ref{duality-lemma-gorenstein-lci}. On peut aussi utiliser le fait qu'il suffit de vérifier l'assertion après passage à la clôture algébrique (Dualité pour les schémas, lemme \ref{duality-lemma-gorenstein-base-change}), puis qu'un anneau local noethérien est de Gorenstein si et seulement si son complété l'est (Complexes dualisants, lemme \ref{dualizing-lemma-flat-under-gorenstein}), et enfin démontrer directement que les anneaux locaux $k[[t]]$ et $k[[x, y]]/(xy)$ sont de Gorenstein.

\medskip\noindent
Pour voir que $X$ est géométriquement réduit, il suffit de montrer que $X_{\overline{k}}$ est réduit (Variétés, lemmes \ref{varieties-lemma-perfect-reduced} et \ref{varieties-lemma-geometrically-reduced}). Or $X_{\overline{k}}$ est une courbe nodale sur un corps algébriquement clos. Les anneaux locaux complets de $X_{\overline{k}}$ sont donc isomorphes soit à $\overline{k}[[t]]$, soit à $\overline{k}[[x, y]]/(xy)$, qui sont réduits, comme voulu.
\end{proof}

\begin{lemma}
\label{lemma-closed-subscheme-nodal-curve}
Soit $k$ un corps. Soit $X$ un $k$-schéma localement algébrique équidimensionnel de dimension $1$ dont les singularités sont au plus nodales. Si $Y \subset X$ est un sous-schéma fermé réduit équidimensionnel de dimension $1$, alors \begin{enumerate} \item les singularités de $Y$ sont au plus nodales; \item si $Z \subset X$ est l'adhérence schématique de $X \setminus Y$, alors \begin{enumerate} \item l'intersection schématique $Y \cap Z$ est la réunion disjointe de spectres d'extensions finies séparables de $k$; \item tout point de $Y \cap Z$ est un nœud de $X$; \item $Y \to \Spec(k)$ est lisse en tout point de $Y \cap Z$.
\end{enumerate}
\end{enumerate}
\end{lemma}

\begin{proof}
Puisque $X$ et $Y$ sont réduits et équidimensionnels de dimension $1$, on voit que $Y$ est la réunion schématique d'un sous-ensemble des composantes irréductibles de $X$ (dans un anneau réduit, $(0)$ est l'intersection des idéaux premiers minimaux). Soit $y \in Y$ un point fermé. Si $y$ appartient au lieu lisse de $X \to \Spec(k)$, alors $y$ se trouve sur une unique composante irréductible de $X$, et $Y$ et $X$ coïncident dans un voisinage ouvert de $y$. Ainsi $Y \to \Spec(k)$ est lisse en $y$. Si $y$ est un nœud de $X$ mais appartient encore à une unique composante irréductible de $X$, le même raisonnement montre que $y$ est un nœud de $Y$. Supposons que $y$ appartienne à plus de $1$ composante irréductible de $X$. Puisque le nombre de branches géométriques de $X$ en $y$ est $2$ d'après le lemme \ref{lemma-nodal}, au plus $2$ composantes irréductibles peuvent passer par $y$, d'après Propriétés, lemme \ref{properties-lemma-number-of-branches-irreducible-components}. Si $Y$ contient ces deux composantes, alors $Y = X$ dans un voisinage ouvert de $y$, et $y$ est un nœud de $Y$. Supposons enfin que $Y$ ne contienne qu'une des composantes irréductibles. Après avoir remplacé $X$ par un voisinage ouvert de $x$, nous pouvons supposer que $Y$ est l'une des deux composantes irréductibles et que $Z$ est l'autre. D'après Propriétés, lemme \ref{properties-lemma-number-of-branches-irreducible-components}, on voit encore que $X$ possède deux branches en $y$, c'est-à-dire que l'anneau local $\mathcal{O}_{X, y}$ possède deux branches et que celles-ci proviennent de $\mathcal{O}_{Y, y}$ et $\mathcal{O}_{Z, y}$. Écrivons $\mathcal{O}_{X, y}^\wedge \cong \kappa(y)[[u, v]]/(uv)$ comme dans la remarque \ref{remark-trivial-quadratic-extension}. Le corps $\kappa(y)$ est fini séparable sur $k$, par exemple d'après le lemme \ref{lemma-nodal}. Ainsi, après avoir éventuellement interverti les rôles de $u$ et $v$, le complété de l'application $\mathcal{O}_{X, y} \to \mathcal{O}_{Y, Y}$ correspond à $\kappa(y)[[u, v]]/(uv) \to \kappa(y)[[u]]$ et celui de l'application $\mathcal{O}_{X, y} \to \mathcal{O}_{Y, Y}$ correspond à $\kappa(y)[[u, v]]/(uv) \to \kappa(y)[[v]]$. L'intersection schématique de $Y \cap Z$ est découpée par la somme de leurs idéaux, qui devient $(u, v)$ dans le complété, c'est-à-dire l'idéal maximal. Les assertions (2)(a) et (2)(b) sont donc claires. Enfin, (2)(c) est satisfaite : le complété de $\mathcal{O}_{Y, y}$ est régulier, donc $\mathcal{O}_{Y, y}$ est régulier (Compléments d'algèbre, lemme \ref{more-algebra-lemma-completion-regular}), et $\kappa(y)/k$ est séparable; on obtient donc la lissité dans un voisinage ouvert d'après Algèbre, lemme \ref{algebra-lemma-separable-smooth}.
\end{proof}




\section{Familles de courbes nodales}
\label{section-families-nodal}

\noindent
Dans le projet Stacks, les courbes sont des variétés irréductibles de dimension $1$. Dans la littérature, toutefois, une « courbe semi-stable » ou une « courbe nodale » n'est généralement pas irréductible et est souvent supposée propre, en particulier dans des expressions telles que « famille de courbes semi-stables », « famille de courbes nodales » ou « famille nodale ». Il nous est donc quelque peu difficile de choisir une terminologie à la fois conforme à la littérature et cohérente en interne. De plus, nous souhaitons réellement commencer par étudier la notion introduite dans le lemme suivant, qui est locale sur la source.

\begin{lemma}
\label{lemma-nodal-family}
Soit $f : X \to S$ un morphisme de schémas. Les conditions suivantes sont équivalentes : \begin{enumerate} \item $f$ est plat et localement de présentation finie, toute fibre non vide $X_s$ est équidimensionnelle de dimension $1$, et les singularités de $X_s$ sont au plus nodales; \item $f$ est syntomique de dimension relative $1$ et le sous-schéma fermé $\text{Sing}(f) \subset X$ défini par le premier idéal de Fitting de $\Omega_{X/S}$ est non ramifié sur $S$.
\end{enumerate}
\end{lemma}

\begin{proof}
Rappelons que la formation de $\text{Sing}(f)$ commute au changement de base; voir Diviseurs, lemme \ref{divisors-lemma-base-change-and-fitting-ideal-omega}. Le lemme résulte donc du lemme \ref{lemma-nodal-lci} et de Morphismes, lemmes \ref{morphisms-lemma-syntomic-flat-fibres} et \ref{morphisms-lemma-unramified-etale-fibres}. (Nous utilisons aussi les triviaux Morphismes, lemmes \ref{morphisms-lemma-syntomic-locally-finite-presentation} et \ref{morphisms-lemma-syntomic-flat}.)
\end{proof}

\begin{definition}
\label{definition-nodal-family}
Soit $f : X \to S$ un morphisme de schémas. On dit que $f$ est {\it à singularités au plus nodales de dimension relative $1$} si $f$ satisfait les conditions équivalentes du lemme \ref{lemma-nodal-family}.
\end{definition}

\noindent
Voici quelques raisons qui expliquent cette terminologie pesante\footnote{Mais n'hésitez pas à écrire au responsable du projet Stacks si vous avez une meilleure suggestion.}. Premièrement, nous voulons éviter toute confusion entre cette notion et les autres notions de la littérature (voir l'introduction de cette section). Deuxièmement, nous pouvons imaginer plusieurs généralisations de cette notion à des morphismes de dimension relative supérieure : on peut, par exemple, demander que les morphismes soient étale-localement des composés de morphismes à singularités au plus nodales, ou demander que leurs fibres soient de dimension supérieure mais n'aient au pire que des points doubles ordinaires.

\begin{lemma}
\label{lemma-smooth-relative-dimension-1}
Un morphisme lisse de dimension relative $1$ est à singularités au plus nodales de dimension relative $1$.
\end{lemma}

\begin{proof}
Omis.
\end{proof}

\begin{lemma}
\label{lemma-base-change-nodal-family}
Soit $f : X \to S$ à singularités au plus nodales de dimension relative $1$. Il en va de même après tout changement de base de $f$.
\end{lemma}

\begin{proof}
Cela tient au fait que le changement de base d'un morphisme syntomique est syntomique (Morphismes, lemme \ref{morphisms-lemma-base-change-syntomic}), que le changement de base d'un morphisme de dimension relative $1$ est de dimension relative $1$ (Morphismes, lemme \ref{morphisms-lemma-base-change-relative-dimension-d}), que la formation de $\text{Sing}(f)$ commute au changement de base (Diviseurs, lemme \ref{divisors-lemma-base-change-and-fitting-ideal-omega}), et que le changement de base d'un morphisme non ramifié est non ramifié (Morphismes, lemme \ref{morphisms-lemma-base-change-unramified}).
\end{proof}

\noindent
Le lemme suivant montre que l'on peut vérifier fibre par fibre qu'un morphisme est à singularités au plus nodales de dimension relative $1$.

\begin{lemma}
\label{lemma-locus-where-nodal}
Soit $f : X \to S$ un morphisme de schémas, plat et localement de présentation finie. Il existe un plus grand sous-schéma ouvert $U \subset X$ tel que $f|_U : U \to S$ soit à singularités au plus nodales de dimension relative $1$. De plus, la formation de $U$ commute à tout changement de base.
\end{lemma}

\begin{proof}
D'après Morphismes, lemme \ref{morphisms-lemma-set-points-where-fibres-lci}, il existe un tel ouvert sur lequel $f$ est syntomique. Nous pouvons donc supposer que $f$ est un morphisme syntomique. En particulier, $f$ est un morphisme de Cohen--Macaulay (Dualité pour les schémas, lemmes \ref{duality-lemma-lci-gorenstein} et \ref{duality-lemma-gorenstein-CM-morphism}). Ainsi $X$ est la réunion disjointe de sous-schémas ouverts et fermés sur lesquels $f$ a une dimension relative fixée; voir Morphismes, lemme \ref{morphisms-lemma-flat-finite-presentation-CM-fibres-relative-dimension}. Cette décomposition est préservée par tout changement de base; voir Morphismes, lemme \ref{morphisms-lemma-base-change-relative-dimension-d}. En ne gardant qu'un morceau, nous pouvons supposer que $f$ est syntomique de dimension relative $1$. Soit $\text{Sing}(f) \subset X$ le sous-schéma fermé défini par le premier idéal de Fitting de $\Omega_{X/S}$. Il existe un plus grand sous-schéma ouvert $W \subset \text{Sing}(f)$ tel que $W \to S$ soit non ramifié, et sa formation commute au changement de base (Morphismes, lemme \ref{morphisms-lemma-set-points-where-fibres-unramified}). Puisque la formation de $\text{Sing}(f)$ commute elle aussi au changement de base (Diviseurs, lemme \ref{divisors-lemma-base-change-and-fitting-ideal-omega}), nous voyons que $$
U = (X \setminus \text{Sing}(f)) \cup W
$$ est le plus grand sous-schéma ouvert de $X$ tel que $f|_U : U \to S$ soit à singularités au plus nodales de dimension relative $1$, et que la formation de $U$ commute au changement de base.
\end{proof}

\begin{lemma}
\label{lemma-nodal-family-precompose-etale}
Soit $f : X \to S$ à singularités au plus nodales de dimension relative $1$. Si $Y \to X$ est un morphisme étale, alors le composé $g : Y \to S$ est à singularités au plus nodales de dimension relative $1$.
\end{lemma}

\begin{proof}
Remarquons que $g$ est plat et localement de présentation finie, comme composé de morphismes plats et localement de présentation finie (utiliser Morphismes, lemmes \ref{morphisms-lemma-etale-locally-finite-presentation}, \ref{morphisms-lemma-etale-flat}, \ref{morphisms-lemma-composition-finite-presentation} et \ref{morphisms-lemma-composition-flat}). Il suffit donc de démontrer que ses fibres ont des singularités au plus nodales. Cela résulte du lemme \ref{lemma-node-etale-local} (et du fait que le composé d'un morphisme étale et d'un morphisme lisse est lisse, d'après Morphismes, lemmes \ref{morphisms-lemma-etale-smooth-unramified} et \ref{morphisms-lemma-composition-smooth}).
\end{proof}

\begin{lemma}
\label{lemma-nodal-family-postcompose-etale}
Soit $S' \to S$ un morphisme étale de schémas. Soit $f : X \to S'$ à singularités au plus nodales de dimension relative $1$. Alors le composé $g : X \to S$ est à singularités au plus nodales de dimension relative $1$.
\end{lemma}

\begin{proof}
Remarquons que $g$ est plat et localement de présentation finie, comme composé de morphismes plats et localement de présentation finie (utiliser Morphismes, lemmes \ref{morphisms-lemma-etale-locally-finite-presentation}, \ref{morphisms-lemma-etale-flat}, \ref{morphisms-lemma-composition-finite-presentation} et \ref{morphisms-lemma-composition-flat}). Il suffit donc de démontrer que les fibres de $g$ ont des singularités au plus nodales. Cela résulte du lemme \ref{lemma-node-over-separable-extension} et du résultat analogue pour les points lisses.
\end{proof}

\begin{lemma}
\label{lemma-nodal-family-etale-local-source}
Soit $f : X \to S$ un morphisme de schémas. Soit $\{U_i \to X\}$ un recouvrement étale. Les conditions suivantes sont équivalentes : \begin{enumerate} \item $f$ est à singularités au plus nodales de dimension relative $1$; \item chaque $U_i \to S$ est à singularités au plus nodales de dimension relative $1$. \end{enumerate} Autrement dit, la propriété d'être à singularités au plus nodales de dimension relative $1$ est locale pour la topologie étale sur la source.
\end{lemma}

\begin{proof}
Une implication a été établie dans le lemme \ref{lemma-nodal-family-precompose-etale}. Pour l'autre, remarquons que les propriétés d'être localement de présentation finie, d'être plat et d'avoir une dimension relative $1$ sont locales pour la topologie étale sur la source (Descente, lemmes \ref{descent-lemma-locally-finite-presentation-fppf-local-source}, \ref{descent-lemma-flat-fpqc-local-source} et \ref{descent-lemma-dimension-at-point}). En prenant les fibres, nous nous ramenons au cas où $S$ est le spectre d'un corps. Le résultat découle alors du lemme \ref{lemma-node-etale-local} (et du fait que la lissité est locale pour la topologie étale sur la source, d'après Descente, lemme \ref{descent-lemma-smooth-smooth-local-source}).
\end{proof}

\begin{lemma}
\label{lemma-nodal-family-fpqc-local-target}
Soit $f : X \to S$ un morphisme de schémas. Soit $\{U_i \to S\}$ un recouvrement fpqc. Les conditions suivantes sont équivalentes : \begin{enumerate} \item $f$ est à singularités au plus nodales de dimension relative $1$; \item chaque $X \times_S U_i \to U_i$ est à singularités au plus nodales de dimension relative $1$. \end{enumerate} Autrement dit, la propriété d'être à singularités au plus nodales de dimension relative $1$ est locale pour la topologie fpqc sur la cible.
\end{lemma}

\begin{proof}
Une implication a été établie dans le lemme \ref{lemma-base-change-nodal-family}. Pour l'autre, remarquons que les propriétés d'être localement de présentation finie, d'être plat et d'avoir une dimension relative $1$ sont locales pour la topologie fpqc sur la cible (Descente, lemmes \ref{descent-lemma-descending-property-locally-finite-presentation} et \ref{descent-lemma-descending-property-flat}, et Morphismes, lemme \ref{morphisms-lemma-dimension-fibre-after-base-change}). En prenant les fibres, nous nous ramenons au cas où $S$ est le spectre d'un corps. Le résultat découle alors du lemme \ref{lemma-node-field-extension} (et du fait que la lissité est locale pour la topologie fpqc sur la cible, d'après Descente, lemme \ref{descent-lemma-descending-property-smooth}).
\end{proof}

\begin{lemma}
\label{lemma-descend-nodal-family}
Soit $S = \lim S_i$ la limite d'un système inductif de schémas dont les morphismes de transition sont affines. Soit $0 \in I$, et soit $f_0 : X_0 \to Y_0$ un morphisme de schémas au-dessus de $S_0$. Supposons $S_0$, $X_0$ et $Y_0$ quasi-compacts et quasi-séparés. Soit $f_i : X_i \to Y_i$ le changement de base de $f_0$ à $S_i$, et soit $f : X \to Y$ le changement de base de $f_0$ à $S$. Si \begin{enumerate} \item $f$ est à singularités au plus nodales de dimension relative $1$; \item $f_0$ est localement de présentation finie, \end{enumerate} alors il existe $i \geq 0$ tel que $f_i$ soit à singularités au plus nodales de dimension relative $1$.
\end{lemma}

\begin{proof}
D'après Limites, lemme \ref{limits-lemma-descend-syntomic}, il existe $i$ tel que $f_i$ soit syntomique. Alors $X_i = \coprod_{d \geq 0} X_{i, d}$ est la réunion disjointe de sous-schémas ouverts et fermés tels que $X_{i, d} \to Y_i$ ait une dimension relative $d$; voir Morphismes, lemme \ref{morphisms-lemma-syntomic-relative-dimension}. Le comportement de la dimension des fibres par changement de base donné dans Morphismes, lemme \ref{morphisms-lemma-dimension-fibre-after-base-change} montre que $X \to X_i$ s'envoie dans $X_{i, 1}$. Il existe alors $i' \geq i$ tel que $X_{i'} \to X_i$ s'envoie dans $X_{i, 1}$; voir Limites, lemme \ref{limits-lemma-limit-contained-in-constructible}. Ainsi $f_{i'} : X_{i'} \to Y_{i'}$ est syntomique de dimension relative $1$ (de nouveau par Morphismes, lemme \ref{morphisms-lemma-dimension-fibre-after-base-change}). Considérons le morphisme $\text{Sing}(f_{i'}) \to Y_{i'}$. Nous savons que son changement de base à $Y$ est non ramifié. D'après Limites, lemme \ref{limits-lemma-descend-unramified}, après avoir augmenté $i'$, le morphisme $\text{Sing}(f_{i'}) \to Y_{i'}$ devient donc non ramifié. Cela achève la démonstration.
\end{proof}

\begin{lemma}
\label{lemma-formal-local-structure-nodal-family}
Soit $f : T \to S$ un morphisme de schémas. Soit $t \in T$, d'image $s \in S$. Supposons que \begin{enumerate} \item $f$ soit plat en $t$; \item $\mathcal{O}_{S, s}$ soit noethérien; \item $f$ soit localement de type fini; \item $t$ soit un nœud déployé de la fibre $T_s$. \end{enumerate} Alors il existe $h \in \mathfrak m_s^\wedge$ et un isomorphisme $$
\mathcal{O}_{T, t}^\wedge \cong
\mathcal{O}_{S, s}^\wedge[[x, y]]/(xy - h)
$$ de $\mathcal{O}_{S, s}^\wedge$-algèbres.
\end{lemma}

\begin{proof}
Remplaçons $S$ par $\Spec(\mathcal{O}_{S, s})$ et $T$ par son changement de base à $\Spec(\mathcal{O}_{S, s})$. Alors $T$ est localement noethérien, et donc $\mathcal{O}_{T, t}$ est noethérien. Posons $A = \mathcal{O}_{S, s}^\wedge$, $\mathfrak m = \mathfrak m_A$ et $B = \mathcal{O}_{T, t}^\wedge$. D'après Compléments d'algèbre, lemme \ref{more-algebra-lemma-flat-completion}, $A \to B$ est plat. Puisque $\mathcal{O}_{T, t}/\mathfrak m_s \mathcal{O}_{T, t} = \mathcal{O}_{T_s, t}$, on voit que $B/\mathfrak m B = \mathcal{O}_{T_s, t}^\wedge$. L'hypothèse (4) et le lemme \ref{lemma-split-node} montrent qu'il existe $\overline{u}, \overline{v} \in B/\mathfrak m B$ tels que l'application $$
(A/\mathfrak m)[[x, y]] \longrightarrow B/\mathfrak m B,\quad
x \longmapsto \overline{u},
x \longmapsto \overline{v}
$$ soit surjective, de noyau $(xy)$.

\medskip\noindent
Supposons donnés $n \geq 1$ et $u, v \in B$ s'envoyant sur $\overline{u}, \overline{v}$ tels que $$
u v = h + \delta
$$ pour certains $h \in A$ et $\delta \in \mathfrak m^nB$. Nous affirmons qu'il existe $u', v' \in B$ avec $u - u', v - v' \in \mathfrak m^n B$ tels que $$
u' v' = h' + \delta'
$$ pour certains $h' \in A$ et $\delta' \in \mathfrak m^{n + 1}B$. Pour le voir, écrivons $\delta = \sum f_i b_i$ avec $f_i \in \mathfrak m^n$ et $b_i \in B$. Écrivons ensuite $b_i = a_i + u b_{i, 1} + v b_{i, 2} + \delta_i$ avec $a_i \in A$, $b_{i, 1}, b_{i, 2} \in B$ et $\delta_i \in \mathfrak m B$. C'est possible parce que le corps résiduel de $B$ coïncide avec celui de $A$, et que les images de $u$ et $v$ dans $B/\mathfrak m B$ engendrent l'idéal maximal. Posons alors $$
u' = u - \sum b_{i, 2}f_i,\quad
v' = v - \sum b_{i, 1}f_i
$$; nous obtenons $$
u'v' = h + \delta - \sum (b_{i, 1}u + b_{i, 2}v)f_i + \sum
c_{ij}f_if_j =
h + \sum a_if_i + \sum f_i \delta_i + \sum c_{ij}f_if_j
$$ pour un certain $c_{i, j} \in B$. Nous avons donc une formule comme ci-dessus avec $h' = h + \sum a_if_i$ et $\delta' = \sum f_i \delta_i + \sum c_{ij}f_if_j$.

\medskip\noindent
En raisonnant par récurrence à partir de relèvements quelconques $u_1, v_1 \in B$ de $\overline{u}, \overline{v}$, le paragraphe précédent fournit des suites d'éléments $u_n, v_n \in B$ et $h_n \in A$ telles que $u_n - u_{n + 1} \in \mathfrak m^n B$, $v_n - v_{n + 1} \in \mathfrak m^n B$, $h_n - h_{n + 1} \in \mathfrak m^n$, et telles que $u_n v_n - h_n \in \mathfrak m^n B$. Puisque $A$ et $B$ sont complets, nous pouvons poser $u_\infty = \lim u_n$, $v_\infty = \lim v_n$ et $h_\infty = \lim h_n$, et nous obtenons alors $u_\infty v_\infty = h_\infty$ dans $B$. Nous disposons donc d'un homomorphisme de $A$-algèbres $$
A[[x, y]]/(xy - h_\infty) \longrightarrow B
$$ envoyant $x$ sur $u_\infty$ et $v$ sur $v_\infty$. C'est une application d'algèbres $A$-plates qui devient un isomorphisme après division par $\mathfrak m$. Elle est surjective modulo $\mathfrak m$, donc surjective par complétude et par Algèbre, lemme \ref{algebra-lemma-completion-generalities}. Nous pouvons alors appliquer Algèbre, lemme \ref{algebra-lemma-mod-injective} pour conclure que c'est un isomorphisme.
\end{proof}

\begin{lemma}
\label{lemma-genus-in-nodal-family-of-curves}
Soit $f : X \to S$ un morphisme de schémas. Supposons que \begin{enumerate} \item $f$ soit propre; \item $f$ soit à singularités au plus nodales de dimension relative $1$; \item les fibres géométriques de $f$ soient connexes. \end{enumerate} Alors (a) $f_*\mathcal{O}_X = \mathcal{O}_S$, et cela reste vrai après tout changement de base; (b) $R^1f_*\mathcal{O}_X$ est un $\mathcal{O}_S$-module fini localement libre dont la formation commute à tout changement de base; et (c) $R^qf_*\mathcal{O}_X = 0$ pour $q \geq 2$.
\end{lemma}

\begin{proof}
La partie (a) résulte de Catégories dérivées des schémas, lemme \ref{perfect-lemma-proper-flat-geom-red-connected}. D'après Catégories dérivées des schémas, lemme \ref{perfect-lemma-proper-flat-h0}, localement sur $S$, nous pouvons écrire $Rf_*\mathcal{O}_X = \mathcal{O}_S \oplus P$, où $P$ est parfait et d'amplitude de Tor contenue dans $[1, \infty)$. Rappelons que la formation de $Rf_*\mathcal{O}_X$ commute à tout changement de base (Catégories dérivées des schémas, lemme \ref{perfect-lemma-flat-proper-perfect-direct-image-general}). Ainsi, pour $s \in S$, on a $$
H^i(P \otimes_{\mathcal{O}_S}^\mathbf{L} \kappa(s)) =
H^i(X_s, \mathcal{O}_{X_s})
\text{ pour }i \geq 1
$$ Ce groupe est nul sauf si $i = 1$, puisque $X_s$ est un schéma noethérien de dimension $1$; voir Cohomologie, proposition \ref{cohomology-proposition-vanishing-Noetherian}. Alors $P = H^1(P)[-1]$ et $H^1(P)$ est fini localement libre, par exemple d'après Compléments d'algèbre, lemme \ref{more-algebra-lemma-lift-perfect-from-residue-field}. Comme tout est compatible au changement de base, nous concluons.
\end{proof}





\section{Structure locale étale des familles nodales}
\label{section-etale-local-nodal}


\noindent
Considérons le morphisme de schémas $$
\Spec(\mathbf{Z}[u, v, a]/(uv - a))
\longrightarrow
\Spec(\mathbf{Z}[a])
$$ Le lemme suivant montre que ce morphisme est un modèle de la structure locale étale d'une famille de courbes nodales.

\begin{lemma}
\label{lemma-etale-local-structure-nodal-family}
Soit $f : X \to S$ un morphisme de schémas. Supposons que $f$ soit à singularités au plus nodales de dimension relative $1$. Soit $x \in X$ un point singulier de la fibre $X_s$. Il existe alors un diagramme commutatif de schémas $$
\xymatrix{
X \ar[d] &
U \ar[rr] \ar[l] \ar[rd] & &
W \ar[r] \ar[ld] &
\Spec(\mathbf{Z}[u, v, a]/(uv - a)) \ar[d] \\
S & &
V \ar[ll] \ar[rr] & & \Spec(\mathbf{Z}[a])
}
$$ dans lequel $X \leftarrow U$, $S \leftarrow V$ et $U \to W$ sont des morphismes étales, et où le carré de droite est cartésien, ainsi qu'un point $u \in U$ s'envoyant sur $x$ dans $X$.
\end{lemma}

\begingroup\emergencystretch=3em
\begin{proof}[Première démonstration par approximation d'Artin] Utilisons d'abord l'approximation noethérienne absolue pour nous ramener au cas de schémas de type fini sur $\mathbf{Z}$. La question est locale sur $X$ et $S$. Nous pouvons donc supposer que $X$ et $S$ sont affines. Écrivons alors $S = \Spec(R)$ et écrivons $R$ comme une colimite filtrante $R = \colim R_i$ de $\mathbf{Z}$-algèbres de type fini. D'après Limites, lemme \ref{limits-lemma-descend-finite-presentation}, nous pouvons trouver $i$ et un morphisme $f_i : X_i \to \Spec(R_i)$ dont le changement de base à $S$ est $f$. Après avoir augmenté $i$, nous pouvons supposer que $f_i$ est à singularités au plus nodales de dimension relative $1$; voir le lemme \ref{lemma-descend-nodal-family}. L'image $x_i \in X_i$ de $x$ sera un point singulier de sa fibre, par exemple parce que la formation de $\text{Sing}(f)$ commute au changement de base (Diviseurs, lemme \ref{divisors-lemma-base-change-and-fitting-ideal-omega}). Si nous pouvons démontrer le lemme pour $f_i : X_i \to S_i$ et $x_i$, il en résulte pour $f : X \to S$ par changement de base. Nous sommes ainsi ramenés au cas étudié dans le paragraphe suivant.

\medskip\noindent
Supposons $S$ de type fini sur $\mathbf{Z}$. Soit $s \in S$ l'image de $x$. Rappelons que $\kappa(x)$ est une extension finie séparable de $\kappa(s)$, par exemple parce que $\text{Sing}(f) \to S$ est non ramifié, ou parce que $x$ est un nœud de la fibre $X_s$ et que nous pouvons appliquer le lemme \ref{lemma-nodal}. Soit en outre $\kappa'/\kappa(x)$ l'algèbre séparable de degré $2$ associée à $\mathcal{O}_{X_s, x}$ dans la remarque \ref{remark-quadratic-extension}. D'après Compléments sur les morphismes, lemme \ref{more-morphisms-lemma-realize-prescribed-residue-field-extension-etale}, nous pouvons choisir un voisinage étale $(V, v) \to (S, s)$ tel que l'extension $\kappa(v)/\kappa(s)$ réalise soit l'extension $\kappa(x)/\kappa(s)$ dans le cas $\kappa' \cong \kappa(x) \times \kappa(x)$, soit l'extension $\kappa'/\kappa(s)$ si $\kappa'$ est un corps. Après avoir remplacé $X$ par $X \times_S V$ et $S$ par $V$, nous nous ramenons à la situation décrite au paragraphe suivant.

\medskip\noindent
Supposons $S$ de type fini sur $\mathbf{Z}$ et $x \in X_s$ un nœud déployé; voir la définition \ref{definition-split-node}. D'après le lemme \ref{lemma-formal-local-structure-nodal-family}, il existe un isomorphisme de $\mathcal{O}_{S, s}$-algèbres $$
\mathcal{O}_{X, x}^\wedge \cong
\mathcal{O}_{S, s}^\wedge[[s, t]]/(st - h)
$$ pour un certain $h \in \mathfrak m_s^\wedge \subset \mathcal{O}_{S, s}^\wedge$. Autrement dit, si l'on considère l'homomorphisme $$
\sigma : \mathbf{Z}[a] \longrightarrow \mathcal{O}_{S, s}^\wedge
$$ envoyant $a$ sur $h$, il existe un isomorphisme de $\mathcal{O}_{S, s}$-algèbres $$
\mathcal{O}_{X, x}^\wedge
\longrightarrow
\mathcal{O}_{Y_\sigma, y_\sigma}^\wedge
$$, où $$
Y_\sigma = \Spec(\mathbf{Z}[u, v, t]/(uv - a))
\times_{\Spec(\mathbf{Z}[a]), \sigma} \Spec(\mathcal{O}_{S, s}^\wedge)
$$ et où $y_\sigma$ est le point de $Y_\sigma$ situé au-dessus du point fermé de $\Spec(\mathcal{O}_{S, s}^\wedge)$ dont les coordonnées $u, v$ sont nulles. Puisque $\mathcal{O}_{S, s}$ est un G-anneau d'après Compléments d'algèbre, proposition \ref{more-algebra-proposition-ubiquity-G-ring}, nous pouvons conclure par Compléments sur les morphismes, lemme \ref{more-morphisms-lemma-relative-map-approximation-pre}.
\end{proof}
\endgroup

\begingroup\emergencystretch=3em
\begin{proof}[Démonstration sans approximation d'Artin] Nous ne faisons qu'esquisser cette démonstration, due à Mohan Swaminathan et calquée sur un argument de \cite[Proposition X.2.1]{ACG}. Contrairement à la précédente, cette démonstration utilise les équations elles-mêmes.

\medskip\noindent
Exactement comme dans la première démonstration, nous nous ramenons au cas où $S$ est de type fini sur $\mathbf{Z}$ et où $x \in X_s$ est un nœud déployé; voir la définition \ref{definition-split-node}. D'après le lemme \ref{lemma-split-node}, il existe un isomorphisme de $\kappa(s)$-algèbres $$
\mathcal{O}_{X_s, x}^\wedge \cong \kappa(s)[[u, v]]/(uv)
$$ Remarquons que l'espace tangent de $X_s$ en $x$ est de dimension $2$. D'après Variétés, lemme \ref{varieties-lemma-immersion-into-affine}, après avoir rétréci au sens de Zariski, nous pouvons donc supposer qu'il existe une immersion fermée $X \to Y$ de schémas sur $S$, avec $Y \to S$ lisse de dimension relative $2$. Puisque $X$ est syntomique de dimension relative $1$ (par définition), $X$ est un diviseur de Cartier effectif de $Y$ (cela résulte de Diviseurs, lemme \ref{divisors-lemma-immersion-lci-into-smooth-regular-immersion}). Après avoir rétréci $Y$ et $X$, nous pouvons donc supposer que $Y = \Spec(B)$ est affine et qu'il existe un non-diviseur de zéro $g \in B$ tel que $X = \Spec(B/gB)$.

\medskip\noindent
Écrivons $S = \Spec(R)$ et notons $\mathfrak p \subset R$ l'idéal premier correspondant à $s$. Soit $\mathfrak q \subset B$ l'idéal premier correspondant à $x$, considéré comme un point de $Y$. Nous avons des applications $$
B_\mathfrak q / \mathfrak p B_\mathfrak q
\longrightarrow
B_\mathfrak q / g B_\mathfrak q + \mathfrak p B_\mathfrak q
\longrightarrow
\kappa(s)[[u, v]]/(uv)
$$ dont la seconde devient un isomorphisme après complétion. Après avoir remplacé $B$ par une localisation principale convenable, nous pouvons supposer qu'il existe $U, V \in B$ qui s'envoient sur $u$ et $v$ à $(u, v)^2$ près. Alors le complété de $B_\mathfrak q/\mathfrak p B_\mathfrak q$ est $\kappa(s)[[U, V]]$. De plus, après avoir rétréci $Y$ et multiplié $g$ par une unité, nous pouvons supposer que $g$ s'envoie sur $UV$ plus des termes d'ordre supérieur dans $\kappa(s)[[U, V]]$. Après avoir encore rétréci $Y$, nous pouvons supposer en outre que \begin{enumerate} \item $\mathfrak q = \mathfrak pB + (U, V)$ \item $\Omega_{B/R}$ est libre sur $\text{d}U$ et $\text{d}V$. \end{enumerate} Écrivons $\text{d}(g) = g_2 \text{d}U + g_1 \text{d}V$ avec $g_1, g_2 \in B$. Remarquons que l'idéal $J = (g_1, g_2) \subset B$ est indépendant du choix des « coordonnées » $U$ et $V$, pour $g$ fixé. En réduisant modulo $\mathfrak p B_\mathfrak q + (U, V)^3B_\mathfrak q$, on voit que $g_1$ et $g_2$ sont congrus à $U$ et $V$ modulo $\mathfrak p B_\mathfrak q + (U, V)^2B_\mathfrak q$. Nous pouvons donc reprendre le raisonnement ci-dessus en remplaçant $U$ par $g_1$ et $V$ par $g_2$. Nous obtenons alors $J = (U, V)$, ce qui implique $g_1, g_2 \in (U, V) = J$ (bien entendu, ces éléments diffèrent des $g_1$ et $g_2$ correspondant à notre choix précédent de coordonnées).

\medskip\noindent
Remarquons que $\Spec(B/J) \to S = \Spec(R)$ est étale en $x$. Après avoir remplacé $S$ par un voisinage étale et rétréci encore $Y$, nous pouvons donc supposer que $R \to B \to B/J$ est un isomorphisme. Notons $r \in R$ l'élément dont l'image dans $B/J$ est celle de $-g$. Nous pouvons alors écrire $$
g + r \equiv a U + b V \bmod (U, V)^2
$$ pour un certain $a, b \in R$. En dérivant et en utilisant $\text{d}g \in J\Omega_{B/R}$, nous obtenons $a = b = 0$. Nous pouvons donc écrire $$
g + r = \alpha U^2 + \beta UV + \gamma V^2
$$ pour un certain $\alpha, \beta, \gamma \in B$. Alors $\beta$ s'envoie sur $1$ dans $\kappa(x)$, tandis que $\alpha$ et $\gamma$ s'envoient sur $0$ dans $\kappa(x)$. Nous pouvons donc supposer que $\beta$ et $\alpha + \beta + \gamma$ sont inversibles dans $B$. Considérons le revêtement $Y'$ de $Y$ donné par $$
Y' = \Spec(B[t]/(t(1 - t) - \alpha \gamma / \beta^2))
$$, et soit $x'$ l'unique point au-dessus de $x$ tel que $t = 1$. Remarquons que $Y' \to Y$ est étale en $x'$. Considérons les éléments $$
T_1 = (\alpha + t \beta)U + ((1 - t)\beta + \gamma)V
\quad\text{et}\quad
T_2 =
\frac{(\alpha + (1 - t) \beta)U + (t\beta + \gamma)V}{\alpha + \beta + \gamma}
$$ de $\Gamma(Y', \mathcal{O}_{Y'})$. On voit alors que $$
g = T_1 T_2 - r
$$ Puisque $T_1$ et $T_2$ sont également des coordonnées au sens précédent, le morphisme $Y' \to \mathbf{A}^2_S$ donné par $T_1$ et $T_2$ est étale dans un voisinage de $x'$, et l'image réciproque $X' \subset Y'$ de $X$ est l'image réciproque du sous-schéma fermé de $\mathbf{A}^2_S = \Spec(R[x_1, x_2])$ défini par l'équation $x_1x_2 - r = 0$, quitte à rétrécir encore $Y'$. Nous laissons au lecteur le soin d'assembler ces éléments pour obtenir le diagramme voulu.
\end{proof}
\endgroup




\section{Autres résultats d'annulation}
\label{section-more-vanishing}

\noindent
Suite de la section \ref{section-vanishing}.

\begin{lemma}
\label{lemma-h1-nonzero-degree-leq-2g-2}
Dans la situation \ref{situation-Cohen-Macaulay-curve}, supposons $X$ intègre et de genre $g$. Soit $\mathcal{L}$ un $\mathcal{O}_X$-module inversible. Soit $Z \subset X$ un sous-schéma fermé de dimension $0$, de faisceau d'idéaux $\mathcal{I} \subset \mathcal{O}_X$. Si $H^1(X, \mathcal{I}\mathcal{L})$ est non nul, alors $$
\deg(\mathcal{L}) \leq 2g - 2 + \deg(Z)
$$, avec inégalité stricte sauf si $\mathcal{I}\mathcal{L} \cong \omega_X$.
\end{lemma}

\begingroup\emergencystretch=4em
\begin{proof}
Toute courbe, par exemple $X$, est de Cohen--Macaulay. Si $H^1(X, \mathcal{I}\mathcal{L})$ est non nul, il existe une application non nulle $\mathcal{I}\mathcal{L} \to \omega_X$; voir le lemme \ref{lemma-duality-dim-1-CM}. Puisque $\mathcal{I}\mathcal{L}$ est sans torsion, cette application est injective. Comme un corps est de Gorenstein et que $X$ est réduit, le lieu de Gorenstein $U \subset X$ de $X$ est non vide; voir Dualité pour les schémas, lemme \ref{duality-lemma-gorenstein}. Ce lemme nous dit aussi que $\omega_X|_U$ est inversible. Nous obtenons ainsi une suite exacte courte $$
0 \to \mathcal{I}\mathcal{L} \to \omega_X \to \mathcal{Q} \to 0
$$ où le support de $\mathcal{Q}$ est de dimension zéro. Par conséquent, \begin{align*}
0 & \leq \dim \Gamma(X, \mathcal{Q})\\
& =
\chi(\mathcal{Q}) \\
& =
\chi(\omega_X) - \chi(\mathcal{I}\mathcal{L}) \\
& =
\chi(\omega_X) - \deg(\mathcal{L}) - \chi(\mathcal{I}) \\
& =
2g - 2 - \deg(\mathcal{L}) + \deg(Z)
\end{align*} d'après les lemmes \ref{lemma-euler} et \ref{lemma-rr}, d'après (\ref{equation-genus}), et d'après Variétés, lemmes \ref{varieties-lemma-chi-tensor-finite} et \ref{varieties-lemma-degree-on-proper-curve}. Nous avons aussi utilisé $\deg(Z) = \dim_k \Gamma(Z, \mathcal{O}_Z) = \chi(\mathcal{O}_Z)$ et la suite exacte courte $0 \to \mathcal{I} \to \mathcal{O}_X \to \mathcal{O}_Z \to 0$. Le lemme en résulte.
\end{proof}
\endgroup

\begin{lemma}
\label{lemma-degree-more-than-2g-2}
\begin{reference}
\cite[Lemme 2]{Jongmin}
\end{reference}
Dans la situation \ref{situation-Cohen-Macaulay-curve}, supposons $X$ intègre et de genre $g$. Soit $\mathcal{L}$ un $\mathcal{O}_X$-module inversible. Soit $Z \subset X$ un sous-schéma fermé de dimension $0$, de faisceau d'idéaux $\mathcal{I} \subset \mathcal{O}_X$. Si $\deg(\mathcal{L}) > 2g - 2 + \deg(Z)$, alors $H^1(X, \mathcal{I}\mathcal{L}) = 0$ et l'une des possibilités suivantes se présente :
\begin{enumerate}
\item $H^0(X, \mathcal{I}\mathcal{L}) \not = 0$;
\item $g = 0$ et $\deg(\mathcal{L}) = \deg(Z) - 1$.
\end{enumerate}
Dans le cas (2), si $Z = \emptyset$, alors $X \cong \mathbf{P}^1_k$ et $\mathcal{L}$ correspond à $\mathcal{O}_{\mathbf{P}^1}(-1)$.
\end{lemma}

\begingroup\emergencystretch=3em
\begin{proof}
L'annulation de $H^1(X, \mathcal{I}\mathcal{L})$ résulte du lemme \ref{lemma-h1-nonzero-degree-leq-2g-2}. Si $H^0(X, \mathcal{I}\mathcal{L}) = 0$, alors $\chi(\mathcal{I}\mathcal{L}) = 0$. La suite exacte courte $0 \to \mathcal{I}\mathcal{L} \to \mathcal{L} \to \mathcal{O}_Z \to 0$ donne $\deg(\mathcal{L}) = g - 1 + \deg(Z)$. Ainsi $g - 1 + \deg(Z) > 2g - 2 + \deg(Z)$, ce qui implique $g = 0$; le cas (2) est donc réalisé. Si $Z = \emptyset$ dans le cas (2), alors $\mathcal{L}^{-1}$ est un faisceau inversible de degré $1$. Il existe par conséquent un isomorphisme $X \to \mathbf{P}^1_k$, et $\mathcal{L}^{-1}$ est l'image réciproque de $\mathcal{O}_{\mathbf{P}^1}(1)$ d'après le lemme \ref{lemma-genus-zero-positive-degree}.
\end{proof}
\endgroup

\begin{lemma}
\label{lemma-degree-more-than-2g}
\begin{reference}
\cite[Lemme 3]{Jongmin}
\end{reference}
Dans la situation \ref{situation-Cohen-Macaulay-curve}, supposons $X$ intègre et de genre $g$. Soit $\mathcal{L}$ un $\mathcal{O}_X$-module inversible. Si $\deg(\mathcal{L}) \geq 2g$, alors $\mathcal{L}$ est engendré par ses sections globales.
\end{lemma}

\begin{proof}
Soit $Z \subset X$ le sous-schéma fermé découpé par les sections globales de $\mathcal{L}$. D'après le lemme \ref{lemma-degree-more-than-2g-2}, on voit que $Z \not = X$. Soit $\mathcal{I} \subset \mathcal{O}_X$ le faisceau d'idéaux définissant $Z$. Considérons la suite exacte courte $$
0 \to \mathcal{I}\mathcal{L}
\to \mathcal{L} \to \mathcal{O}_Z \to 0
$$ Si $Z \not = \emptyset$, alors $H^1(X, \mathcal{I}\mathcal{L})$ est non nul, comme le montre la suite exacte longue de cohomologie. D'après le lemme \ref{lemma-duality-dim-1-CM}, cela donne une application non nulle, donc injective, $$
\mathcal{I}\mathcal{L}
\longrightarrow
\omega_X
$$ En particulier, nous obtenons une application injective $H^0(X, \mathcal{L}) = H^0(X, \mathcal{I}\mathcal{L})
\to H^0(X, \omega_X)$. C'est impossible, puisque $$
\dim_k H^0(X, \mathcal{L}) = \dim_k H^1(X, \mathcal{L}) +
\deg(\mathcal{L}) + 1 - g \geq g + 1
$$ et $\dim H^0(X, \omega_X) = g$ d'après (\ref{equation-genus}).
\end{proof}

\begin{lemma}
\label{lemma-degree-more-than-2g-1-and-Z}
Dans la situation \ref{situation-Cohen-Macaulay-curve}, supposons $X$ intègre et de genre $g$. Soit $\mathcal{L}$ un $\mathcal{O}_X$-module inversible. Soit $Z \subset X$ un sous-schéma fermé non vide de dimension $0$. Si $\deg(\mathcal{L}) \geq 2g - 1 + \deg(Z)$, alors $\mathcal{L}$ est engendré par ses sections globales et $H^0(X, \mathcal{L}) \to H^0(X, \mathcal{L}|_Z)$ est surjective.
\end{lemma}

\begin{proof}
L'engendrement par les sections globales résulte du lemme \ref{lemma-degree-more-than-2g}. Si $\mathcal{I} \subset \mathcal{O}_X$ est le faisceau d'idéaux de $Z$, alors $H^1(X, \mathcal{I}\mathcal{L}) = 0$ d'après le lemme \ref{lemma-h1-nonzero-degree-leq-2g-2}. D'où la surjectivité.
\end{proof}

\begin{lemma}
\label{lemma-degree-at-least-2g+1}
Dans la situation \ref{situation-Cohen-Macaulay-curve}, supposons $X$ géométriquement intègre sur $k$ et de genre $g$. Soit $\mathcal{L}$ un $\mathcal{O}_X$-module inversible. Si $\deg(\mathcal{L}) \geq 2g + 1$, alors $\mathcal{L}$ est très ample.
\end{lemma}

\begin{proof}
D'après le lemme \ref{lemma-degree-more-than-2g}, $\mathcal{L}$ est engendré par ses sections globales et détermine donc un morphisme $f : X \to \mathbf{P}^n_k$, où $n = h^0(X,\mathcal{L}) - 1$. Dire que $\mathcal{L}$ est très ample revient à montrer que $f$ est une immersion fermée. Il suffit de vérifier que le changement de base de $f$ à une clôture algébrique $\overline{k}$ de $k$ est une immersion fermée (Descente, lemme \ref{descent-lemma-descending-property-closed-immersion}). Nous pouvons donc supposer $k$ algébriquement clos; par hypothèse, $X$ reste intègre. Le lemme \ref{lemma-degree-more-than-2g-1-and-Z} montre que, pour tout sous-schéma fermé de dimension $0$, noté $Z\subset X$, et de degré 2, l'application de restriction $H^0(X, \mathcal{L}) \to H^0(X, \mathcal{L}|_Z)$ est surjective. D'après Variétés, lemme \ref{varieties-lemma-variant-separate-points-tangent-vectors}, $\mathcal{L}$ est très ample.
\end{proof}


\begin{lemma}
\label{lemma-vanishing-on-gorenstein}
\begin{reference}
Version faible de \cite[Lemme 4]{Jongmin}
\end{reference}
Soit $k$ un corps. Soit $X$ un schéma propre sur $k$, réduit, connexe et de dimension $1$. Soit $\mathcal{L}$ un $\mathcal{O}_X$-module inversible. Soit $Z \subset X$ un sous-schéma fermé de dimension $0$, de faisceau d'idéaux $\mathcal{I} \subset \mathcal{O}_X$. Si $H^1(X, \mathcal{I}\mathcal{L}) \not = 0$, il existe un sous-schéma fermé réduit connexe $Y \subset X$ de dimension $1$ tel que
$$
\deg(\mathcal{L}|_Y) \leq -2\chi(Y, \mathcal{O}_Y) + \deg(Z \cap Y)
$$, où $Z \cap Y$ est l'intersection schématique.
\end{lemma}

\begin{proof}
Si $H^1(X, \mathcal{I}\mathcal{L})$ est non nul, il existe une application non nulle $\varphi : \mathcal{I}\mathcal{L} \to \omega_X$; voir le lemme \ref{lemma-duality-dim-1-CM}. Soit $Y \subset X$ la réunion des composantes irréductibles $C$ de $X$ telles que $\varphi$ soit non nul au point générique de $C$. Alors $Y$ est un sous-schéma fermé réduit. Soit $\mathcal{J} \subset \mathcal{O}_X$ le faisceau d'idéaux de $Y$. Puisque $\mathcal{J}\mathcal{I}\mathcal{L}$ n'a pas de point associé immergé, en tant que sous-module de $\mathcal{L}$, et que $\varphi$ est nul aux points génériques du support de $\mathcal{J}$, par le choix de $Y$ et puisque $X$ est réduit, on voit que $\varphi$ se factorise sous la forme $$
\mathcal{I}\mathcal{L} \to
\mathcal{I}\mathcal{L}/\mathcal{J}\mathcal{I}\mathcal{L} \to \omega_X
$$ Nous pouvons considérer $\mathcal{I}\mathcal{L}/\mathcal{J}\mathcal{I}\mathcal{L}$ comme l'image directe d'un faisceau cohérent sur $Y$, que nous noterons encore de la même manière par abus de notation. Puisque $\omega_Y = \SheafHom(\mathcal{O}_Y, \omega_X)$ d'après le lemme \ref{lemma-closed-immersion-dim-1-CM}, nous obtenons une application $$
\mathcal{I}\mathcal{L}/
\mathcal{J}\mathcal{I}\mathcal{L} 
\to \omega_Y
$$ de $\mathcal{O}_Y$-modules, injective aux points génériques de $Y$. Soit $\mathcal{I}' \subset \mathcal{O}_Y$ le faisceau d'idéaux de $Z \cap Y$. Il existe une application $\mathcal{I}\mathcal{L}/\mathcal{J}\mathcal{I}\mathcal{L} \to
\mathcal{I}'\mathcal{L}|_Y$ dont le noyau est supporté en des points fermés. Comme $\omega_Y$ est un module de Cohen--Macaulay, l'application précédente se factorise par une application injective $\mathcal{I}'\mathcal{L}|_Y \to
\omega_Y$. Nous obtenons ainsi une suite exacte $$
0 \to \mathcal{I}'\mathcal{L}|_Y \to \omega_Y \to \mathcal{Q} \to 0
$$ de faisceaux cohérents sur $Y$, où $\mathcal{Q}$ est supporté en dimension $0$; on utilise ici le fait que $\omega_Y$ est un module inversible aux points génériques de $Y$. Nous concluons que $$
0 \leq \dim \Gamma(Y, \mathcal{Q}) =
\chi(\mathcal{Q}) = \chi(\omega_Y) - \chi(\mathcal{I}'\mathcal{L}) =
-2\chi(\mathcal{O}_Y) - \deg(\mathcal{L}|_Y) + \deg(Z \cap Y)
$$ d'après le lemme \ref{lemma-euler} et Variétés, lemme \ref{varieties-lemma-chi-tensor-finite}. Si $Y$ est connexe, cela démontre le lemme. Sinon, nous répétons la dernière partie de l'argument pour l'une des composantes connexes de $Y$.
\end{proof}

\begin{lemma}
\label{lemma-global-generation}
Soit $k$ un corps. Soit $X$ un schéma propre sur $k$, réduit, connexe et de dimension $1$. Soit $\mathcal{L}$ un $\mathcal{O}_X$-module inversible. Supposons que, pour tout sous-schéma fermé réduit connexe $Y \subset X$ de dimension $1$, on ait $$
\deg(\mathcal{L}|_Y) \geq 2\dim_k H^1(Y, \mathcal{O}_Y)
$$ Alors $\mathcal{L}$ est engendré par ses sections globales.
\end{lemma}

\begin{proof}
Raisonnons par récurrence sur le nombre de composantes irréductibles de $X$. Si $X$ est irréductible, le lemme résulte du lemme \ref{lemma-degree-more-than-2g} appliqué à $X$ considéré comme un schéma sur le corps $k' = H^0(X, \mathcal{O}_X)$. Supposons $X$ non irréductible. Avant de poursuivre, si $k$ est fini, remplaçons $k$ par une extension purement transcendante $K$. Cela est permis par Variétés, lemmes \ref{varieties-lemma-globally-generated-base-change}, \ref{varieties-lemma-degree-base-change}, \ref{varieties-lemma-geometrically-reduced-any-base-change} et \ref{varieties-lemma-bijection-irreducible-components}, par Cohomologie des schémas, lemme \ref{coherent-lemma-flat-base-change-cohomology}, par le lemme \ref{lemma-sanity-check-duality} et par le fait élémentaire que $K$ est géométriquement intègre sur $k$.

\medskip\noindent
Supposons, pour aboutir à une contradiction, que $\mathcal{L}$ ne soit pas engendré par ses sections globales. Nous pouvons alors choisir un faisceau cohérent d'idéaux $\mathcal{I} \subset \mathcal{O}_X$ tel que $H^0(X, \mathcal{I}\mathcal{L}) = H^0(X, \mathcal{L})$ et tel que $\mathcal{O}_X/\mathcal{I}$ soit non nul et de support de dimension $0$. On peut, par exemple, prendre pour $\mathcal{I}$ le faisceau d'idéaux d'un point fermé quelconque du lieu d'annulation commun des sections globales de $\mathcal{L}$. Considérons la suite exacte courte $$
0 \to \mathcal{I}\mathcal{L} \to \mathcal{L} \to 
\mathcal{L}/\mathcal{I}\mathcal{L} \to 0
$$ Puisque le support de $\mathcal{L}/\mathcal{I}\mathcal{L}$ est de dimension $0$, $\mathcal{L}/\mathcal{I}\mathcal{L}$ est engendré par ses sections globales (Variétés, lemme \ref{varieties-lemma-chi-tensor-finite}). La suite exacte courte et le fait que $H^0(X, \mathcal{I}\mathcal{L}) = H^0(X, \mathcal{L})$ donnent une injection $H^0(X, \mathcal{L}/\mathcal{I}\mathcal{L}) \to H^1(X, \mathcal{I}\mathcal{L})$.

\medskip\noindent
Rappelons que le $k$-espace vectoriel $H^1(X, \mathcal{I}\mathcal{L})$ est dual de $\Hom(\mathcal{I}\mathcal{L}, \omega_X)$. Choisissons $\varphi : \mathcal{I}\mathcal{L} \to \omega_X$. D'après le lemme \ref{lemma-vanishing-on-gorenstein}, on a $H^1(X, \mathcal{L}) = 0$. Par conséquent, \begingroup\footnotesize $$
\dim_k H^0(X, \mathcal{I}\mathcal{L}) = \dim_k H^0(X, \mathcal{L}) =
\deg(\mathcal{L}) + \chi(\mathcal{O}_X) > \dim_k H^1(X, \mathcal{O}_X) =
\dim_k H^0(X, \omega_X)
$$\endgroup Nous en concluons que $\varphi$ n'est pas injective sur les sections globales; en particulier, $\varphi$ n'est pas injective. Pour chaque point générique $\eta \in X$ d'une composante irréductible de $X$, notons $V_\eta \subset \Hom(\mathcal{I}\mathcal{L}, \omega_X)$ le sous-$k$-espace vectoriel formé des $\varphi$ qui sont nuls en $\eta$. Puisque tout point associé de $\mathcal{I}\mathcal{L}$ est un point générique de $X$, ce qui précède montre que $\Hom(\mathcal{I}\mathcal{L}, \omega_X) = \bigcup V_\eta$. Comme $X$ possède un nombre fini de points génériques et que $k$ est infini, nous concluons $\Hom(\mathcal{I}\mathcal{L}, \omega_X) = V_\eta$ pour un certain $\eta$. Soit $\eta \in C \subset X$ la composante irréductible correspondante. Soit $Y \subset X$ la réunion des autres composantes irréductibles de $X$. Alors $Y$ est un sous-schéma fermé réduit non vide, distinct de $X$. Soit $\mathcal{J} \subset \mathcal{O}_X$ le faisceau d'idéaux de $Y$. Gardons à l'esprit que le support de $\mathcal{J}$ est $C$.

\medskip\noindent
Soit $\varphi : \mathcal{I}\mathcal{L} \to \omega_X$ quelconque. Puisque $\mathcal{J}\mathcal{I}\mathcal{L}$ n'a pas de point associé immergé, en tant que sous-module de $\mathcal{L}$, et que $\varphi$ est nul au point générique $\eta$ du support de $\mathcal{J}$, on voit que $\varphi$ se factorise sous la forme $$
\mathcal{I}\mathcal{L} \to
\mathcal{I}\mathcal{L}/\mathcal{J}\mathcal{I}\mathcal{L} \to \omega_X
$$ Nous pouvons considérer $\mathcal{I}\mathcal{L}/\mathcal{J}\mathcal{I}\mathcal{L}$ comme l'image directe d'un faisceau cohérent sur $Y$, que nous noterons encore de la même manière par abus de notation. Puisque $\omega_Y = \SheafHom(\mathcal{O}_Y, \omega_X)$ d'après le lemme \ref{lemma-closed-immersion-dim-1-CM}, nous obtenons une factorisation $$
\mathcal{I}\mathcal{L} \to
\mathcal{I}\mathcal{L}/
\mathcal{J}\mathcal{I}\mathcal{L} 
\xrightarrow{\varphi'} \omega_Y \to \omega_X
$$ de $\varphi$. Soit $\mathcal{I}' \subset \mathcal{O}_Y$ l'image de $\mathcal{I} \subset \mathcal{O}_X$. Il existe une application surjective $\mathcal{I}\mathcal{L}/\mathcal{J}\mathcal{I}\mathcal{L} \to
\mathcal{I}'\mathcal{L}|_Y$ dont le noyau est supporté en des points fermés. Comme $\omega_Y$ est un module de Cohen--Macaulay sur $Y$, l'application $\varphi'$ se factorise par une application $\varphi'' : \mathcal{I}'\mathcal{L}|_Y \to \omega_Y$. Nous avons donc des diagrammes commutatifs $$
\vcenter{
\xymatrix{
0 \ar[r] &
\mathcal{I}\mathcal{L} \ar[r] \ar[d] &
\mathcal{L} \ar[r] \ar[d] &
\mathcal{L}/\mathcal{I}\mathcal{L} \ar[r] \ar[d] &
0 \\
0 \ar[r] &
\mathcal{I}'\mathcal{L}|_Y \ar[r] &
\mathcal{L}|_Y \ar[r] &
\mathcal{L}|_Y/\mathcal{I}'\mathcal{L}|_Y \ar[r] &
0
}
}
\quad\text{et}\quad
\vcenter{
\xymatrix{
\mathcal{I}\mathcal{L} \ar[r]_\varphi \ar[d] & \omega_X \\
\mathcal{I}'\mathcal{L}|_Y \ar[r]^{\varphi''} & \omega_Y \ar[u]
}
}
$$ Nous pouvons maintenant achever la démonstration comme suit. Puisque, pour tout $\varphi$, nous avons $\varphi''$ et puisque $\omega_X \in \textit{Coh}(\mathcal{O}_X)$ représente le foncteur $\mathcal{F} \mapsto \Hom_k(H^1(X, \mathcal{F}), k)$, on voit que $H^1(X, \mathcal{I}\mathcal{L}) \to H^1(Y, \mathcal{I}'\mathcal{L}|_Y)$ est injective. Comme l'application de bord $H^0(X, \mathcal{L}/\mathcal{I}\mathcal{L}) \to H^1(X, \mathcal{I}\mathcal{L})$ est injective, le composé $$
H^0(X, \mathcal{L}/\mathcal{I}\mathcal{L}) \to
H^0(X, \mathcal{L}|_Y/\mathcal{I}'\mathcal{L}|_Y) \to
H^1(X, \mathcal{I}'\mathcal{L}|_Y)
$$ l'est aussi. Puisque $\mathcal{L}/\mathcal{I}\mathcal{L} \to
\mathcal{L}|_Y/\mathcal{I}'\mathcal{L}|_Y$ est une application surjective de modules cohérents dont les supports sont de dimension $0$, la première application $H^0(X, \mathcal{L}/\mathcal{I}\mathcal{L}) \to
H^0(X, \mathcal{L}|_Y/\mathcal{I}'\mathcal{L}|_Y)$ est surjective, donc bijective. Mais, par récurrence, $\mathcal{L}|_Y$ est engendré par ses sections globales, y compris si $Y$ est non connexe, et l'application de bord $$
H^0(X, \mathcal{L}|_Y/\mathcal{I}'\mathcal{L}|_Y) \to
H^1(X, \mathcal{I}'\mathcal{L}|_Y)
$$ ne peut donc pas être injective. Cette contradiction achève la démonstration.
\end{proof}





\section{Contraction des queues rationnelles}
\label{section-contracting-rational-tails}

\noindent
Dans cette section, nous étudions le cas le plus simple de contraction d'un schéma visant à améliorer les propriétés de positivité de son faisceau canonique.

\begin{example}[Contraction d'une queue rationnelle] \label{example-rational-tail} Soit $k$ un corps. Soit $X$ un schéma propre sur $k$, de dimension $1$ et tel que $H^0(X, \mathcal{O}_X) = k$. Supposons que les singularités de $X$ soient au plus nodales. Une {\it queue rationnelle} est une composante irréductible $C \subset X$, considérée comme un sous-schéma fermé intègre, qui possède les propriétés suivantes : \begin{enumerate} \item $X' \not = \emptyset$, où $X' \subset X$ est l'adhérence schématique de $X \setminus C$; \item l'intersection schématique $C \cap X'$ est un unique point réduit $x$; \item $H^0(C, \mathcal{O}_C)$ s'envoie isomorphiquement sur le corps résiduel de $x$; \item $C$ est de genre zéro. \end{enumerate} Comme au moins deux composantes irréductibles de $X$ passent par $x$, nous en concluons que $x$ est un nœud. Posons $k' = H^0(C, \mathcal{O}_C) = \kappa(x)$. Alors $k'/k$ est une extension finie séparable de corps (lemme \ref{lemma-nodal}). Il existe un morphisme canonique $$
c : X \longrightarrow X'
$$ qui induit l'identité sur $X'$ et envoie $C$ sur $x \in X'$ par le morphisme canonique $C \to \Spec(k') = x$. Cela résulte de Morphismes, lemme \ref{morphisms-lemma-scheme-theoretic-union}, puisque $X$ est la réunion schématique de $C$ et $X'$, car $X$ est réduit. De plus, nous affirmons que $$
c_*\mathcal{O}_X = \mathcal{O}_{X'}
\quad\text{et}\quad
R^1c_*\mathcal{O}_X = 0
$$ Pour le voir, notons $i_C : C \to X$, $i_{X'} : X' \to X$ et $i_x : x \to X$ les immersions et utilisons la suite exacte $$
0 \to \mathcal{O}_X \to
i_{C, *}\mathcal{O}_C \oplus i_{X', *}\mathcal{O}_{X'} \to
i_{x, *}\kappa(x) \to 0
$$ de Morphismes, lemme \ref{morphisms-lemma-scheme-theoretic-union}. La suite exacte longue des images directes supérieures montre qu'il suffit d'établir $H^0(C, \mathcal{O}_C) = k'$ et $H^1(C, \mathcal{O}_C) = 0$, ce qui résulte des hypothèses. Remarquons que $X'$ est lui aussi un schéma propre sur $k$, de dimension $1$, dont les singularités sont au plus nodales (lemme \ref{lemma-closed-subscheme-nodal-curve}), que $H^0(X', \mathcal{O}_{X'}) = k$, et que $X'$ a le même genre que $X$. On dit que $c : X \to X'$ est la {\it contraction d'une queue rationnelle}.
\end{example}

\begin{lemma}
\label{lemma-rational-tail-negative}
Soit $k$ un corps. Soit $X$ un schéma propre sur $k$, de dimension $1$ et tel que $H^0(X, \mathcal{O}_X) = k$. Supposons que les singularités de $X$ soient au plus nodales. Soit $C \subset X$ une queue rationnelle (exemple \ref{example-rational-tail}). Alors $\deg(\omega_X|_C) < 0$.
\end{lemma}

\begin{proof}
Soit $X' \subset X$ comme dans l'exemple. Nous avons une suite exacte courte $$
0 \to \omega_C \to \omega_X|_C \to \mathcal{O}_{C \cap X'} \to 0
$$ Voir les lemmes \ref{lemma-closed-subscheme-reduced-gorenstein}, \ref{lemma-facts-about-nodal-curves} et \ref{lemma-closed-subscheme-nodal-curve}. Avec $k'$ comme dans l'exemple, on voit que $\deg(\omega_C) = -2[k' : k]$ puisque $C \cong \mathbf{P}^1_{k'}$, d'après la proposition \ref{proposition-projective-line}, et que $\deg(C \cap X') = [k' : k]$. Ainsi $\deg(\omega_X|_C) = -[k' : k]$, qui est négatif.
\end{proof}

\begin{lemma}
\label{lemma-rational-tail-field-extension}
Soit $k$ un corps. Soit $X$ un schéma propre sur $k$, de dimension $1$ et tel que $H^0(X, \mathcal{O}_X) = k$. Supposons que les singularités de $X$ soient au plus nodales. Soit $C \subset X$ une queue rationnelle (exemple \ref{example-rational-tail}). Pour toute extension de corps $K/k$, le changement de base $C_K \subset X_K$ est une réunion disjointe finie de queues rationnelles.
\end{lemma}

\begin{proof}
Soient $x \in C$ et $k' = \kappa(x)$ comme dans l'exemple. Remarquons que $C \cong \mathbf{P}^1_{k'}$ d'après la proposition \ref{proposition-projective-line}. Puisque $k'/k$ est fini séparable, $k' \otimes_k K = K'_1 \times \ldots \times K'_n$ est un produit fini d'extensions finies séparables $K'_i/K$. Posons $C_i = \mathbf{P}^1_{K'_i}$ et notons $x_i \in C_i$ l'image réciproque de $x$. Alors $C_K = \coprod C_i$ et $X'_K \cap C_i = x_i$, comme voulu.
\end{proof}

\begin{lemma}
\label{lemma-no-rational-tail}
Soit $k$ un corps. Soit $X$ un schéma propre sur $k$, de dimension $1$ et tel que $H^0(X, \mathcal{O}_X) = k$. Supposons que les singularités de $X$ soient au plus nodales. Si $X$ ne possède pas de queue rationnelle (exemple \ref{example-rational-tail}), alors, pour tout sous-schéma fermé réduit connexe $Y \subset X$, $Y \not = X$, de dimension $1$, on a $\deg(\omega_X|_Y) \geq \dim_k H^1(Y, \mathcal{O}_Y)$.
\end{lemma}

\begin{proof}
Soit $Y \subset X$ comme dans l'énoncé. Alors $k' = H^0(Y, \mathcal{O}_Y)$ est un corps et une extension finie de $k$, et $[k' : k]$ divise tous les invariants numériques ci-dessous associés à $Y$ et aux faisceaux cohérents sur $Y$; voir Variétés, lemme \ref{varieties-lemma-divisible}. Soit $Z \subset X$ comme dans le lemme \ref{lemma-closed-subscheme-reduced-gorenstein}. Nous utiliserons sans autre mention les résultats de ce lemme et des lemmes \ref{lemma-facts-about-nodal-curves} et \ref{lemma-closed-subscheme-nodal-curve}. Nous obtenons une suite exacte courte $$
0 \to \omega_Y \to \omega_X|_Y \to \mathcal{O}_{Y \cap Z} \to 0
$$ Voir le lemme \ref{lemma-closed-subscheme-reduced-gorenstein}. Nous en concluons que $$
\deg(\omega_X|_Y) = \deg(Y \cap Z) + \deg(\omega_Y) =
\deg(Y \cap Z) - 2\chi(Y, \mathcal{O}_Y)
$$ Par conséquent, si le lemme est faux, alors $$
2[k' : k] > \deg(Y \cap Z) + \dim_k H^1(Y, \mathcal{O}_Y)
$$ Puisque $Y \cap Z$ est non vide, et compte tenu de la divisibilité mentionnée ci-dessus, cela n'est possible que si $Y \cap Z$ est un unique point $k'$-rationnel du lieu lisse de $Y$ et si $H^1(Y, \mathcal{O}_Y) = 0$. Si $Y$ est irréductible, cela implique que $Y$ est une queue rationnelle. Si $Y$ est réductible, puisque $\deg(\omega_X|_Y) = -[k' : k]$, il existe une composante irréductible $C$ de $Y$ telle que $\deg(\omega_X|_C) < 0$; voir Variétés, lemme \ref{varieties-lemma-degree-in-terms-of-components}. L'analyse précédente appliquée à $C$ montre alors que $C$ est une queue rationnelle.
\end{proof}

\begin{lemma}
\label{lemma-no-rational-tail-semiample-genus-geq-2}
Soit $k$ un corps. Soit $X$ un schéma propre sur $k$, de dimension $1$ et tel que $H^0(X, \mathcal{O}_X) = k$. Supposons que les singularités de $X$ soient au plus nodales et que $X$ ne possède pas de queue rationnelle (exemple \ref{example-rational-tail}). Si \begin{enumerate} \item le genre de $X$ est $0$, alors $X$ est isomorphe à une conique plane irréductible et $\omega_X^{\otimes -1}$ est très ample; \item le genre de $X$ est $1$, alors $\omega_X \cong \mathcal{O}_X$; \item le genre de $X$ est $\geq 2$, alors $\omega_X^{\otimes m}$ est engendré par ses sections globales pour $m \geq 2$.
\end{enumerate}
\end{lemma}

\begin{proof}
D'après le lemme \ref{lemma-facts-about-nodal-curves}, $X$ est de Gorenstein, c'est-à-dire que $\omega_X$ est un $\mathcal{O}_X$-module inversible.

\medskip\noindent
Si le genre de $X$ est nul, alors $\deg(\omega_X) < 0$. Par conséquent, si $X$ possède plus d'une composante irréductible, nous obtenons une contradiction avec le lemme \ref{lemma-no-rational-tail}. Dans le cas irréductible, $X$ est isomorphe à une conique plane irréductible et $\omega_X^{\otimes -1}$ est très ample d'après le lemme \ref{lemma-genus-zero}.

\medskip\noindent
Si le genre de $X$ est $1$, alors $\omega_X$ possède une section globale et $\deg(\omega_X|_C) = 0$ pour toute composante irréductible. En effet, $\deg(\omega_X|_C) \geq 0$ pour toute composante irréductible $C$ d'après le lemme \ref{lemma-no-rational-tail}, la somme de ces nombres est $0$ d'après le lemme \ref{lemma-genus-gorenstein}, et nous pouvons appliquer Variétés, lemme \ref{varieties-lemma-degree-in-terms-of-components}. Alors $\omega_X \cong \mathcal{O}_X$ d'après Variétés, lemme \ref{varieties-lemma-no-sections-dual-nef}.

\medskip\noindent
Supposons que le genre $g$ de $X$ soit supérieur ou égal à $2$. Si $X$ est irréductible, le résultat découle du lemme \ref{lemma-degree-more-than-2g}. Supposons $X$ réductible. D'après le lemme \ref{lemma-no-rational-tail}, les inégalités du lemme \ref{lemma-global-generation} sont satisfaites pour tout $Y \subset X$ comme dans l'énoncé, sauf pour $Y = X$. En examinant la démonstration du lemme \ref{lemma-global-generation}, on voit que, dans le cas réductible, les seules inégalités utilisées pour $Y = X$ sont $$
\deg(\omega_X^{\otimes m}) > -2 \chi(\mathcal{O}_X)
\quad\text{et}\quad
\deg(\omega_X^{\otimes m}) + \chi(\mathcal{O}_X) > \dim_k H^1(X, \mathcal{O}_X)
$$ Comme elles sont toutes deux satisfaites sous les hypothèses $g \geq 2$ et $m \geq 2$, nous concluons.
\end{proof}

\begin{lemma}
\label{lemma-contracting-rational-tails}
Soit $k$ un corps. Soit $X$ un schéma propre sur $k$, de dimension $1$ et tel que $H^0(X, \mathcal{O}_X) = k$. Supposons que les singularités de $X$ soient au plus nodales. Considérons une suite $$
X = X_0 \to X_1 \to \ldots \to X_n = X'
$$ de contractions de queues rationnelles (exemple \ref{example-rational-tail}), poursuivie jusqu'à ce qu'il n'en reste aucune. Alors \begin{enumerate} \item si le genre de $X$ est $0$, $X'$ est une conique plane irréductible; \item si le genre de $X$ est $1$, alors $\omega_{X'} \cong \mathcal{O}_X$; \item si le genre de $X$ est $> 1$, alors $\omega_{X'}^{\otimes m}$ est engendré par ses sections globales pour $m \geq 2$. \end{enumerate} Si le genre de $X$ est $\geq 1$, le morphisme $X \to X'$ est indépendant des choix et sa formation commute aux extensions du corps de base.
\end{lemma}

\begin{proof}
Contractons successivement les queues rationnelles jusqu'à ce qu'il n'en reste aucune. Les assertions (1), (2) et (3) résultent alors du lemme \ref{lemma-no-rational-tail-semiample-genus-geq-2}.

\medskip\noindent
Unicité. Pour voir que $f : X \to X'$ est indépendant des choix effectués, il suffit de montrer que toute queue rationnelle $C \subset X$ est envoyée sur un point par $X \to X'$; nous omettons quelques détails. Dans le cas contraire, il existe une section $s \in \Gamma(X', \omega_{X'}^{\otimes 2})$ qui ne s'annule pas au point générique de la composante irréductible $f(C)$. Puisque, pour chacune des contractions $X_i \to X_{i + 1}$, nous disposons d'une section $X_{i + 1} \to X_i$, il existe une section $X' \to X$ de $f$. Nous avons alors une suite exacte $$
0 \to \omega_{X'} \to \omega_X \to \omega_X|_{X''} \to 0
$$, où $X'' \subset X$ est la réunion des composantes irréductibles contractées par $f$. Voir le lemme \ref{lemma-closed-subscheme-reduced-gorenstein}. Nous obtenons donc une application $\omega_{X'}^{\otimes 2} \to \omega_X^{\otimes 2}$; en prenant l'image de $s$, nous obtenons une section de $\omega_X^{\otimes 2}$ qui ne s'annule pas au point générique de $C$. Cela contredit le fait que la restriction de $\omega_X$ à une queue rationnelle est de degré négatif (lemme \ref{lemma-rational-tail-negative}).

\medskip\noindent
L'assertion concernant les extensions du corps de base résulte du lemme \ref{lemma-rational-tail-field-extension}. Nous omettons quelques détails.
\end{proof}





\section{Contraction des ponts rationnels}
\label{section-contracting-rational-bridges}

\noindent
Dans cette section, nous étudions le deuxième cas le plus simple, après celui de la section \ref{section-contracting-rational-tails}, de contraction d'un schéma visant à améliorer les propriétés de positivité de son faisceau canonique.

\begin{example}[Contraction d'un pont rationnel] \label{example-rational-bridge} Soit $k$ un corps. Soit $X$ un schéma propre sur $k$, de dimension $1$ et tel que $H^0(X, \mathcal{O}_X) = k$. Supposons que les singularités de $X$ soient au plus nodales. Un {\it pont rationnel} est une composante irréductible $C \subset X$, considérée comme un sous-schéma fermé intègre, qui possède les propriétés suivantes : \begin{enumerate} \item $X' \not = \emptyset$, où $X' \subset X$ est l'adhérence schématique de $X \setminus C$; \item l'intersection schématique $C \cap X'$ est de degré $2$ sur $H^0(C, \mathcal{O}_C)$; \item $C$ est de genre zéro. \end{enumerate} Posons $k' = H^0(C, \mathcal{O}_C)$ et $k'' = H^0(C \cap X', \mathcal{O}_{C \cap X'})$. Alors $k'$ est un corps (Variétés, lemme \ref{varieties-lemma-proper-geometrically-reduced-global-sections}) et $\dim_{k'}(k'') = 2$. Comme au moins deux composantes irréductibles de $X$ passent par chaque point de $C \cap X'$, nous concluons que ces points sont des nœuds de $X$ et des points lisses à la fois de $C$ et de $X'$ (lemme \ref{lemma-closed-subscheme-nodal-curve}). Ainsi $k'/k$ est une extension finie séparable de corps, et $k''/k'$ est soit une extension séparable de corps de degré $2$, soit $k'' = k' \times k'$ (lemme \ref{lemma-nodal}). D'après la section \ref{section-pushouts}, il existe un carré cocartésien $$
\xymatrix{
C \cap X' \ar[r] \ar[d] &
X' \ar[d]^a \\
\Spec(k') \ar[r] &
Y
}
$$ doté de nombreuses bonnes propriétés, que nous utiliserons toutes ci-dessous sans autre mention. Soit $y \in Y$ l'image de $\Spec(k') \to Y$. Alors $$
\mathcal{O}_{Y, y}^\wedge \cong k'[[s, t]]/(st)
\quad\text{ou}\quad
\mathcal{O}_{Y, y}^\wedge \cong
\{f \in k''[[s]] : f(0) \in k'\}
$$, selon que $C \cap X'$ possède $2$ ou $1$ points. Cela résulte du lemme \ref{lemma-complete-local-ring-pushout} et du fait que $\mathcal{O}_{X', p} \cong \kappa(p)[[t]]$ pour $p \in C \cap X'$ d'après Compléments d'algèbre, lemme \ref{more-algebra-lemma-power-series-over-residue-field}. Nous voyons ainsi que $y \in Y$ est un nœud; voir les lemmes \ref{lemma-nodal} et \ref{lemma-2-branches-delta-1}, et en particulier la discussion du cas II dans la démonstration de (2) $\Rightarrow$ (1) du lemme \ref{lemma-2-branches-delta-1}. Les singularités de $Y$ sont donc au plus nodales.

\medskip\noindent
Nous pouvons prolonger le diagramme commutatif ci-dessus en un diagramme $$
\xymatrix{
C \cap X' \ar[r] \ar[d] &
X' \ar[d]^a \ar[r] & X \ar[ld]^c &
C \ar[ld] \ar[l] \\
\Spec(k') \ar[r] &
Y &
\Spec(k') \ar[l]
}
$$ dans lequel les deux flèches horizontales inférieures sont identiques. En effet, $X$ est la réunion schématique de $X'$ et $C$, donc un carré cocartésien d'après Morphismes, lemme \ref{morphisms-lemma-scheme-theoretic-union}, et les morphismes $C \to Y$ et $X' \to Y$ coïncident sur $C \cap X'$. Enfin, nous affirmons que $$
c_*\mathcal{O}_X = \mathcal{O}_Y
\quad\text{et}\quad
R^1c_*\mathcal{O}_X = 0
$$ Pour le voir, utilisons la suite exacte $$
0 \to \mathcal{O}_X \to
\mathcal{O}_C \oplus \mathcal{O}_{X'} \to
\mathcal{O}_{C \cap X'} \to 0
$$ de Morphismes, lemme \ref{morphisms-lemma-scheme-theoretic-union}. La suite exacte longue des images directes supérieures est $$
0 \to c_*\mathcal{O}_X \to c_*\mathcal{O}_C \oplus c_*\mathcal{O}_{X'} \to
c_*\mathcal{O}_{C \cap X'} \to
R^1c_*\mathcal{O}_X \to
R^1c_*\mathcal{O}_C \oplus R^1c_*\mathcal{O}_{X'}
$$ Puisque $c|_{X'} = a$ est affine, on a $R^1c_*\mathcal{O}_{X'} = 0$. Comme $c|_C$ se factorise sous la forme $C \to \Spec(k') \to X$ et que $C$ est de genre zéro, on obtient $R^1c_*\mathcal{O}_C = 0$. Puisque $\mathcal{O}_{X'} \to \mathcal{O}_{C \cap X'}$ est surjective et $c|_{X'}$ affine, $c_*\mathcal{O}_{X'} \to c_*\mathcal{O}_{C \cap X'}$ est surjective. Cela démontre $R^1c_*\mathcal{O}_X = 0$. Enfin, on a $\mathcal{O}_Y = c_*\mathcal{O}_X$ d'après la suite exacte et la description du faisceau structural du carré cocartésien dans Compléments sur les morphismes, proposition \ref{more-morphisms-proposition-pushout-along-closed-immersion-and-integral}.

\medskip\noindent
Tout cela signifie que $Y$ est lui aussi un schéma propre sur $k$, de dimension $1$ et tel que $H^0(Y, \mathcal{O}_Y) = k$, dont les singularités sont au plus nodales (lemme \ref{lemma-closed-subscheme-nodal-curve}), et que $Y$ a le même genre que $X$. On dit que $c : X \to Y$ est la {\it contraction d'un pont rationnel}.
\end{example}

\begin{lemma}
\label{lemma-rational-bridge-zero}
Soit $k$ un corps. Soit $X$ un schéma propre sur $k$, de dimension $1$ et tel que $H^0(X, \mathcal{O}_X) = k$. Supposons que les singularités de $X$ soient au plus nodales. Soit $C \subset X$ un pont rationnel (exemple \ref{example-rational-bridge}). Alors $\deg(\omega_X|_C) = 0$.
\end{lemma}

\begin{proof}
Soit $X' \subset X$ comme dans l'exemple. Nous avons une suite exacte courte $$
0 \to \omega_C \to \omega_X|_C \to \mathcal{O}_{C \cap X'} \to 0
$$ Voir les lemmes \ref{lemma-closed-subscheme-reduced-gorenstein}, \ref{lemma-facts-about-nodal-curves} et \ref{lemma-closed-subscheme-nodal-curve}. Avec $k''/k'/k$ comme dans l'exemple, on voit que $\deg(\omega_C) = -2[k' : k]$, puisque $C$ est de genre $0$ (lemme \ref{lemma-rr}), et que $\deg(C \cap X') = [k'' : k] = 2[k' : k]$. Ainsi $\deg(\omega_X|_C) = 0$.
\end{proof}

\begin{lemma}
\label{lemma-rational-bridge-field-extension}
Soit $k$ un corps. Soit $X$ un schéma propre sur $k$, de dimension $1$ et tel que $H^0(X, \mathcal{O}_X) = k$. Supposons que les singularités de $X$ soient au plus nodales. Soit $C \subset X$ un pont rationnel (exemple \ref{example-rational-bridge}). Pour toute extension de corps $K/k$, le changement de base $C_K \subset X_K$ est une réunion disjointe finie de ponts rationnels.
\end{lemma}

\begin{proof}
Soit $k''/k'/k$ comme dans l'exemple. Puisque $k'/k$ est fini séparable, $k' \otimes_k K = K'_1 \times \ldots \times K'_n$ est un produit fini d'extensions finies séparables $K'_i/K$. La décomposition en produit correspondante $k'' \otimes_k K = \prod K''_i$ fournit des extensions d'algèbres séparables de degré $2$, notées $K''_i/K'_i$. Posons $C_i = C_{K'_i}$. Alors $C_K = \coprod C_i$, et chaque $C_i$ est donc de genre $0$, considéré comme une courbe sur $K'_i$, car $H^1(C_K, \mathcal{O}_{C_K}) = 0$ par changement de base plat. Enfin, $X'_K \cap C_i = \Spec(K''_i)$ est de degré $2$ sur $K'_i$, comme voulu.
\end{proof}

\begin{lemma}
\label{lemma-rational-bridge-canonical}
Soit $c : X \to Y$ la contraction d'un pont rationnel (exemple \ref{example-rational-bridge}). Alors $c^*\omega_Y \cong \omega_X$.
\end{lemma}

\begin{proof}
On peut le démontrer par un calcul direct, mais nous préférons utiliser la caractérisation de $\omega_X$ comme le $\mathcal{O}_X$-module cohérent qui représente le foncteur $\textit{Coh}(\mathcal{O}_X) \to \textit{Sets}$, $\mathcal{F} \mapsto \Hom_k(H^1(X, \mathcal{F}), k) =
H^1(X, \mathcal{F})^\vee$; voir le lemme \ref{lemma-duality-dim-1-CM} ou Dualité pour les schémas, lemme \ref{duality-lemma-dualizing-module-proper-over-A}.

\medskip\noindent
Plus précisément, notons $\mathcal{C}_Y$ la catégorie dont les objets sont les $\mathcal{O}_Y$-modules inversibles et dont les flèches sont les homomorphismes de $\mathcal{O}_Y$-modules. Notons $\mathcal{C}_X$ la catégorie dont les objets sont les $\mathcal{O}_X$-modules inversibles $\mathcal{L}$ tels que $\mathcal{L}|_C \cong \mathcal{O}_C$ et dont les flèches sont les homomorphismes de $\mathcal{O}_Y$-modules. Nous affirmons que le foncteur $$
c^* : \mathcal{C}_Y \to \mathcal{C}_X
$$ est une équivalence de catégories. En effet, d'après Compléments sur les morphismes, lemme \ref{more-morphisms-lemma-bijection-on-Pic}, il est essentiellement surjectif. La formule de projection (Cohomologie, lemme \ref{cohomology-lemma-projection-formula}) montre ensuite que $c_*c^*\mathcal{N} = \mathcal{N}$, et donc que $c^*$ est une équivalence dont un quasi-inverse est donné par $c_*$.

\medskip\noindent
Nous affirmons que $\omega_X$ est un objet de $\mathcal{C}_X$. En effet, nous avons une suite exacte courte $$
0 \to \omega_C \to \omega_X|_C \to \mathcal{O}_{C \cap X'} \to 0
$$ Voir le lemme \ref{lemma-closed-subscheme-reduced-gorenstein}. En prenant les degrés, nous trouvons $\deg(\omega_X|_C) = 0$ (nous omettons un petit détail). Ainsi $\omega_X|_C$ est trivial d'après le lemme \ref{lemma-genus-zero-pic}, et $\omega_X$ est un objet de $\mathcal{C}_X$.

\medskip\noindent
Puisque $R^1c_*\mathcal{O}_X = 0$, la formule de projection montre que $R^1c_*c^*\mathcal{N} = 0$ pour $\mathcal{N} \in \Ob(\mathcal{C}_Y)$. La suite spectrale de Leray (Cohomologie, lemme \ref{cohomology-lemma-apply-Leray}) montre donc que le diagramme $$
\xymatrix{
\mathcal{C}_Y \ar[rr]_{c^*} \ar[dr]_{H^1(Y, -)^\vee} & &
\mathcal{C}_X \ar[ld]^{H^1(X, -)^\vee} \\
& \textit{Sets}
}
$$ de catégories et de foncteurs est commutatif. Puisque $\omega_Y \in \Ob(\mathcal{C}_Y)$ représente la flèche sud-est et que $\omega_X \in \Ob(\mathcal{C}_X)$ représente la flèche sud-ouest, nous concluons par le lemme de Yoneda (Catégories, lemme \ref{categories-lemma-yoneda}).
\end{proof}

\begin{lemma}
\label{lemma-no-rational-bridge-ample-genus-geq-2}
Soit $k$ un corps. Soit $X$ un schéma propre sur $k$, de dimension $1$ et tel que $H^0(X, \mathcal{O}_X) = k$. Supposons que \begin{enumerate} \item les singularités de $X$ soient au plus nodales; \item $X$ ne possède pas de queue rationnelle (exemple \ref{example-rational-tail}); \item $X$ ne possède pas de pont rationnel (exemple \ref{example-rational-bridge}); \item le genre $g$ de $X$ soit $\geq 2$. \end{enumerate} Alors $\omega_X$ est ample.
\end{lemma}

\begin{proof}
Il suffit de montrer que $\deg(\omega_X|_C) > 0$ pour toute composante irréductible $C$ de $X$; voir Variétés, lemme \ref{varieties-lemma-ampleness-in-terms-of-degrees-components}. Si $X = C$ est irréductible, cela résulte de $g \geq 2$ et du lemme \ref{lemma-genus-gorenstein}. Sinon, posons $k' = H^0(C, \mathcal{O}_C)$. C'est un corps et une extension finie de $k$, et $[k' : k]$ divise tous les invariants numériques ci-dessous associés à $C$ et aux faisceaux cohérents sur $C$; voir Variétés, lemme \ref{varieties-lemma-divisible}. Soit $X' \subset X$ l'adhérence de $X \setminus C$ comme dans le lemme \ref{lemma-closed-subscheme-reduced-gorenstein}. Nous utiliserons sans autre mention les résultats de ce lemme et des lemmes \ref{lemma-facts-about-nodal-curves} et \ref{lemma-closed-subscheme-nodal-curve}. Nous obtenons une suite exacte courte $$
0 \to \omega_C \to \omega_X|_C \to \mathcal{O}_{C \cap X'} \to 0
$$ Voir le lemme \ref{lemma-closed-subscheme-reduced-gorenstein}. Nous en concluons que $$
\deg(\omega_X|_C) = \deg(C \cap X') + \deg(\omega_C) =
\deg(C \cap X') - 2\chi(C, \mathcal{O}_C)
$$ Par conséquent, si le lemme est faux, alors $$
2[k' : k] \geq \deg(C \cap X') + 2\dim_k H^1(C, \mathcal{O}_C)
$$ Puisque $C \cap X'$ est non vide, et compte tenu de la divisibilité mentionnée ci-dessus, cela n'est possible que si l'une des deux possibilités suivantes se présente : \begin{enumerate} \item[(a)] $C \cap X'$ est un unique point $k'$-rationnel de $C$ et $H^1(C, \mathcal{O}_C) = 0$; \item[(b)] $C \cap X'$ est de degré $2$ sur $k'$ et $H^1(C, \mathcal{O}_C) = 0$. \end{enumerate} La première possibilité signifie que $C$ est une queue rationnelle et la seconde que $C$ est un pont rationnel. Comme les deux sont exclues, la démonstration est achevée.
\end{proof}

\begin{lemma}
\label{lemma-contracting-rational-bridges}
Soit $k$ un corps. Soit $X$ un schéma propre sur $k$, de dimension $1$, tel que $H^0(X, \mathcal{O}_X) = k$ et de genre $g \geq 2$. Supposons que les singularités de $X$ soient au plus nodales et que $X$ ne possède aucune queue rationnelle. Considérons une suite $$
X = X_0 \to X_1 \to \ldots \to X_n = X'
$$ de contractions de ponts rationnels (exemple \ref{example-rational-bridge}), poursuivie jusqu'à ce qu'il n'en reste aucun. Alors $\omega_{X'}$ est ample. Le morphisme $X \to X'$ est indépendant des choix et sa formation commute aux extensions du corps de base.
\end{lemma}

\begin{proof}
Contractons successivement les ponts rationnels jusqu'à ce qu'il n'en reste aucun. Alors $\omega_{X'}$ est ample d'après le lemme \ref{lemma-no-rational-bridge-ample-genus-geq-2}.

\medskip\noindent
Notons $f : X \to X'$ le composé. Le lemme \ref{lemma-rational-bridge-canonical} et un raisonnement par récurrence donnent $f^*\omega_{X'} = \omega_X$. On a $f_*\mathcal{O}_X = \mathcal{O}_{X'}$, puisque cette égalité est vraie pour la contraction d'un pont rationnel. La formule de projection donne donc $f_*f^*\mathcal{L} = \mathcal{L}$ pour tout $\mathcal{O}_{X'}$-module inversible $\mathcal{L}$. Ainsi $$
\Gamma(X', \omega_{X'}^{\otimes m}) = \Gamma(X, \omega_X^{\otimes m})
$$ pour tout $m$. Puisque $X'$ est le Proj de la somme directe de ces modules, d'après Morphismes, lemme \ref{morphisms-lemma-proper-ample-is-proj}, nous concluons que le morphisme $X \to X'$ est entièrement canonique.

\medskip\noindent
Soit $K/k$ une extension de corps. Alors $\omega_{X_K}$ est l'image réciproque de $\omega_X$ (lemme \ref{lemma-sanity-check-duality}), et $\Gamma(X, \omega_X^{\otimes m}) \otimes_k K$ est égal à $\Gamma(X_K, \omega_{X_K}^{\otimes m})$ d'après Cohomologie des schémas, lemme \ref{coherent-lemma-flat-base-change-cohomology}. Ainsi, la formation de $f : X \to X'$ commute au changement de base d'après $K/k$ et les arguments ci-dessus. Nous omettons quelques détails.
\end{proof}






\section{Contraction vers une courbe stable}
\label{section-contracting-to-stable}

\noindent
Dans cette section, nous combinons les morphismes de contraction obtenus dans les sections \ref{section-contracting-rational-tails} et \ref{section-contracting-rational-bridges}. Supposons donc que $k$ soit un corps et soit $X$ un schéma propre sur $k$, de dimension $1$, tel que $H^0(X, \mathcal{O}_X) = k$ et de genre $g \geq 2$. Supposons que les singularités de $X$ soient au plus nodales. En composant le morphisme du lemme \ref{lemma-contracting-rational-tails} avec celui du lemme \ref{lemma-contracting-rational-bridges}, nous obtenons un morphisme $$
c : X \longrightarrow Y
$$ tel que $Y$ soit lui aussi un schéma propre sur $k$, de dimension $1$, dont les singularités sont au plus nodales, tel que $k = H^0(Y, \mathcal{O}_Y)$, de genre $g$, tel que $\mathcal{O}_Y = c_*\mathcal{O}_X$ et $R^1c_*\mathcal{O}_X = 0$, et tel que $\omega_Y$ soit ample sur $Y$. Le lemme \ref{lemma-characterize-contraction-to-stable} montre que ces conditions caractérisent en fait ce morphisme.

\begin{lemma}
\label{lemma-contract-gorenstein-canonical}
Soit $k$ un corps. Soit $c : X \to Y$ un morphisme de schémas propres sur $k$. Supposons que \begin{enumerate} \item $\mathcal{O}_Y = c_*\mathcal{O}_X$ et $R^1c_*\mathcal{O}_X = 0$; \item $X$ et $Y$ soient réduits, de Gorenstein et de dimension $1$; \item $\exists\ m \in \mathbf{Z}$, avec $H^1(X, \omega_X^{\otimes m}) = 0$ et $\omega_X^{\otimes m}$ engendré par ses sections globales. \end{enumerate} Alors $c^*\omega_Y \cong \omega_X$.
\end{lemma}

\begin{proof}
Les fibres de $c$ sont géométriquement connexes d'après Compléments sur les morphismes, théorème \ref{more-morphisms-theorem-stein-factorization-Noetherian}. En particulier, $c$ est surjectif. Il existe un nombre fini de points fermés $y = y_1, \ldots, y_r$ de $Y$ au-dessus desquels $X_y$ est de dimension $1$, et sur $Y \setminus \{y_1, \ldots, y_r\}$ le morphisme $c$ est un isomorphisme. Nous omettons quelques détails; indication : hors de $\{y_1, \ldots, y_r\}$, le morphisme $c$ est fini; voir Cohomologie des schémas, lemme \ref{coherent-lemma-characterize-finite}.

\medskip\noindent
Construisons soigneusement une application $b : c^*\omega_Y \to \omega_X$. Notons $f : X \to \Spec(k)$ et $g : Y \to \Spec(k)$ les morphismes structuraux. On a $f^!k = \omega_X[1]$ et $g^!k = \omega_Y[1]$; voir le lemme \ref{lemma-duality-dim-1} et sa démonstration. Alors $f^! = c^! \circ g^!$, et donc $c^!\omega_Y = \omega_X$. Il existe ainsi un isomorphisme fonctoriel $$
\Hom_{D(\mathcal{O}_X)}(\mathcal{F}, \omega_X)
\longrightarrow
\Hom_{D(\mathcal{O}_Y)}(Rc_*\mathcal{F}, \omega_Y)
$$ pour tout $\mathcal{O}_X$-module cohérent $\mathcal{F}$, par définition de $c^!$\footnote{Comme restriction de l'adjoint à droite de Dualité pour les schémas, lemme \ref{duality-lemma-twisted-inverse-image}, à $D^+_\QCoh(\mathcal{O}_Y)$.}. Cet isomorphisme est induit par une application trace $t : Rc_*\omega_X \to \omega_Y$, la co-unité de l'adjonction. Par la formule de projection (Cohomologie, lemme \ref{cohomology-lemma-projection-formula}), l'application canonique $a : \omega_Y \to Rc_*c^*\omega_Y$ est un isomorphisme. En combinant ce qui précède, nous obtenons une application canonique $b : c^*\omega_Y \to \omega_X$ telle que $$
t \circ Rc_*(b) = a^{-1}
$$ En particulier, la restriction de $b$ à $c^{-1}(Y \setminus \{y_1, \ldots, y_r\})$ est un isomorphisme, car il s'agit d'une application entre modules inversibles dont le composé avec une autre application est l'isomorphisme $a^{-1}$.

\medskip\noindent
Choisissons $m \in \mathbf{Z}$ comme dans (3) et considérons l'application $$
b^{\otimes m} :
\Gamma(Y, \omega_Y^{\otimes m})
\longrightarrow
\Gamma(X, \omega_X^{\otimes m})
$$ Cette application est injective parce que $Y$ est réduit et par la dernière propriété de $b$ mentionnée lors de sa construction. Le théorème de Riemann--Roch (lemme \ref{lemma-rr}) donne $\chi(X, \omega_X^{\otimes m}) =\chi(Y, \omega_Y^{\otimes m})$. Ainsi $$
\dim_k \Gamma(Y, \omega_Y^{\otimes m}) \geq
\dim_k \Gamma(X, \omega_X^{\otimes m}) = \chi(X, \omega_X^{\otimes m})
$$, et nous concluons que $b^{\otimes m}$ induit un isomorphisme sur les sections globales. Donc $b^{\otimes m} : c^*\omega_Y^{\otimes m} \to \omega_X^{\otimes m}$ est surjective, puisque des générateurs de $\omega_X^{\otimes m}$ appartiennent à son image. Par conséquent, $b^{\otimes m}$ est un isomorphisme, puis $b$ est un isomorphisme.
\end{proof}

\begin{lemma}
\label{lemma-characterize-contraction-to-stable}
Soit $k$ un corps. Soit $X$ un schéma propre sur $k$, de dimension $1$, tel que $H^0(X, \mathcal{O}_X) = k$ et de genre $g \geq 2$. Supposons que les singularités de $X$ soient au plus nodales. Il existe un unique morphisme, à unique isomorphisme près, $$
c : X \longrightarrow Y
$$ de schémas sur $k$ possédant les propriétés suivantes : \begin{enumerate} \item $Y$ est propre sur $k$, $\dim(Y) = 1$, et les singularités de $Y$ sont au plus nodales; \item $\mathcal{O}_Y = c_*\mathcal{O}_X$ et $R^1c_*\mathcal{O}_X = 0$; \item $\omega_Y$ est ample sur $Y$.
\end{enumerate}
\end{lemma}

\begin{proof}
Existence. Un morphisme possédant toutes les propriétés énumérées existe en combinant les lemmes \ref{lemma-contracting-rational-tails} et \ref{lemma-contracting-rational-bridges}, comme expliqué dans l'introduction de cette section. De plus, il peut s'écrire comme un composé $$
X \to X_1 \to X_2 \ldots \to X_n \to X_{n + 1} \to \ldots \to X_{n + n'}
$$ dans lequel les $n$ premiers morphismes sont des contractions de queues rationnelles et les $n'$ derniers des contractions de ponts rationnels. Remarquons que la propriété (2) est satisfaite par toute contraction d'une queue rationnelle (exemple \ref{example-rational-tail}) et par toute contraction d'un pont rationnel (exemple \ref{example-rational-bridge}). On voit aisément que cette propriété est préservée par composition.

\medskip\noindent
Unicité. Soit $c : X \to Y$ un morphisme satisfaisant les conditions (1), (2) et (3). Nous allons montrer qu'il existe un unique isomorphisme $X_{n + n'} \to Y$ compatible avec les morphismes $X \to X_{n + n'}$ et $c$.

\medskip\noindent
Avant de commencer la démonstration, faisons quelques observations sur $c$. Tout d'abord, les fibres de $c$ sont géométriquement connexes d'après Compléments sur les morphismes, théorème \ref{more-morphisms-theorem-stein-factorization-Noetherian}. En particulier, $c$ est surjectif. Pour tout point fermé $y \in Y$, la fibre $X_y$ satisfait $$
H^1(X_y, \mathcal{O}_{X_y}) = 0
\quad\text{et}\quad
H^0(X_y, \mathcal{O}_{X_y}) = \kappa(y)
$$ La première égalité résulte de Compléments sur les morphismes, lemme \ref{more-morphisms-lemma-check-h1-fibre-zero}, et la seconde de Compléments sur les morphismes, lemme \ref{more-morphisms-lemma-h1-fibre-zero-check-h0-kappa}. Ainsi, soit $X_y = x$, où $x$ est l'unique point de $X$ s'envoyant sur $y$ et possède le même corps résiduel que $y$, soit $X_y$ est un schéma propre de dimension $1$ sur $\kappa(y)$. Dans le second cas, remarquons que $X_y$ est de Cohen--Macaulay (lemme \ref{lemma-automatic}). Toutefois, puisque $X$ est réduit, $X_y$ doit être réduit en tous ses points génériques (nous omettons les détails), et donc $X_y$ est réduit d'après Propriétés, lemme \ref{properties-lemma-criterion-reduced}. Il s'ensuit que les singularités de $X_y$ sont au plus nodales (lemme \ref{lemma-closed-subscheme-nodal-curve}). Remarquons que $X_y$ est de genre zéro, d'après ce qui précède. Enfin, il n'existe qu'un nombre fini de points $y$ au-dessus desquels la fibre $X_y$ est de dimension $1$; notons-les $\{y_1, \ldots, y_r\}$. Le sous-schéma $c^{-1}(Y \setminus \{y_1, \ldots, y_r\})$ s'envoie isomorphiquement sur $Y \setminus \{y_1, \ldots, y_r\}$ par $c$. Nous omettons quelques détails; indication : hors de $\{y_1, \ldots, y_r\}$, le morphisme $c$ est fini; voir Cohomologie des schémas, lemme \ref{coherent-lemma-characterize-finite}.

\medskip\noindent
Soit $C \subset X$ une queue rationnelle. Nous affirmons que $c$ envoie $C$ sur un point. Supposons le contraire pour aboutir à une contradiction. Alors l'image de $C$ est une composante irréductible $D \subset Y$. Rappelons que $H^0(C, \mathcal{O}_C) = k'$ est une extension finie séparable de $k$, et que $C$ possède un point $k'$-rationnel $x$ qui est aussi l'unique point d'intersection de $C$ avec le « reste » de $X$. La discussion générale ci-dessus montre que $C \setminus \{x\} \subset c^{-1}(Y \setminus \{y_1, \ldots, y_r\})$ s'envoie isomorphiquement sur un ouvert $V$ de $D$. Soit $y = c(x) \in D$. Remarquons que $y$ est l'unique point de $D$ qui rencontre le « reste » de $Y$. Si $y \not \in \{y_1, \ldots, y_r\}$, alors $C \cong D$, et il est clair que $D$ est une queue rationnelle de $Y$, ce qui contredit l'amplitude de $\omega_Y$ (lemme \ref{lemma-rational-tail-negative}). Ainsi $y \in \{y_1, \ldots, y_r\}$ et $\dim(X_y) = 1$. Alors $x \in X_y \cap C$, et $x$ est un point lisse de $X_y$ et de $C$ (lemme \ref{lemma-closed-subscheme-nodal-curve}). Si $y \in D$ est un point singulier de $D$, alors $y$ est un nœud et $Y = D$, car aucune autre composante de $Y$ ne peut passer par $y$ d'après le lemme \ref{lemma-closed-subscheme-nodal-curve}. Alors $X = X_y \cup C$, ce qui signifie que $g = 0$, puisque ce nombre est égal au genre de $X_y$ d'après la discussion de l'exemple \ref{example-rational-tail}; contradiction. Si $y \in D$ est un point lisse de $D$, alors $C \to D$ est un isomorphisme, car le modèle projectif non singulier est unique et $C$ et $D$ sont birationnels; voir la section \ref{section-curves-function-fields}. Alors $D$ est une queue rationnelle de $Y$, ce qui contredit l'amplitude de $\omega_Y$.

\medskip\noindent
Supposons $n \geq 1$. Si $C \subset X$ est la queue rationnelle contractée par $X \to X_1$, le paragraphe précédent montre que $C$ est envoyé sur un point de $Y$. Ainsi $c : X \to Y$ se factorise par $X \to X_1$, puisque $X$ est le carré cocartésien de $C$ et $X_1$; voir la discussion de l'exemple \ref{example-rational-tail}. Après avoir remplacé $X$ par $X_1$, nous avons diminué $n$. Par récurrence, nous pouvons supposer $n = 0$, c'est-à-dire que $X$ ne possède aucune queue rationnelle.

\medskip\noindent
Supposons $n = 0$, c'est-à-dire que $X$ ne possède aucune queue rationnelle. Alors $\omega_X^{\otimes 2}$ et $\omega_X^{\otimes 3}$ sont engendrés par leurs sections globales d'après le lemme \ref{lemma-no-rational-tail-semiample-genus-geq-2}. Il s'ensuit que $H^1(X, \omega_X^{\otimes 3}) = 0$ d'après le lemme \ref{lemma-vanishing-twist}. En appliquant le lemme \ref{lemma-contract-gorenstein-canonical} avec $m = 3$, nous obtenons $c^*\omega_Y \cong \omega_X$. Nous avons aussi $\omega_X = (X \to X_{n'})^*\omega_{X_{n'}}$ d'après le lemme \ref{lemma-rational-bridge-canonical} et par récurrence. En appliquant la formule de projection à $c$ et à $X \to X_{n'}$, nous concluons que $$
\Gamma(X_{n'}, \omega_{X_{n'}}^{\otimes m}) =
\Gamma(X, \omega_X^{\otimes m}) =
\Gamma(Y, \omega_Y^{\otimes m})
$$ pour tout $m$. Puisque $X_{n'}$ et $Y$ sont les Proj des sommes directes correspondantes, d'après Morphismes, lemme \ref{morphisms-lemma-proper-ample-is-proj}, nous obtenons un isomorphisme canonique $X_{n'} = Y$, comme voulu. Nous omettons la vérification que c'est l'unique isomorphisme rendant le diagramme commutatif.
\end{proof}

\begin{lemma}
\label{lemma-tricanonical}
Soit $k$ un corps. Soit $X$ un schéma propre sur $k$, de dimension $1$, tel que $H^0(X, \mathcal{O}_X) = k$ et de genre $g \geq 2$. Supposons que les singularités de $X$ soient au plus nodales et que $\omega_X$ soit ample. Alors $\omega_X^{\otimes 3}$ est très ample et $H^1(X, \omega_X^{\otimes 3}) = 0$.
\end{lemma}

\begin{proof}
En combinant Variétés, lemme \ref{varieties-lemma-ampleness-in-terms-of-degrees-components}, et les lemmes \ref{lemma-rational-tail-negative} et \ref{lemma-rational-bridge-zero}, nous voyons que $X$ ne contient ni queue rationnelle ni pont rationnel. Alors $\omega_X^{\otimes 3}$ est engendré par ses sections globales d'après le lemme \ref{lemma-contracting-rational-tails}. Choisissons une $k$-base $s_0, \ldots, s_n$ de $H^0(X, \omega_X^{\otimes 3})$. Nous obtenons un morphisme $$
\varphi_{\omega_X^{\otimes 3}, (s_0, \ldots, s_n)} :
X \longrightarrow \mathbf{P}^n_k
$$ Voir Constructions, section \ref{constructions-section-projective-space}. Le lemme affirme que ce morphisme est une immersion fermée. Pour le vérifier, nous pouvons remplacer $k$ par sa clôture algébrique; voir Descente, lemme \ref{descent-lemma-descending-property-closed-immersion}. Nous pouvons donc supposer $k$ algébriquement clos.

\medskip\noindent
Supposons $k$ algébriquement clos. Nous allons utiliser Variétés, lemme \ref{varieties-lemma-variant-separate-points-tangent-vectors}, pour démontrer le lemme. Soit $Z \subset X$ un sous-schéma fermé de degré $2$ sur $Z$, de faisceau d'idéaux $\mathcal{I} \subset \mathcal{O}_X$. Il faut montrer que $$
H^0(X, \mathcal{L}) \to H^0(Z, \mathcal{L}|_Z)
$$ est surjective. Il suffit donc d'établir $H^1(X, \mathcal{I}\mathcal{L}) = 0$. Pour cela, nous utiliserons le lemme \ref{lemma-vanishing-on-gorenstein}. Il suffit ainsi de montrer que $$
3\deg(\omega_X|_Y) > -2\chi(Y, \mathcal{O}_Y) + \deg(Z \cap Y)
$$ pour tout sous-schéma fermé réduit connexe $Y \subset X$. Puisque $k$ est algébriquement clos et que $Y$ est connexe et réduit, on a $H^0(Y, \mathcal{O}_Y) = k$ (Variétés, lemme \ref{varieties-lemma-proper-geometrically-reduced-global-sections}). Par conséquent, $\chi(Y, \mathcal{O}_Y) = 1 - \dim H^1(Y, \mathcal{O}_Y)$. Il reste donc à montrer $$
3\deg(\omega_X|_Y) > -2 + 2\dim H^1(Y, \mathcal{O}_Y) + \deg(Z \cap Y)
$$, ce qui résulte du lemme \ref{lemma-no-rational-tail}, sauf éventuellement si $Y = X$ ou si $\deg(\omega_X|_Y) = 0$. Puisque $\omega_X$ est ample, la seconde possibilité ne se présente pas; voir le premier lemme cité dans cette démonstration. Enfin, si $Y = X$, le théorème de Riemann--Roch (lemme \ref{lemma-rr}) et le fait que $g \geq 2$ montrent que l'inégalité est satisfaite. Le même argument avec $Z = \emptyset$ donne $H^1(X, \omega_X^{\otimes 3}) = 0$.
\end{proof}







\section{Champs de vecteurs}
\label{section-vector-fields}

\noindent
Dans cette section, nous étudions l'espace des champs de vecteurs sur une courbe. Les champs de vecteurs correspondent aux automorphismes infinitésimaux; voir Compléments sur les morphismes, section \ref{more-morphisms-section-action-by-derivations}. Ils jouent donc un rôle important dans la théorie des espaces de modules.

\medskip\noindent
Soit $k$ un corps algébriquement clos. Soit $X$ un schéma de type fini sur $k$. Soit $x \in X$ un point fermé. Nous dirons qu'un élément $D \in \text{Der}_k(\mathcal{O}_X, \mathcal{O}_X)$ {\it fixe $x$} si $D(\mathcal{I}) \subset \mathcal{I}$, où $\mathcal{I} \subset \mathcal{O}_X$ est le faisceau d'idéaux de $x$.

\begin{lemma}
\label{lemma-smooth-vector-fields}
Soit $k$ un corps algébriquement clos. Soit $X$ une courbe lisse, propre et connexe sur $k$. Soit $g$ le genre de $X$. \begin{enumerate} \item Si $g \geq 2$, alors $\text{Der}_k(\mathcal{O}_X, \mathcal{O}_X)$ est nul; \item si $g = 1$ et si $D \in \text{Der}_k(\mathcal{O}_X, \mathcal{O}_X)$ est non nul, alors $D$ ne fixe aucun point fermé de $X$; \item si $g = 0$ et si $D \in \text{Der}_k(\mathcal{O}_X, \mathcal{O}_X)$ est non nul, alors $D$ fixe au plus $2$ points fermés de $X$.
\end{enumerate}
\end{lemma}

\begin{proof}
Rappelons que nous avons une $k$-dérivation universelle $d : \mathcal{O}_X \to \Omega_{X/k}$, et donc $D = \theta \circ d$ pour une certaine application $\mathcal{O}_X$-linéaire $\theta : \Omega_{X/k} \to \mathcal{O}_X$. Rappelons que $\Omega_{X/k} \cong \omega_X$; voir le lemme \ref{lemma-duality-dim-1}. Le théorème de Riemann--Roch donne $\deg(\omega_X) = 2g - 2$ (lemme \ref{lemma-rr}). Ainsi $\theta$ est nécessairement nul si $g > 1$, d'après Variétés, lemme \ref{varieties-lemma-check-invertible-sheaf-trivial}. Cela démontre la partie (1). Si $g = 1$, un $\theta$ non nul ne s'annule nulle part; si $g = 0$, un $\theta$ non nul s'annule suivant un diviseur de degré $2$. Les parties (2) et (3) en résultent dès que l'on montre que l'annulation de $\theta$ en un point fermé $x \in X$ équivaut à dire que $D$ fixe $x$, au sens défini ci-dessus. Soit $z \in \mathcal{O}_{X, x}$ une uniformisante. Alors $dz$ est un élément de base de $\Omega_{X, x}$; voir le lemme \ref{lemma-uniformizer-works}. Puisque $D(z) = \theta(dz)$, nous concluons.
\end{proof}

\begin{lemma}
\label{lemma-nodal-vector-fields}
Soit $k$ un corps algébriquement clos. Soit $X$ un schéma propre connexe de dimension $1$ sur $k$, à singularités au plus nodales. Soit $\nu : X^\nu \to X$ la normalisation. Soit $S \subset X^\nu$ l'ensemble des points où $\nu$ n'est pas un isomorphisme. Alors
$$
\text{Der}_k(\mathcal{O}_X, \mathcal{O}_X) =
\{D' \in \text{Der}_k(\mathcal{O}_{X^\nu}, \mathcal{O}_{X^\nu}) \mid
D' \text{ fixe tout }x^\nu \in S\}
$$
\end{lemma}

\begin{proof}
Soit $x \in X$ un nœud. Soient $x', x'' \in X^\nu$ les images réciproques de $x$. Tout nœud est déployé puisque $k$ est algébriquement clos; voir la définition \ref{definition-split-node} et le lemme \ref{lemma-split-node}. Soient $u \in \mathcal{O}_{X^\nu, x'}$ et $v \in \mathcal{O}_{X^\nu, x''}$ des uniformisantes. Remarquons que nous avons une suite exacte $$
0 \to \mathcal{O}_{X, x} \to
\mathcal{O}_{X^\nu, x'} \times \mathcal{O}_{X^\nu, x''} \to k \to 0
$$ Cela résulte du lemme \ref{lemma-multicross}. Nous pouvons donc considérer $u$ et $v$ comme des éléments de $\mathcal{O}_{X, x}$ tels que $uv = 0$.

\medskip\noindent
Soit $D \in \text{Der}_k(\mathcal{O}_X, \mathcal{O}_X)$. Alors $0 = D(uv) = vD(u) + uD(v)$. Puisque $(u)$ est l'annulateur de $v$ dans $\mathcal{O}_{X, x}$ et réciproquement, on voit que $D(u) \in (u)$ et $D(v) \in (v)$. Comme $\mathcal{O}_{X^\nu, x'} = k + (u)$, nous concluons que nous pouvons prolonger $D$ en $\mathcal{O}_{X^\nu, x'}$ et que, de plus, ce prolongement fixe $x'$. Cela produit un $D'$ dans le membre de droite de l'égalité. Réciproquement, étant donné $D'$ qui fixe $x'$ et $x''$, nous voyons que $D'$ préserve le sous-anneau $\mathcal{O}_{X, x} \subset
\mathcal{O}_{X^\nu, x'} \times \mathcal{O}_{X^\nu, x''}$; c'est ainsi que l'on passe du membre de droite au membre de gauche de l'égalité.
\end{proof}

\begin{lemma}
\label{lemma-stable-vector-fields}
Soit $k$ un corps algébriquement clos. Soit $X$ un schéma propre connexe de dimension $1$ sur $k$, à singularités au plus nodales. Supposons que le genre de $X$ soit au moins $2$ et que $X$ ne possède ni queue rationnelle ni pont rationnel. Alors $\text{Der}_k(\mathcal{O}_X, \mathcal{O}_X) = 0$.
\end{lemma}

\begin{proof}
Soit $D \in \text{Der}_k(\mathcal{O}_X, \mathcal{O}_X)$. Soit $X^\nu$ la normalisation de $X$. Soit $D' \in \text{Der}_k(\mathcal{O}_{X^\nu}, \mathcal{O}_{X^\nu})$ l'élément correspondant à $D$ par le lemme \ref{lemma-nodal-vector-fields}. Soit $C \subset X^\nu$ une composante irréductible. Si le genre de $C$ est $> 1$, alors $D'|_{\mathcal{O}_C} = 0$ d'après la partie (1) du lemme \ref{lemma-smooth-vector-fields}. Si le genre de $C$ est $1$, il existe au moins un point fermé $c$ de $C$ qui s'envoie sur un nœud de $X$, car sinon $X \cong C$ serait de genre $1$. Par la correspondance, cela signifie que $D'|_{\mathcal{O}_C}$ fixe $c$, et il est donc nul d'après la partie (2) du lemme \ref{lemma-smooth-vector-fields}. Enfin, si le genre de $C$ est nul, il existe au moins $3$ points fermés deux à deux distincts $c_1, c_2, c_3 \in C$ qui s'envoient sur des nœuds de $X$. Sinon, soit $X$ serait $C$ avec deux points recollés, deux points de $C$ s'envoyant sur le même nœud, soit $C$ serait un pont rationnel, deux points s'envoyant sur des nœuds distincts de $X$, soit $C$ serait une queue rationnelle, un point s'envoyant sur un nœud de $X$. Ces trois possibilités sont exclues puisque $C$ est de genre $\geq 2$ et ne possède ni pont rationnel ni queue rationnelle. Ainsi $D'|_{\mathcal{O}_C}$ fixe $c_1, c_2, c_3$ et est donc nul d'après la partie (3) du lemme \ref{lemma-smooth-vector-fields}.
\end{proof}





\begin{multicols}{2}[\section{Autres chapitres}]
\noindent
Préliminaires
\begin{enumerate}
\item \hyperref[introduction-section-phantom]{Introduction}
\item \hyperref[conventions-section-phantom]{Conventions}
\item \hyperref[sets-section-phantom]{Théorie des ensembles}
\item \hyperref[categories-section-phantom]{Catégories}
\item \hyperref[topology-section-phantom]{Topologie}
\item \hyperref[sheaves-section-phantom]{Faisceaux sur les espaces}
\item \hyperref[sites-section-phantom]{Sites et faisceaux}
\item \hyperref[stacks-section-phantom]{Champs}
\item \hyperref[fields-section-phantom]{Corps}
\item \hyperref[algebra-section-phantom]{Algèbre commutative}
\item \hyperref[brauer-section-phantom]{Groupes de Brauer}
\item \hyperref[homology-section-phantom]{Algèbre homologique}
\item \hyperref[derived-section-phantom]{Catégories dérivées}
\item \hyperref[simplicial-section-phantom]{Méthodes simpliciales}
\item \hyperref[more-algebra-section-phantom]{Compléments d'algèbre}
\item \hyperref[smoothing-section-phantom]{Lissage des morphismes d'anneaux}
\item \hyperref[modules-section-phantom]{Faisceaux de modules}
\item \hyperref[sites-modules-section-phantom]{Modules sur les sites}
\item \hyperref[injectives-section-phantom]{Objets injectifs}
\item \hyperref[cohomology-section-phantom]{Cohomologie des faisceaux}
\item \hyperref[sites-cohomology-section-phantom]{Cohomologie sur les sites}
\item \hyperref[dga-section-phantom]{Algèbre différentielle graduée}
\item \hyperref[dpa-section-phantom]{Algèbre à puissances divisées}
\item \hyperref[sdga-section-phantom]{Faisceaux différentiels gradués}
\item \hyperref[hypercovering-section-phantom]{Hyperrecouvrements}
\end{enumerate}
Schémas
\begin{enumerate}
\setcounter{enumi}{25}
\item \hyperref[schemes-section-phantom]{Schémas}
\item \hyperref[constructions-section-phantom]{Constructions de schémas}
\item \hyperref[properties-section-phantom]{Propriétés des schémas}
\item \hyperref[morphisms-section-phantom]{Morphismes de schémas}
\item \hyperref[coherent-section-phantom]{Cohomologie des schémas}
\item \hyperref[divisors-section-phantom]{Diviseurs}
\item \hyperref[limits-section-phantom]{Limites de schémas}
\item \hyperref[varieties-section-phantom]{Variétés}
\item \hyperref[topologies-section-phantom]{Topologies sur les schémas}
\item \hyperref[descent-section-phantom]{Descente}
\item \hyperref[perfect-section-phantom]{Catégories dérivées de schémas}
\item \hyperref[more-morphisms-section-phantom]{Compléments sur les morphismes}
\item \hyperref[flat-section-phantom]{Compléments sur la platitude}
\item \hyperref[groupoids-section-phantom]{Schémas en groupoïdes}
\item \hyperref[more-groupoids-section-phantom]{Compléments sur les schémas en groupoïdes}
\item \hyperref[etale-section-phantom]{Morphismes étales de schémas}
\end{enumerate}
Thèmes de théorie des schémas
\begin{enumerate}
\setcounter{enumi}{41}
\item \hyperref[chow-section-phantom]{Homologie de Chow}
\item \hyperref[intersection-section-phantom]{Théorie de l'intersection}
\item \hyperref[pic-section-phantom]{Schémas de Picard des courbes}
\item \hyperref[weil-section-phantom]{Théories cohomologiques de Weil}
\item \hyperref[adequate-section-phantom]{Modules adéquats}
\item \hyperref[dualizing-section-phantom]{Complexes dualisants}
\item \hyperref[duality-section-phantom]{Dualité pour les schémas}
\item \hyperref[discriminant-section-phantom]{Discriminants et différentes}
\item \hyperref[derham-section-phantom]{Cohomologie de de Rham}
\item \hyperref[local-cohomology-section-phantom]{Cohomologie locale}
\item \hyperref[algebraization-section-phantom]{Géométrie algébrique et formelle}
\item \hyperref[curves-section-phantom]{Courbes algébriques}
\item \hyperref[resolve-section-phantom]{Résolution des surfaces}
\item \hyperref[models-section-phantom]{Réduction semi-stable}
\item \hyperref[functors-section-phantom]{Foncteurs et morphismes}
\item \hyperref[equiv-section-phantom]{Catégories dérivées de variétés}
\item \hyperref[pione-section-phantom]{Groupes fondamentaux de schémas}
\item \hyperref[etale-cohomology-section-phantom]{Cohomologie étale}
\item \hyperref[crystalline-section-phantom]{Cohomologie cristalline}
\item \hyperref[proetale-section-phantom]{Cohomologie pro-étale}
\item \hyperref[relative-cycles-section-phantom]{Cycles relatifs}
\item \hyperref[more-etale-section-phantom]{Compléments de cohomologie étale}
\item \hyperref[trace-section-phantom]{La formule des traces}
\end{enumerate}
Espaces algébriques
\begin{enumerate}
\setcounter{enumi}{64}
\item \hyperref[spaces-section-phantom]{Espaces algébriques}
\item \hyperref[spaces-properties-section-phantom]{Propriétés des espaces algébriques}
\item \hyperref[spaces-morphisms-section-phantom]{Morphismes d'espaces algébriques}
\item \hyperref[decent-spaces-section-phantom]{Espaces algébriques décents}
\item \hyperref[spaces-cohomology-section-phantom]{Cohomologie des espaces algébriques}
\item \hyperref[spaces-limits-section-phantom]{Limites d'espaces algébriques}
\item \hyperref[spaces-divisors-section-phantom]{Diviseurs sur les espaces algébriques}
\item \hyperref[spaces-over-fields-section-phantom]{Espaces algébriques sur des corps}
\item \hyperref[spaces-topologies-section-phantom]{Topologies sur les espaces algébriques}
\item \hyperref[spaces-descent-section-phantom]{Descente et espaces algébriques}
\item \hyperref[spaces-perfect-section-phantom]{Catégories dérivées d'espaces}
\item \hyperref[spaces-more-morphisms-section-phantom]{Compléments sur les morphismes d'espaces}
\item \hyperref[spaces-flat-section-phantom]{Platitude sur les espaces algébriques}
\item \hyperref[spaces-groupoids-section-phantom]{Groupoïdes dans les espaces algébriques}
\item \hyperref[spaces-more-groupoids-section-phantom]{Compléments sur les groupoïdes dans les espaces}
\item \hyperref[bootstrap-section-phantom]{Amorçage}
\item \hyperref[spaces-pushouts-section-phantom]{Sommes amalgamées d'espaces algébriques}
\end{enumerate}
Thèmes de géométrie
\begin{enumerate}
\setcounter{enumi}{81}
\item \hyperref[spaces-chow-section-phantom]{Groupes de Chow des espaces}
\item \hyperref[groupoids-quotients-section-phantom]{Quotients de groupoïdes}
\item \hyperref[spaces-more-cohomology-section-phantom]{Compléments sur la cohomologie des espaces}
\item \hyperref[spaces-simplicial-section-phantom]{Espaces simpliciaux}
\item \hyperref[spaces-duality-section-phantom]{Dualité pour les espaces}
\item \hyperref[formal-spaces-section-phantom]{Espaces algébriques formels}
\item \hyperref[restricted-section-phantom]{Algébrisation des espaces formels}
\item \hyperref[spaces-resolve-section-phantom]{Résolution des surfaces revisitée}
\end{enumerate}
Théorie des déformations
\begin{enumerate}
\setcounter{enumi}{89}
\item \hyperref[formal-defos-section-phantom]{Théorie formelle des déformations}
\item \hyperref[defos-section-phantom]{Théorie des déformations}
\item \hyperref[cotangent-section-phantom]{Le complexe cotangent}
\item \hyperref[examples-defos-section-phantom]{Problèmes de déformation}
\end{enumerate}
Champs algébriques
\begin{enumerate}
\setcounter{enumi}{93}
\item \hyperref[algebraic-section-phantom]{Champs algébriques}
\item \hyperref[examples-stacks-section-phantom]{Exemples de champs}
\item \hyperref[stacks-sheaves-section-phantom]{Faisceaux sur les champs algébriques}
\item \hyperref[criteria-section-phantom]{Critères de représentabilité}
\item \hyperref[artin-section-phantom]{Axiomes d'Artin}
\item \hyperref[quot-section-phantom]{Espaces de Quot et de Hilbert}
\item \hyperref[stacks-properties-section-phantom]{Propriétés des champs algébriques}
\item \hyperref[stacks-morphisms-section-phantom]{Morphismes de champs algébriques}
\item \hyperref[stacks-limits-section-phantom]{Limites de champs algébriques}
\item \hyperref[stacks-cohomology-section-phantom]{Cohomologie des champs algébriques}
\item \hyperref[stacks-perfect-section-phantom]{Catégories dérivées de champs}
\item \hyperref[stacks-introduction-section-phantom]{Introduction aux champs algébriques}
\item \hyperref[stacks-more-morphisms-section-phantom]{Compléments sur les morphismes de champs}
\item \hyperref[stacks-geometry-section-phantom]{La géométrie des champs}
\end{enumerate}
Thèmes de théorie des espaces de modules
\begin{enumerate}
\setcounter{enumi}{107}
\item \hyperref[moduli-section-phantom]{Champs de modules}
\item \hyperref[moduli-curves-section-phantom]{Espaces de modules de courbes}
\end{enumerate}
Divers
\begin{enumerate}
\setcounter{enumi}{109}
\item \hyperref[examples-section-phantom]{Exemples}
\item \hyperref[exercises-section-phantom]{Exercices}
\item \hyperref[guide-section-phantom]{Guide bibliographique}
\item \hyperref[desirables-section-phantom]{Desiderata}
\item \hyperref[coding-section-phantom]{Style de programmation}
\item \hyperref[obsolete-section-phantom]{Obsolète}
\item \hyperref[fdl-section-phantom]{Licence de documentation libre GNU}
\item \hyperref[index-section-phantom]{Index généré automatiquement}
\end{enumerate}
\end{multicols}


\bibliography{my}
\bibliographystyle{amsalpha}

\end{document}
